\RequirePackage{fix-cm}
\documentclass[11pt,a4paper]{article}
\usepackage[margin=25mm]{geometry}
\usepackage[T1]{fontenc}
\usepackage{lmodern,microtype,amsmath,amssymb,amsthm,mathtools}
\usepackage{booktabs,array,enumitem,needspace,graphicx,adjustbox,etoolbox,listings,xcolor}
\usepackage[colorlinks=true,linkcolor=black,citecolor=magenta,urlcolor=black]{hyperref}
\hypersetup{bookmarksdepth=3,bookmarksnumbered=true,pdftitle={Report on regular Hadamard matrices of order 11},pdfauthor={Teo and GPT}}
\newtheorem{theorem}{Theorem}[section]
\newtheorem{proposition}[theorem]{Proposition}
\newtheorem{lemma}[theorem]{Lemma}
\newtheorem{corollary}[theorem]{Corollary}
\newtheorem{conjecture}[theorem]{Conjecture}
\theoremstyle{plain}
\newtheorem{definition}[theorem]{Definition}
\theoremstyle{definition}
\newtheorem{remark}[theorem]{Remark}
\newtheorem{example}[theorem]{Example}
\theoremstyle{plain}
\newtheorem{introtheorem}{Theorem}
\renewcommand{\theintrotheorem}{\Alph{introtheorem}}
\setlength{\parskip}{2pt}
\setlength{\emergencystretch}{3em}
\renewcommand{\arraystretch}{1.12}
\setcounter{MaxMatrixCols}{20}
\setcounter{secnumdepth}{2}
\setcounter{tocdepth}{1}
\setlist{itemsep=0pt,topsep=2pt,parsep=0pt,partopsep=0pt}
\lstset{basicstyle=\ttfamily\fontsize{9}{10.5}\selectfont,breaklines=true,breakatwhitespace=false,columns=fullflexible,keepspaces=true,showstringspaces=false,aboveskip=0.4em,belowskip=0.4em}
\makeatletter
\renewcommand\subsection{\@startsection{subsection}{2}{\z@}  {\medskipamount}  {-1em}  {\normalfont\normalsize\bfseries}}
\makeatother
\pagestyle{plain}
\title{Report on regular Hadamard matrices of order 11}
\author{Teo and GPT}
\date{}

\begin{document}
\pagenumbering{arabic}
\maketitle
\begin{abstract}
We report on the classification of regular complex Hadamard matrices of order $11$. Write $\mathrm A=11$, $\mathrm B=7+2+2$, $\mathrm C=5+3+3$, $\mathrm D=5+2+2+2$, $\mathrm E=3+3+3+2$ and $\mathrm F=3+2+2+2+2$ for the possible prime-cycle profiles of a row quotient. A regular triple is homogeneous when its three row quotients admit the same profile. We obtain finite affine covers, up to ordinary equivalence, of all homogeneous $3\times11$ partial matrices, with sharp maximum dimensions $0,3,4,3,4,6$ for the branches $\mathrm{AAA},\mathrm{BBB},\mathrm{CCC},\mathrm{DDD},\mathrm{EEE},\mathrm{FFF}$, respectively. The six-dimensional families have two explicit normal forms and generically admit no regular fourth row, with this being our main theorem here. The proofs retain arbitrary unit phases and use exact, reproducible finite computations. The full order-$11$ classification remains open. We also obtain further extension results, notably in the $\mathrm{DDD}$ case.
\end{abstract}
\begingroup
\setlength{\parskip}{0pt}
\makeatletter
\patchcmd{\l@section}{\addvspace{1.0em \@plus\p@}}{\addvspace{0pt}}{}{}
\makeatother
\hypersetup{linkcolor=magenta}
\tableofcontents
\endgroup
\section*{Introduction}
\addcontentsline{toc}{section}{Introduction}

In \cite{regular78} we classified regular complex Hadamard matrices at orders $7$ and $8$. The present report concerns the next prime order, $11$. The two projects began together, with the expectation that primality would make the order-$11$ problem comparable in difficulty to the order-$8$ problem, although for different reasons. At order $8$ the successive partial-matrix calculations eventually gave a complete classification. At order $11$ they produced considerably more continuous families and exceptional parameter sets than we had anticipated. We record the results obtained, together with their proofs and computational certificates, and leave the full classification open.

A complex Hadamard matrix has entries of modulus one and orthogonal rows. It is regular when each row quotient splits into rotated prime cycles, that is, rotated copies of the complete set of $p$th roots of unity for primes $p$. This regularity condition was introduced in \cite{bbs} in the study of the quantum permutation algebras associated to Hadamard matrices. The same definition applies to a rectangular matrix with orthogonal rows. We use ordinary equivalence: row and column permutations and multiplication of individual rows and columns by unit scalars. Conjugation and transposition are separate operations. General background can be found in \cite{inv,tz1}; the relation with vanishing sums of roots of unity is discussed in \cite{lle}.

For a row quotient of length $11$, the possible cycle profiles are
\[
\begin{array}{lll}
\mathrm{A}=11,&\mathrm{B}=7+2+2,&\mathrm{C}=5+3+3,\\
\mathrm{D}=5+2+2+2,&\mathrm{E}=3+3+3+2,&\mathrm{F}=3+2+2+2+2.
\end{array}
\]
A regular triple has three row quotients. We call it homogeneous if all three admit the same profile. Thus $\mathrm{DDD}$, for instance, means that each quotient has a chosen pentagon and three chosen antipodal pairs. A quotient can admit more than one profile, and the supports of its cycles are part of a presentation rather than an invariant of the matrix.

Fixing these supports and their root labels gives integer linear congruences in the phases. Smith normal form describes the solution set as finitely many tori translated by torsion points. After removing row and column phases, their entries are roots of unity times Laurent monomials in free unit parameters. We call these affine families. Our first result gives a complete finite affine cover of every homogeneous triple branch and the sharp dimensions of these covers.

\par\addvspace{\medskipamount}\noindent
\begin{minipage}{\linewidth}
\begin{introtheorem}\label{intro:homogeneous}
Every homogeneous regular $3\times11$ partial Hadamard matrix is ordinarily equivalent to a member of a finite affine cover. The largest dimensions in the six branches are
\[
\begin{array}{c|rrrrrr}
\text{chosen profile}&\mathrm{AAA}&\mathrm{BBB}&\mathrm{CCC}&\mathrm{DDD}&\mathrm{EEE}&\mathrm{FFF}\\\hline
\text{largest dimension}&0&3&4&3&4&6,
\end{array}
\]
and every bound is attained. The six-dimensional chosen presentations have two explicit $\mathrm{FFF}$ normal forms. In each of these families, outside an explicit finite union of proper monomial loci, there is no regular fourth row.
\end{introtheorem}
\end{minipage}
\par\addvspace{\medskipamount}

The two normal forms make the dimension-six phenomenon quite concrete. Put
\[
T_0(a)=\begin{pmatrix}1&1&1&1\\1&-1&a&-a\\1&1&-1&-1\end{pmatrix},
\qquad
T_1(a)=\begin{pmatrix}1&1&1&1\\1&-1&a&-a\\1&-1&-a&a\end{pmatrix}.
\]
With $F_3=(w^{ij})_{i,j=0}^2$ and $w=e^{2\pi i/3}$, they are
\[
H_h(a,b,s,t,p,q)=
\left[\ F_3\ \middle|\ \operatorname{diag}(1,s,t)T_0(a)
\ \middle|\ \operatorname{diag}(1,p,q)T_h(b)\ \right],\qquad h=0,1.
\]
All six parameters range independently over $\mathbb T$. Every quotient consists of a triangle and four antipodal pairs. The assertion about a fourth row allows every new cycle profile and arbitrary unit phases. It therefore says more than the failure of an extension preserving the displayed blocks. The exceptional set is sufficient for the exclusion and is not claimed to be minimal. Nor does the theorem identify different points of the affine covers up to ordinary equivalence.

The other homogeneous branches have different extension behaviour. The maximal-dimensional $\mathrm{EEE}$ families also stop generically at three rows. Common-support $\mathrm{BBB}$ families have generic maximum three or four, according to their seventh-root labels. The maximal-dimensional $\mathrm{CCC}$ families have generic maximum five; their fourth-row circles have a six-cycle compatibility graph. These results concern the stated families on explicit open dense sets. The exceptional parameters and the remaining lower-dimensional presentations are part of the open problem.

The next result is uniform in its parameters. Let $F_5=(w^{ij})_{i,j=0}^4$, with $w=e^{2\pi i/5}$. Joining its rows to five rows of a quaternary Hadamard matrix of order six gives a regular five-row matrix of width eleven. The two blocks can be rotated independently in each row. A new row preserving their pentagon support is ruled out immediately by the five rows of $F_5$; the issue is to exclude every other cycle decomposition.

\par\addvspace{\medskipamount}\noindent
\begin{minipage}{\linewidth}
\begin{introtheorem}\label{intro:five}
Put
\[
V=\begin{pmatrix}
1&1&1&1&1&1\\
1&-1&-i&i&i&-i\\
1&-i&-1&-i&i&i\\
1&i&-i&-1&-i&i\\
1&i&i&-i&-1&-i
\end{pmatrix},
\qquad
H(t)=\left[\ F_5\ \middle|\ \operatorname{diag}(1,t_1,t_2,t_3,t_4)V\ \right].
\]
For every $t\in\mathbb T^4$, the matrix $H(t)$ is regular and has no regular sixth row, whatever the profiles of the new quotients.
\end{introtheorem}
\end{minipage}
\par\addvspace{\medskipamount}

Thus the whole four-parameter family is maximal at five rows. Regularity is essential: we exhibit a parameter value with an orthogonal unimodular sixth row whose quotient with one old row is not regular. A second family $G(a,b,c)$, with three different pentagon supports among its first three rows, is also uniformly maximal at five rows. Its seed has minimal Butson level $60$, and its generic members are inequivalent to members of $H(t)$.

We pursue the extension problem more systematically in the $\mathrm{DDD}$ case. The decomposition $\mathrm D=5+2+2+2$ combines a pentagon with three antipodal pairs, whose phase relations avoid the triangle orientations occurring in $\mathrm{CCC}$. This initially made $\mathrm{DDD}$ an attractive case for explicit calculations. The subsequent construction of continuous five-row families provided a further reason to investigate it: these families exhibit substantial extensions beyond the triple classification. We begin with the triples arising from zero-dimensional presentations.

The $\mathrm{DDD}$ census has $452$ feasible classes of chosen supports. Of these, $218$ have dimension zero, $190$ have dimension one, $41$ have dimension two, and $3$ have dimension three. This does not count matrix-equivalence classes: all root labels and torsion components must still be expanded. For the zero-dimensional presentations we carry out this expansion and solve the entire extension problem.

\par\addvspace{\medskipamount}\noindent
\begin{minipage}{\linewidth}
\begin{introtheorem}\label{intro:finite}
Let $Z_{\mathrm D}$ be the set of ordinary equivalence classes of regular triples admitting a zero-dimensional chosen $\mathrm{DDD}$ presentation. It has $54836$ elements, all with representatives over $\mu_{40}$. Their maximum regular extension sizes are
\[
\begin{array}{c|rrr}
\text{maximum number of rows}&3&4&5\\\hline
\text{classes in }Z_{\mathrm D}&4000&44292&6544,
\end{array}
\]
where added rows may have arbitrary unit entries. In particular, no member of $Z_{\mathrm D}$ belongs to a regular square matrix of order $11$.
\end{introtheorem}
\end{minipage}
\par\addvspace{\medskipamount}

A zero-dimensional chosen presentation need not be isolated in the full regular locus. The set $Z_{\mathrm D}$ includes every triple with such a presentation, even if it also lies in a continuous family with other chosen supports. For continuous $\mathrm{DDD}$ presentations we obtain complete exclusions for some support classes, and generic or restricted exclusions for others. We retain these calculations in the report, with their parameter domains stated explicitly. They show both the effectiveness of the method and the amount of work hidden by a small number of support classes.

The Fourier branch is more rigid. A regular square of order $11$ containing an $11$-cycle quotient is equivalent to $F_{11}$. We include the argument and the finite compatibility calculation, since they also supply exclusions used throughout the report. The broader prime-order questions arising from this branch were subsequently developed in \cite{prime}. None of these results reduces every regular square of order $11$ to the homogeneous cases: the mixed triples and their completions remain a separate obstacle.

The computational proofs keep unit phases as variables. Finite enumeration concerns cycle supports, root labels and exact integer equations, rather than an arbitrary bounded grid of roots of unity. Extensions are covered by phase graphs, and their simultaneous compatibility is tested by exact cyclotomic identities, integer lattice consistency and certified exclusions. The explicitly identified level-$120$ result is a restricted Butson calculation. Each computation is accompanied by its finite input, coverage argument and reproducible checks.

We still regard a full classification by these methods as plausible. The present calculations do not, however, establish that the remaining work is short, and its scale exceeded what we could reasonably complete with our 2026 workflow and resources. This report records a stopping point with definite theorems and explicit outstanding cases, so that the classification can be resumed. Section~\ref{sec:conventions} fixes the conventions and Section~\ref{sec:fourier} treats the Fourier branch. Sections~\ref{sec:heptagons_pentagons} and~\ref{sec:triangles_pairs} treat the other homogeneous cases, culminating in the combined theorem. The uniform five-row theorems follow in Section~\ref{sec:order11part2}. Section~\ref{sec:extensionsDDD} studies extensions of $\mathrm{DDD}$ triples, Section~\ref{sec:mixed} records the mixed cases, and Section~\ref{sec:N11classification} describes the remaining problems and the code.


\section{Regular partial Hadamard matrices}\label{sec:order11part1}\label{sec:conventions}
Write $\mathbb T=\{z\in\mathbb C:|z|=1\}$ and $\mu_m=\{z\in\mathbb T:z^m=1\}$. A rotated prime cycle is the multiset $q\mu_p$, where $q\in\mathbb T$ and $p$ is prime. Such a cycle has sum zero.

\par\addvspace{\medskipamount}\noindent
\begin{minipage}{\linewidth}
\begin{definition}
A partial Hadamard matrix $H\in M_{K\times N}(\mathbb C)$ has entries of modulus one and satisfies $HH^*=NI_K$. We use $\langle x,y\rangle=\sum_j x_j\bar y_j$. Its row quotients are the vectors $(H_{ij}\bar H_{kj})_{j=0}^{N-1}$ for $i\ne k$. It is regular if each row quotient is a disjoint union, with multiplicities, of rotated prime cycles.
\end{definition}
\end{minipage}
\par\addvspace{\medskipamount}

Throughout, ordinary equivalence allows row and column permutations and multiplication of individual rows and columns by elements of $\mathbb T$. It does not include conjugation. A matrix is dephased if its first row and first column consist of ones. A partial Butson matrix~\cite{but} over $\mu_m$ has all its entries in $\mu_m$; in the square case of order $N$ we write $BH(N,m)$. The minimal Butson level of an equivalence class is the least such $m$, when it exists.

We use the cross ratios $H_{ij}H_{k\ell}\bar H_{i\ell}\bar H_{kj}$. Their multiset is invariant under ordinary equivalence, and each entry of a dephased matrix is one of these ratios. Thus a cross ratio of order $m$ forces every Butson representative to have level divisible by $m$. We write $F_N=(w^{ij})_{i,j=0}^{N-1}$ for the Fourier matrix, with $w=e^{2\pi i/N}$. Throughout the paper, $w$ denotes a root of unity whose order is specified locally; other required roots are written as powers of $w$, sometimes using $\omega=w^k$.

At width $11$ the possible profiles are
$$\begin{array}{lll}\mathrm{A}=11,&\mathrm{B}=7+2+2,&\mathrm{C}=5+3+3,\\\mathrm{D}=5+2+2+2,&\mathrm{E}=3+3+3+2,&\mathrm{F}=3+2+2+2+2.\end{array}$$
For $3$ rows we write $(1;u;v)$, where $1$ is the all-ones row, and $z=u/v$, with division taken entrywise. A label such as $\mathrm{DDD}$ specifies a chosen cycle decomposition for each of $u,v,z$. Profiles need not be unique. A chosen support records the column positions of these cycles; its classification is distinct from ordinary matrix equivalence. Different support families can meet, and a zero-dimensional chosen presentation need not be isolated in the full regular locus.

An affine family has entries which are fixed roots of unity times Laurent monomials in free unit parameters. Setting those parameters equal to one gives its Butson seed. The following phase-system argument, as in \cite{regular78}, explains why such families arise.

\par\addvspace{\medskipamount}\noindent
\begin{minipage}{\linewidth}
\begin{proposition}\label{n11_phase:affine}
For fixed cycle supports and root labels, the solution set of the phase equations is a finite union of tori translated by torsion points. Consequently regular partial Hadamard matrices of fixed size admit a finite cover by affine families with Butson seeds.
\end{proposition}
\end{minipage}
\par\addvspace{\medskipamount}

\begin{proof}
The entry phases, together with any auxiliary cycle phases, satisfy an integer system
$$Ax=b\pmod1,\qquad x\in(\mathbb R/\mathbb Z)^n,\quad b\in\mathbb Q^m/\mathbb Z^m.$$
Choose unimodular matrices $U,V$ with $UAV=\operatorname{diag}(d_1,\ldots,d_r,0)$, $d_i>0$. Consistency requires $(Ub)_i=0$ modulo $1$ in the zero rows; all solutions are
$$x=Vy,\qquad y_i=\frac{(Ub)_i+k_i}{d_i}\ (i\le r),\quad 0\le k_i<d_i,\qquad y_i=t_i\ (i>r).$$
Each choice of the integers $k_i$ gives a torus with a rational base point and a primitive integer parameter basis. Projecting to the entry phases preserves the asserted form. There are finitely many decorated supports, so their union gives the required cover.
\end{proof}

Smith and Hermite normal forms can be computed over the integers as in \cite{KB}. The same argument applies when old entries are prescribed monomials in free parameters: include those parameters among the phase variables. Every torsion branch must be retained.

Our finite counts come from exhaustive integer calculations. The proofs specify the finite input, the equations and the coverage argument; the accompanying code supplies the complete lists and checks. Modular nonvanishing is used only for exclusion. Positive regularity assertions use exact prime-cycle covers. Except in the explicitly stated level-$120$ calculation, all phase variables range over the full unit circle.

\section{The Fourier branch}\label{sec:fourier}
\subsection{The branch containing an 11-cycle quotient}\label{sec:branch_A}

The Fourier results below were obtained in this project and led to the prime-order study \cite{prime}. We retain the proofs because the triple-profile restriction is used in the extension arguments.

Put $w=e^{2\pi i/11}$. An $11$-cycle quotient allows the normalization $(1;(w^j)_j;\ldots)$. Our aim is the following square classification.

\par\addvspace{\medskipamount}\noindent
\begin{minipage}{\linewidth}
\begin{theorem}\label{n11_A_completion:main}
Let $H$ be a regular complex Hadamard matrix of order $11$. If some quotient of 2 rows is a single $11$-cycle, then $H$ is equivalent to $F_{11}=(w^{ab})_{a,b=0}^{10}$.
\end{theorem}
\end{minipage}
\par\addvspace{\medskipamount}

We first determine the possible $3$-row patterns. A rigidity argument then forces every quotient in a square completion to be an $11$-cycle, leaving a finite Fourier compatibility calculation. The proof of the announced theorem is completed at the end of this subsection.

Using the equivalence operations above, an $11$-cycle quotient allows us to write
$$H=(1;u;v),\qquad u_j=w^{j},\quad j\in\mathbb{Z}/11\mathbb{Z}.$$
Put $y_j=v_j/u_j$. Thus the 3 quotient profiles are those of $u,y,v$, and $u_j y_j=v_j$.

\par\addvspace{\medskipamount}\noindent
\begin{minipage}{\linewidth}
\begin{lemma}\label{n11_A_branch:overlap}
If 2 of the quotients have profile $\mathrm{A}$, so does the third. If $y,v$ contain no $11$-cycles, a cycle in $y$ and a cycle in $v$ meet in at most one column position.
\end{lemma}
\end{minipage}
\par\addvspace{\medskipamount}

\begin{proof}
2 $11$-cycle quotient vectors have all their entries in rotated copies of $\mu_{11}$. Their ratio also lies in a rotated $\mu_{11}$. By regularity, this third quotient decomposes into prime cycles. Each such cycle must have length $11$, so the quotient is a single $11$-cycle.

For the second assertion, suppose columns $j\ne k$ belong to the same $p$-cycle of $y$ and the same $q$-cycle of $v$. Then
$$w^{j-k}=\frac{v_j/v_k}{y_j/y_k}\in\mu_{\operatorname{lcm}(p,q)}.$$
As $p,q\ne11$, this contradicts $j\ne k$ modulo $11$.
\end{proof}

The overlap restriction leaves only a short list of profile patterns. The labels on their support graphs also give a finite affine description.

\par\addvspace{\medskipamount}\noindent
\begin{minipage}{\linewidth}
\begin{theorem}\label{n11_A_branch:profiles}
The possible chosen profile triples of regular $3\times11$ partial Hadamard matrices with an $11$-cycle quotient are precisely:
$$\mathrm{AAA},\qquad \mathrm{AEE},\qquad \mathrm{AEF},\qquad \mathrm{AFF}.$$
All 4 occur. This branch admits an exhaustive finite parametrization by fixed matrices and affine one-parameter families. The parametrization below is redundant; it is not a list of inequivalent matrices or families.
\end{theorem}
\end{minipage}
\par\addvspace{\medskipamount}

\begin{proof}
If either remaining quotient has profile $\mathrm{A}$, Lemma~\ref{n11_A_branch:overlap} gives $\mathrm{AAA}$. Suppose now that neither does. Make a bipartite graph whose left vertices are the cycles of $y$, whose right vertices are the cycles of $v$, and whose edges are the 11 columns. The lemma says that it is simple. Its degrees are the respective cycle lengths. Every vertex has degree at least 2. A $7$-cycle cannot occur, since there are at most 5 vertices on the other side. A $5$-cycle requires 5 vertices on the other side, leaving only $\mathrm{ACF}$ and $\mathrm{ADF}$ as possible cases involving a $5$-cycle.

These 2 cases are also impossible. In either graph, the $5$-cycle is the unique vertex of degree $5$; all other vertices have degree $2$ or $3$. Along graph edges, ratios of the cycle phases are products of roots of orders dividing $330$. Within a connected component, factor out one common phase and project these roots onto their $\mu_5$ component. After removing the $5$-cycle vertex, this component is constant on each remaining connected component. Its 5 neighbors must have 5 different $\mu_5$ components, because the 5 entries in the $5$-cycle do. They must therefore belong to 5 different connected components after removal. Each neighbor has another left neighbor, because its original degree is at least 2. This requires at least 5 other left vertices. There are only 2 in profile $\mathrm{C}$ and 3 in profile $\mathrm{D}$, a contradiction. This proves the profile exclusion used in Theorem~\ref{n11_A_branch:profiles}.

\medskip
For the graph parametrization put $w=e^{\pi i/33}$ and $\omega=w^{11}=e^{\pi i/3}$. Thus $w^6$ is the primitive $11$th root used above. For $\mathrm{AEE},\mathrm{AEF},\mathrm{AFF}$, all the other cycles have length $2$ or $3$. In each connected graph component, the edge equations force the cycle phases, relative to one common phase, to lie in $\mu_{66}=\mu_{11}\mu_6$. Consequently the data are:
\par\noindent
\begin{minipage}{\linewidth}
\begin{enumerate}
\item A simple bipartite graph with 11 edges and the prescribed degrees.
\item An edge color $r_e\in\mathbb{Z}/6\mathbb{Z}$ such that the colors at a degree-$p$ vertex form one rotated $p$-cycle in $\mu_6$.
\item Labels $b_i,c_k\in\mathbb{Z}/11\mathbb{Z}$ on the left and right vertices, with the following multiset identity:
$$\{c_k-b_i:(i,k)\text{ is an edge}\}=\mathbb{Z}/11\mathbb{Z}.$$
\item One phase $t_s\in\mathbb{T}$ for each connected component.
\end{enumerate}
\end{minipage}\par
For the edge $e=(i,k)$ of color $r_e$ in component $s$, set
$$j=c_k-b_i,\qquad y_j=t_sw^{6b_i}\omega^{r_e},\qquad v_j=t_sw^{6c_k}\omega^{r_e}.$$
The differences enumerate the columns exactly once. Every vertex supplies its required prime cycle, and $v_j/u_j=y_j$. These conditions are sufficient as well as necessary. A common factor of all the $t_s$ is removed by a row factor. Thus a graph with $c$ connected components supplies $c-1$ continuous parameters.

\medskip
Enumerate the simple bipartite graphs with the prescribed degrees, then all cycle colors and vertex labels satisfying these conditions. The color enumeration removes a common rotation in $\mu_6$ and orders cycles of equal length before matching their entries; the removed rotation is absorbed in the component phases. After graph isomorphism and the label normalization specified below, the retained presentations are listed in the next table. Here $123$ denotes the neighbor set $\{1,2,3\}$ of a left vertex. Rows of a graph are listed in degree order. Colors and labels are retained finite presentations, not equivalence-class counts.
\begin{center}
\small
\begin{tabular}{@{}clrrr@{}}
\toprule
Profiles & Left neighbor sets & Components & Colors & Labels\\
\midrule
$\mathrm{EE}$ & $(123,123,124,34)$ & 1 & 4 & 40\\
$\mathrm{EF}$ & $(123,234,145,15)$ & 1 & 2 & 368\\
$\mathrm{FF}$ & $(123,12,13,45,45)$ & 2 & 6 & 480\\
$\mathrm{FF}$ & $(123,12,14,35,45)$ & 1 & 2 & 3744\\
$\mathrm{FF}$ & $(123,14,24,15,35)$ & 1 & 1 & 2592\\
$\mathrm{FF}$ & $(234,12,13,15,45)$ & 1 & 1 & 2880\\
\bottomrule
\end{tabular}
\end{center}
The label normalization is precise: choose the first right vertex $k$, and its first 2 left neighbors $i,j$, and impose
$$b_i=0,\qquad b_j=1,\qquad c_k=0.$$
The edge differences force $b_i\ne b_j$, so independent translations on the 2 sides followed by a common nonzero field scaling achieve this. In every other connected component, set its first vertex label to zero; its removed common label is absorbed into $t_s$. To recover all presentations, restore the 10 common nonzero field scalings. They are not silently identified as Hadamard equivalences.

Only the third row of the table has more than one connected component. Its components have 7 and 4 edges. It gives affine one-parameter families; all the other graph presentations have no continuous parameter. The $\mathrm{AAA}$ branch is the finite complete-mapping problem used again in the square-completion calculation below: after normalization it consists of permutations $\pi$ with $\pi(0)=0$ and $j+\pi(j)$ also a permutation. There are $3441$ such normalized maps.
\end{proof}

\medskip
With the same $w=e^{\pi i/33}$, let $|q|=1$. An explicit family from the disconnected $\mathrm{FF}$ graph is
$$H(q)_{0j}=1,\qquad H(q)_{1j}=w^{6j},\qquad H(q)_{2j}=w^{a_j}q^{d_j},\quad 0\le j\le10,$$
where
\begin{align*}
(a_j)&=(0,23,56,40,2,41,22,35,8,7,44),\\
(d_j)&=(0,0,0,0,1,1,0,1,1,0,0).
\end{align*}
Thus the $q$ factors occur precisely in columns $4,5,7,8$, with columns numbered from zero. Writing the third row as $v$ and the second as $u$, explicit cycle supports are
$$\begin{array}{c|l}
v & \{0,6,10\};\ \{1,2\};\ \{3,9\};\ \{4,7\};\ \{5,8\}\\[2pt]
v/u & \{0,2,3\};\ \{1,10\};\ \{6,9\};\ \{4,5\};\ \{7,8\}.
\end{array}$$
In each line, the first block is a $3$-cycle and the remaining 4 blocks are antipodal pairs. On each block the power of $q$ is constant. These supports prove regularity for every $q\in\mathbb{T}$, and the middle row is the full $11$-cycle. In particular $H(q)H(q)^*=11I_3$ identically in $q$. The code independently verifies every Laurent coefficient of each row scalar product by reduction modulo $\Phi_{66}$ over the integers. The cross ratio $qw^{2}$ shows that each equivalence class meets this family in finitely many parameter values \cite{ha1,tz1}.

\medskip
We now return to $w=e^{2\pi i/11}$ and consider a regular square matrix with a distinguished $11$-cycle quotient. Normalize its distinguished rows to $1$ and $u=(w^{j})_{j=0}^{10}$. For any other row $v$, put $y=v/u$. Lemma~\ref{n11_A_branch:overlap} and the profile exclusion above show that either both $y,v$ are $11$-cycles, or both have profiles $\mathrm{E}$ or $\mathrm{F}$. The same overlap graph therefore controls every row of the square matrix.

In the latter case their graph has at most 10 vertices and minimum degree 2. Each connected component has at least 4 vertices, so there are at most 2 components. Within a component, ratios of cycle phases belong to $\mu_{66}$. Consequently the entries of $v$ lie in the union of at most 2 rotated copies of $\mu_{66}$. The same conclusion holds for an $\mathrm{A}$ row, which needs just one copy.

\par\addvspace{\medskipamount}\noindent
\begin{minipage}{\linewidth}
\begin{lemma}\label{n11_A_completion:smallprimes}
Every row quotient of $H$ decomposes using only cycles of lengths $2$, $3$, and $11$.
\end{lemma}
\end{minipage}
\par\addvspace{\medskipamount}

\begin{proof}
Each row lies in at most 2 cosets of $\mu_{66}$, so the quotient of any 2 rows lies in at most 4 such cosets. A rotated $5$-cycle or $7$-cycle intersects a single coset of $\mu_{66}$ in at most one value: the ratio of 2 distinct common values would lie both in $\mu_{66}$ and in $\mu_5$ or $\mu_7$. 4 cosets cannot contain either cycle. Regularity and the list of prime partitions of $11$ give the result.
\end{proof}

This argument uses only the graph size and the phase relations. It does not require the detailed enumeration of the 6 3-row support graphs.

\medskip
We first record an elementary fact about flat matrices; here flat means that all entries have modulus one.

\par\addvspace{\medskipamount}\noindent
\begin{minipage}{\linewidth}
\begin{lemma}\label{n11_A_completion:blocks}
Let $K$ be a flat $n\times n$ matrix whose rows are pairwise orthogonal or proportional, and whose columns have the same property. All row proportionality classes and all column proportionality classes have the same cardinality $d$. In particular $\operatorname{rank}(K)=n/d$ divides $n$.
\end{lemma}
\end{minipage}
\par\addvspace{\medskipamount}

\begin{proof}
The number of row classes and the number of column classes both equal the rank, say $k$. Remove the proportionality phases and retain one row and column from each class to obtain a flat invertible $k\times k$ matrix $B$. Let $R=\operatorname{diag}(r_i)$ and $C=\operatorname{diag}(c_j)$ record the row and column class sizes. Orthogonality of distinct classes gives
$$\mathrm{BCB}^*=nI_k,\qquad B^*RB=nI_k.$$
Thus $CB^*=B^*R$, or $RB=BC$. Since all entries of $B$ are nonzero, $r_i=c_j$ for every $i,j$. Their common value is $d$, and $n=kd$.
\end{proof}

We use this rank argument after retaining the entry phases formally. Primality then forces the proportionality classes to collapse.

\par\addvspace{\medskipamount}\noindent
\begin{minipage}{\linewidth}
\begin{proposition}\label{n11_A_completion:rigidity}
Let $p>3$ be prime. Suppose a complex Hadamard matrix of order $p$ has every row quotient decomposable using only $2$-, $3$-, and $p$-cycles. If one row quotient is a $p$-cycle, then every row quotient is a $p$-cycle. In particular the matrix is equivalent to a Butson matrix $BH(p,p)$.
\end{proposition}
\end{minipage}
\par\addvspace{\medskipamount}

\begin{proof}
Let $G\subset\mathbb{T}$ be the finitely generated group generated by all matrix entries and by $\mu_6$ and $\mu_p$. For this proof put $w=e^{2\pi i/(3p)}$, $\omega=w^p=e^{2\pi i/3}$ and $O=\mathbb Z[\omega]$. We retain the phases formally, imposing only the ordinary arithmetic of the sixth roots, in the commutative ring
$$R=\mathbb{Z}[G]\otimes_{\mathbb{Z}[\mu_6]}O.$$
Here the second factor uses the usual inclusion of $\mu_6$ into $O$. The ring $R$ is a free $O$-module on representatives of $G/\mu_6$. Its involution sends each group element to its inverse and conjugates its coefficient. Write $[z]$ for the formal image of $z\in G$.

For $f=\sum_g a_g[g]$ in this basis put $\|f\|_1=\sum_g|a_g|$. Changing representatives multiplies coefficients by sixth roots, so the norm is independent of those choices. Every nonzero Eisenstein integer has absolute value at least one, since
$$|a+b\omega|^2=a^2-ab+b^2\in\mathbb{Z}_{>0}.$$

The root $w^3$ generates $\mu_p$. Set
$$S=1+[w^3]+\cdots+[w^3]^{p-1},\qquad e=S/p\in R\otimes\mathbb{Q}.$$
Then $e$ is a central idempotent. Let $M$ be the formal matrix obtained by replacing each entry $h_{ij}$ by $[h_{ij}]$. Every $2$- or $3$-cycle sum vanishes in $R$, and a $p$-cycle sum is a monomial times $S$. Therefore
$$(1-e)MM^*=p(1-e)I_p.$$
In the commutative ring with identity $1-e$, a square matrix with a right inverse also has a left inverse, by the determinant and adjugate identity. It follows that
$$(1-e)M^*M=p(1-e)I_p.$$
Hence every off-diagonal column scalar product $f$ satisfies $f=ef$, and therefore $[w^3]f=f$.

Multiplication by $w^3$ acts freely on $G/\mu_6$, with orbits of length $p$. On each nonzero orbit the coefficients of $f$ have the same absolute value, at least one. Each such orbit thus contributes at least $p$ to $\|f\|_1$. On the other hand, $f$ is a sum of precisely $p$ group monomials, so $\|f\|_1\le p$. Consequently either $f=0$, or there is one orbit and every coefficient has absolute value one. In the latter case there are $p$ nonzero cosets and exactly $p$ original summands. There is one summand in each coset, and invariance under multiplication by $w^3$ identifies these summands as a single rotated $p$-cycle. This is an assertion about the original column quotient multiset, not just its evaluated sum.

Choose a unitary character $\chi:G\to\mathbb{T}$ which is the identity on $\mu_6$ and trivial on $\mu_p$. Such a character exists: the torsion group of $G$ is finite cyclic, and the Chinese remainder theorem gives an exponent congruent to $1$ modulo $6$ and to $0$ modulo $p$; choose arbitrary unitary images for the free generators. Apply $\chi$ entrywise to obtain $K$.

Distinct rows of $K$ are orthogonal when their original quotient is a sum of $2$- and $3$-cycles, and proportional when it is a $p$-cycle. The formal column argument proves the same orthogonal-or-proportional property for the columns of $K$. The given $p$-cycle row quotient becomes a proportional row pair, so $\operatorname{rank}(K)<p$. By Lemma~\ref{n11_A_completion:blocks}, the rank divides the prime $p$; it is therefore one. Thus every pair of rows of $K$ is proportional. No original quotient can have been a sum of $2$- and $3$-cycles, since that would have made this pair orthogonal. All original row quotients are $p$-cycles. After dephasing, every entry is a $p$th root of unity.
\end{proof}

\begin{proof}[Proof of Theorem~\ref{n11_A_completion:main}]
Apply Proposition~\ref{n11_A_completion:rigidity} with $p=11$, using Lemma~\ref{n11_A_completion:smallprimes}. It excludes every mixed regular square completion in the distinguished-$11$-cycle branch, including all values of the free phases in the previously obtained 3-row affine families.

\medskip\label{sec:A_fourier}
It remains to classify $BH(11,11)$. This uniqueness is also included in the small-prime Butson classification of \cite{hkm}; we retain the direct compatibility calculation, which gives information about partial extensions. For the related circulant Butson problem see \cite{hsc}. Here again put $w=e^{2\pi i/11}$. After dephasing and permuting columns, its first 2 rows are $1$ and $(w^{j})_{j=0}^{10}$. Any further dephased row is $(w^{a_j})$, where
$$a_0=0,\qquad (a_j)_{j=0}^{10}\text{ and }(a_j-j)_{j=0}^{10} \text{ are permutations of }\mathbb{Z}/11\mathbb{Z}.$$
Indeed, a vanishing sum of 11 $11$th roots has one of each root: its exponent-count polynomial has degree at most 10 and is a multiple of $\Phi_{11}=1+x+\cdots+x^{10}$.

There are exactly $3441$ such candidate rows. Join 2 candidates $a,b$ when $(a_j-b_j)_{j=0}^{10}$ is a permutation. A square completion is exactly a clique of size 9 in this compatibility graph. The graph has $2016$ edges, with degree distribution
$$\begin{array}{c|rrrr}
\text{degree}&0&3&8&162\\\hline
\text{number of vertices}&2640&660&135&6.
\end{array}$$
The 9 Fourier candidates are $a_j=kj$, $k=2,\ldots,10$. 3 have degree 8 and 6 have degree $162$; they form a complete graph $K_9$. All other vertices are excluded from a $9$-clique as follows:
\par\noindent
\begin{minipage}{\linewidth}
\begin{itemize}
\item $3300$ have degree less than 8.
\item The remaining $132$ have degree exactly 8, but their 8 neighbors are not pairwise adjacent. For each one, the certificate lists an explicit nonadjacent pair of neighbors. A $9$-clique containing that vertex would require its entire neighborhood to be complete.
\end{itemize}
\end{minipage}\par
Thus the Fourier $K_9$ is the unique square completion. This proves Theorem~\ref{n11_A_completion:main}.
\end{proof}

The same Python verification gives the following additional information about the non-Fourier candidate rows. Here extensions are required to have all their row quotients be $11$-cycles, and counts are normalized candidate rows, not equivalence classes:
\begin{center}
\begin{tabular}{@{}rr@{}}
\toprule
Number of candidates & Largest total row count in an all-$11$-cycle extension\\
\midrule
2640 & 3\\
660 & 4\\
132 & 5\\
\bottomrule
\end{tabular}
\end{center}
In particular, an all-$11$-cycle partial matrix with at least 6 rows is a Fourier submatrix up to the stated normalization and equivalence. The mixed cases need not fail at these row counts; their exclusion above uses the square-matrix rigidity proposition.

\section{The heptagon and pentagon cases}\label{sec:homogeneous}\label{sec:heptagons_pentagons}
The next two sections treat the $5$ homogeneous chosen profiles $\mathrm{BBB},\mathrm{CCC},\mathrm{DDD},\mathrm{EEE},\mathrm{FFF}$. Together with the finite $\mathrm{AAA}$ branch of Section~\ref{sec:fourier}, they give all homogeneous triples. For each profile we describe the supports and phase equations, establish an exhaustive affine cover, and then study the extensions that have been determined. We begin with the heptagon and pentagon cases. Theorem~\ref{report:homogeneous}, at the end of Section~\ref{sec:triangles_pairs}, gathers all six branches into the sharp dimension statement announced in the Introduction.

The dimension is that of the family of matrices after removing row and column phases. Our chosen normalizations leave only finite permutation ambiguities, which do not change it. Counts of support classes, connected phase components and ordinary matrix-equivalence classes will always be identified separately. No disjointness of the resulting affine covers is assumed.

\subsection{The homogeneous BBB case}\label{sec:BBB}

\medskip
We begin the non-$\mathrm{A}$ reductions with the homogeneous type $\mathrm{BBB}$: each of the 3 chosen quotients of a $3\times11$ matrix has profile $7+2+2$. The seventh roots can first be collapsed, leaving a small phase system, and then restored by a separate finite label problem.

Write the 3 rows as $1,u,v$, and put $z=u/v$ entrywise. Denote their chosen support partitions by $P,Q,R$, each consisting of one 7-element block and 2 unordered pairs. Rotate the 7-cycles of $u,v$ to $\mu_7$.

\par\addvspace{\medskipamount}\noindent
\begin{minipage}{\linewidth}
\begin{lemma}\label{n11_noA_BBB:common}
The 3 7-element support blocks have a common column. The 7-cycle of $z$ is therefore also $\mu_7$ after the preceding normalization.
\end{lemma}
\end{minipage}
\par\addvspace{\medskipamount}

\begin{proof}
The 7-blocks of $u$ and $v$ overlap in at least 3 columns. At all these columns $z$ belongs to $\mu_7$. An antipodal pair of $z$ contains at most one such column, since $-1\notin\mu_7$. The 2 pairs cannot cover all 3 columns, so one belongs to the 7-block of $z$. A rotated 7-cycle meeting $\mu_7$ equals $\mu_7$ as a multiset.
\end{proof}

\medskip
The common column fixes the seventh-root normalization. We can now separate the continuous phase equations from the finite root labels.

\par\addvspace{\medskipamount}\noindent
\begin{minipage}{\linewidth}
\begin{theorem}\label{n11_noA_BBB:main}
Every regular $3\times11$ matrix of chosen type $\mathrm{BBB}$ is equivalent to a member of a finite union of affine families of dimension at most 3, each containing a partial Butson matrix over $\mu_{28}$. The bound 3 is attained.

There are 10 relative support types for the first 2 quotients. For these representatives and all $990$ choices of the third quotient support, exactly $2571$ support pairs are feasible. Their seventh-power phase systems have $2776$ connected components that lift to genuine triples, distributed as
$$\begin{array}{c|rrrr}
\text{dimension}&0&1&2&3\\\hline
\text{components}&1788&856&128&4.
\end{array}$$
These are counts in the specified support normalization. In particular, the $1788$ zero-dimensional presentations are not asserted to be isolated points in the full space of partial Hadamard matrices.
\end{theorem}
\end{minipage}
\par\addvspace{\medskipamount}

\begin{proof}
Set $U=u^7$, $V=v^7$, entrywise. Each 7-block becomes 7 copies of $1$; each pair remains antipodal. A partition has the form $P=(P_0,P_1,P_2)$, with $|P_0|=7$, $|P_1|=|P_2|=2$. For an ordered pair block $(a,b)$ with $a<b$, write its entries as $e^{2\pi i x},-e^{2\pi i x}$. Thus
$$\arg U/(2\pi)=\epsilon_P+D_Px,\qquad \arg V/(2\pi)=\epsilon_Q+D_Qx,\quad x\in(\mathbb R/\mathbb{Z})^4,$$
where $\epsilon$ is $0$ on the 7-block and at the first point of each pair, and $1/2$ at the second point. The rows of $D_P$ select $x_1,x_2$ on its 2 pairs, and those of $D_Q$ select $x_3,x_4$. Require $U/V=1$ on $R_0$, and antipodal entries on $R_1,R_2$. These are 9 integer linear congruences
$$A x=b\pmod1,\qquad b\in\{0,1/2\}^9.$$
They are necessary; taking seventh powers alone does not preserve the distinctness of the roots in a 7-cycle.

Fix $P=((0,1,2,3,4,5,6),(7,8),(9,10))$. There are $\binom{11}{4}\,3=990$ partitions of this kind. The stabilizer of $P$ is $S_7\times(C_2\wr S_2)$, and it has 10 orbits on $Q$. Its arbitrary permutations of the 7-block are legitimate here because the root labels will be restored later. Every $R$ is then considered without further reduction.

For $27$ support pairs the common-column condition fails. The remaining $9873$ congruence systems are solved over the full torus, without imposing a root-order bound. Exactly $2995$ are consistent, with $3210$ connected components. The solver writes each component as
$$x=c+Mt,\quad c\in(\tfrac14\mathbb{Z}/\mathbb{Z})^4,\quad M\in M_{4\times d}(\mathbb{Z}),\quad t\in(\mathbb R/\mathbb{Z})^d.$$
Its integer basis is primitive. All components have $d\le3$.

Exact minor calculations bound every nonzero Smith invariant by $2$. Hence every consistent component meets the quarter-root grid. The independent eighth-root-grid check verifies coverage by the reported primitive components, each with $8^d$ grid points and dimension $4-\operatorname{rank}A$. This proves continuous completeness.

Put $w=e^{2\pi i/7}$. For each of the 3 support partitions introduce 9 labels in $\mathbb{F}_7$: one label for each column of its 7-block, and one label shared by the 2 columns of each pair. Call the resulting 11-column vectors $a,b,h$ respectively. Impose
$$a_j-b_j=h_j\quad(j=0,\ldots,10),$$
and require the 7 labels on each 7-block to be a permutation of $\mathbb{F}_7$. These are the entire finite lifting conditions.

For a common column $j_0$ we may normalize $a_{j_0}=b_{j_0}=0$. A common nonzero field scaling then sets $a_{j_1}=1$ at a second column of $P_0$. This scaling is used only for existence testing. It is not assumed to be a Hadamard equivalence: the general normal form permits every label assignment, or equivalently restores all 6 nonzero scalings.

Of the $2995$ consistent support pairs, $424$ fail this finite test. Each has an exact linear certificate: the column equations, the gauge conditions, and the necessary identities that each 7-block's labels sum to zero force 2 labels in a 7-block to be equal. Gaussian elimination is over $\mathbb{F}_7$. For every remaining support pair, a direct label witness is provided and checked entrywise. The search is not relied upon for proving nonexistence.

There are $2571$ liftable support pairs. Removing the impossible pairs leaves the $2776$ phase components in Theorem~\ref{n11_noA_BBB:main}. Lifting is independent of which phase component is chosen. Explicitly, for a component $x=c+Mt$ set
\begin{align*}
u_j(t)&=w^{a_j}\exp\bigl(2\pi i[\epsilon_{P,j}+(D_P(3c+Mt))_j]\bigr),\\
v_j(t)&=w^{b_j}\exp\bigl(2\pi i[\epsilon_{Q,j}+(D_Q(3c+Mt))_j]\bigr).
\end{align*}
The factor $3$ in $3c$ uses $7\cdot3\equiv1\pmod4$. The fourth-root constants have the prescribed collapsed supports, and the seventh-root labels restore each genuine 7-cycle. Pair signs are preserved. The basis equations ensure that every cycle has a constant parameter exponent on its support. Thus every displayed matrix is regular for every $t$, and its seed at $t=0$ has entries in $\mu_{28}$.

Conversely, let $y$ be the vector of the original 4 pair phases, so that $7y=c+Mt$ modulo $1$. Then $y-3c-Mt/7$ has coordinates in $\frac17\mathbb Z/\mathbb Z$. These residues are exactly the finite labels above; rescaling the free variables proves coverage.

\medskip
The maximal dimension is already attained by a simple block construction:
$$\begin{array}{c|c}
\text{columns }j=0,\ldots,6&\text{last 4 columns}\\\hline
1&(1,1,1,1)\\
w^{j}&s(1,-1,q,-q)\\
w^{2j}&t(1,1,-1,-1)
\end{array}
\qquad q,s,t\in\mathbb{T}.$$
The first block supplies the 7-cycles; the last block supplies 2 antipodal pairs for each quotient. The parameters can be recovered from cross ratios, so this is a genuine 3-dimensional family modulo the finite ambiguities of monomial equivalence.
\end{proof}

A less separated example uses $w=e^{\pi i/14}$. Its first row is all ones, and its other rows are $w$ raised entrywise to
$$\begin{pmatrix}
0&4&8&12&16&20&24&5&19&25&11\\
0&8&16&4&12&21&9&20&7&24&23
\end{pmatrix}.$$
With columns numbered from zero, certificates for $u,v,u/v$ are
$$\begin{array}{c|c|c|c}
&\text{7-cycle}&\text{pair}&\text{pair}\\\hline
u&0123456&78&9\,10\\
v&0123479&58&6\,10\\
u/v&012348\,10&57&69
\end{array}$$
Each 7-cycle consists of all seventh roots; each pair has exponent difference $14$ modulo $28$. The cross ratio from the first 2 rows and columns $0,7$ is a primitive $28$th root. Hence this example cannot be equivalent to a partial Butson matrix over $\mu_{14}$.

\medskip
We next test extensions of the separated construction, allowing all cycle supports for the new row. Put $w=e^{2\pi i/7}$. Let $f$ be a permutation of $\mathbb F_7$ such that $f(0)=0$ and $j\mapsto f(j)-j$ is also a permutation. Define
\[
T_f(q,s,t)=
\begin{array}{c|c}
\text{columns }j=0,\ldots,6&\text{last 4 columns}\\\hline
1&(1,1,1,1)\\
w^j&s(1,-1,q,-q)\\
w^{f(j)}&t(1,1,-1,-1)
\end{array},\qquad q,s,t\in\mathbb T.
\]
The $3$ chosen $7$-cycles have the same support. There are $19$ such normalized permutations $f$: the $5$ linear maps $f(j)=aj$, $a=2,\ldots,6$, and $14$ others. These are labelled presentations, not equivalence classes.

\par\addvspace{\medskipamount}\noindent
\begin{minipage}{\linewidth}
\begin{theorem}\label{bbb_common:extension}
There is an explicit proper closed subset $E\subset\mathbb T^3$, given by a union of $3038$ monomial equations, with the following property. For $(q,s,t)\notin E$, every regular fourth row extending $T_f(q,s,t)$, normalized to have first entry $1$, is
\[
x_k(z)=\bigl((w^{kj})_{j=0}^6\mid z(1,-1,-q,q)\bigr),\qquad z\in\mathbb T,
\]
where $f(j)=aj$ is linear and $k\in\{2,3,4,5,6\}\setminus\{a\}$. Consequently the maximum number of rows of a regular extension is $4$ for these $5$ linear patterns and $3$ for the other $14$ patterns.

The equations defining $E$ have the form
$$q^a s^b t^c=e^{\pi i d/105},\qquad
0<\max(|a|,|b|,|c|)\le4,\quad d\in\mathbb Z/210\mathbb Z.$$
Their complete list is supplied with the computation. Every displayed fourth row is valid for all $q,s,t,z\in\mathbb T$; only the exhaustiveness and maximality assertions require $(q,s,t)\notin E$.
\end{theorem}
\end{minipage}
\par\addvspace{\medskipamount}

\begin{proof}
Exhausting the $6!=720$ normalized permutations gives the $19$ possibilities above. Two such permutations $f,g$ are adjacent when $f-g$ is a permutation. The resulting graph is a complete graph on the $5$ linear maps and $14$ isolated vertices. This settles extensions which preserve the $7+4$ split. To exclude other supports we use the following complete phase calculation.

Write the rows of $T_f(q,s,t)$ as $1,u,v$. By Theorem~\ref{n11_A_branch:profiles}, none of the quotients involving a new row $x$ can have profile $\mathrm{A}$, since the old quotients have profile $\mathrm{B}$. Choose a prime-cycle presentation of $x$ with profile $P\in\{\mathrm{B},\mathrm{C},\mathrm{D},\mathrm{E},\mathrm{F}\}$. Enumerate every partition of its $11$ columns into the prescribed cycle sizes, and every root-label permutation on each block, fixing the first label to zero. Equal-size blocks are unordered; their independent phases are unrestricted.

Next enumerate every prime-cycle cover of $x/u$. A cycle imposes ratios between the phases of the initial blocks of $x$. Integrating along a forest expresses these phases as monomials in $q,s,t$, roots in $\mu_{210}$, and one free unit phase per component. An edge closing a cycle imposes a monomial equation on the old parameters. The generic scan keeps identities, rejects inconsistent constant equations, and records every nonconstant equation before rejecting that branch. Repeat the same procedure for $x/v$, for each of the $19$ choices of $f$.

Let $E$ be the union of the equations recorded by this calculation. Outside $E$, every rejected branch is impossible, including branches whose later equations were not explored. Forest integration retains arbitrary unit phases, so this is a full-torus exclusion, not a root-grid test. The recorded equations have the bounds stated in the theorem. Since their integer characters are nonzero, each defines a proper closed subset with empty interior, and the finite union $E$ is proper.

The following table gives the complete enumeration. A row record retains its free phases and monomial coefficients; the last column is summed over the $19$ choices of $f$.
$$\begin{array}{c|r|r|r}
P&\text{initial presentations}&\text{records against }1,u&\text{records against }1,u,v\\\hline
\mathrm{B}&712800&2964&1420\\
\mathrm{C}&443520&0&0\\
\mathrm{D}&166320&0&0\\
\mathrm{E}&123200&0&0\\
\mathrm{F}&34650&59136&0
\end{array}$$
For block sizes $p_1,\ldots,p_r$ with multiplicities $m_p$, the initial count is $11!/(\prod_i p_i\prod_p m_p!)$. Thus all $1480490$ presentations are included. Two implementations, using forest potentials and explicit component potentials respectively, with opposite column and prime orders, agree on the complete sets of pair and triple records. A separate verifier checks exact prime-cycle covers for every retained record and enumerates the $19$ labels independently.

All $1420$ final records belong to the $5$ linear patterns, with $284$ records for each. For a fixed $a$, these consist of the $4$ circles $x_k(z)$ and $280$ fixed specializations of those circles. Entrywise monomial comparison verifies this containment. The $3038$ exceptional equations are retained in \texttt{exceptional\_union.txt} of the \texttt{bbb\_common} component; a line $a\ b\ c\ d$ has exactly the meaning stated above. The union is a sufficient exceptional set and is not asserted to be minimal.

Direct subtraction on the first block and pairing on the last block prove that $x_k(z)$ gives all $6$ quotients of the resulting $4$ rows profile $\mathrm{B}$, for arbitrary parameters. Finally,
$$\langle x_k(z),x_l(z')\rangle=
\begin{cases}
4z\overline{z'},&k\ne l,\\
7+4z\overline{z'},&k=l.
\end{cases}$$
Neither expression vanishes. Hence no two fourth-row candidates can be added together, proving the maximum $4$. For a nonlinear $f$ there is no fourth row at all.
\end{proof}

This theorem concerns the displayed common-support families. It does not settle their exceptional parameters or the other $\mathrm{BBB}$ support presentations. The pair stage also retains many mixed $\mathrm{B}/\mathrm{F}$ records which disappear after imposing the third row; those records must not be discarded in a general study of the $\mathrm{B}$ branch.

\par\addvspace{\medskipamount}\noindent
\begin{minipage}{\linewidth}
\begin{proposition}\label{bbb_common:BDD}
The mixed chosen profile $\mathrm{BDD}$ occurs over $\mu_{70}$.
\end{proposition}
\end{minipage}
\par\addvspace{\medskipamount}

\begin{proof}
Put $w=e^{\pi i/35}$. Take the first row to be $1$ and the other rows $u,x$ to be $w$ raised entrywise to
$$\begin{pmatrix}
0&10&20&30&40&50&60&34&69&13&48\\
0&14&28&42&9&44&25&56&49&60&14
\end{pmatrix}.$$
Numbering columns from zero, exact cycle certificates are
$$\begin{array}{c|c|c}
\text{quotient}&\text{long cycle}&\text{pairs}\\\hline
u&0123456&78,\ 9\,10\\
x&01237&45,\ 69,\ 8\,10\\
x/u&2578\,10&06,\ 14,\ 39
\end{array}$$
The long cycles have exponent steps $10,14,14$ modulo $70$ respectively; each pair has difference $35$. Thus all $3$ quotients are regular. The row $u$ is the second row of $T_f(q,s,t)$ at $q=w^{49}$, $s=w^{34}$. This is a witness for the pair extension problem, not a fourth row for a prescribed $T_f$.
\end{proof}

\subsection{The homogeneous CCC case}\label{sec:CCC}

\medskip
For type $\mathrm{CCC}$, write the rows as $1,u,v$ and put $z=u/v$. The chosen support partitions $P,Q,R$ have block sizes $(5,3,3)$, and each block carries a rotated prime cycle. The pentagon labels will be separated from the 4 continuous triangle phases.

For the enumeration fix
$$P=((0,1,2,3,4),(5,6,7),(8,9,10)).$$
There are $\binom{11}{5}\binom63/2=4620$ partitions of this shape. The stabilizer $S_5\times(S_3\wr S_2)$ of $P$ has $21$ orbits on $Q$. For one representative of each orbit we test every $R$, giving $97020$ support presentations. Arbitrary within-block column permutations are legitimate because no root labels have yet been assigned.

The tensor $n_{ijk}=|P_i\cap Q_j\cap R_k|$ determines the ordered supports up to column permutation. Its minimum under the $48$ quotient permutations and triangle swaps classifies the chosen supports.

\medskip
Rotate the pentagons of $u,v$ so that both are $\mu_5$. On a triangle in ascending column order every root assignment has the form
$$e^{2\pi i x}(1,w^{\sigma},w^{2\sigma}), \qquad w=e^{2\pi i/3},\quad \sigma\in\{1,-1\}.$$
There are 2 triangle phases for $u$ and 2 for $v$. Write them as $x\in(\mathbb R/\mathbb{Z})^4$. Let $D_P,D_Q$ be the $11\times4$ matrices whose rows are zero on the pentagon and select the appropriate one of these 4 variables on the triangles. Set $W=D_P-D_Q$.

For each block $S$ of $R$, let $j_S=\min S$. The requirement that the quotient has its prescribed relative roots on $S$ gives one equation for each $j\in S\setminus\{j_S\}$. There are 8 equations, with coefficient matrix
$$A_{(S,j)}=W_j-W_{j_S}\in\mathbb{Z}^4.$$
The right-hand side has the form
\[
Ax=f/5+g/3\pmod1,
\]
where $f$ records pentagon root labels and $g$ records the 6 triangle orientations in $P,Q,R$. Taking entrywise fifteenth powers gives the homogeneous system $Ax=0$ for the collapsed phases, with a rescaled variable.

\par\addvspace{\medskipamount}\noindent
\begin{minipage}{\linewidth}
\begin{lemma}\label{n11_noA_CCC:smith}
For every one of the $97020$ support presentations the nonzero Smith invariants of $A$ are all $1$, except possibly one invariant equal to $2$ or $3$. The kernel dimension is $4-\operatorname{rank}A$.
\end{lemma}
\end{minipage}
\par\addvspace{\medskipamount}

\begin{proof}
One program computes the gcds of the integer minors, of order at most $4$, and hence the rank and largest Smith invariant. Independently, another program supplies integer matrices $U\in\mathrm{GL}_8(\mathbb{Z})$, $V\in\mathrm{GL}_4(\mathbb{Z})$ with
$$UAV=\operatorname{diag}(d_0,\ldots,d_{r-1},0).$$
The verifier checks the displayed identity, $|\det U|=|\det V|=1$, positivity and divisibility of the $d_i$, and agreement with the minor calculation. All $97020$ certificates pass. The possible products $\prod d_i$ are exactly $1,2,3$.
\end{proof}

The integer invariants allow the quintic labels and the triangle phases to be treated separately. Keeping every torsion branch gives the following cover.

\par\addvspace{\medskipamount}\noindent
\begin{minipage}{\linewidth}
\begin{theorem}\label{n11_noA_CCC:main}
Every regular $3\times11$ matrix of chosen type $\mathrm{CCC}$ is equivalent to a member of a finite union of affine families of dimension at most 4, each with a seed whose entries lie in $\mu_{30}$ or $\mu_{45}$. The dimension bound is attained. Examples requiring Butson levels $30$ and $45$ occur.

There are $1617$ classes of chosen support triples under column permutations, permutations of the 3 quotient partitions, and interchanges of the 2 triangle blocks in each partition. Exactly $325$ are feasible. Their phase dimensions are distributed as follows:
$$\begin{array}{c|rrrrr}
\text{dimension}&0&1&2&3&4\\\hline
\text{feasible support classes}&262&46&14&2&1.
\end{array}$$
Of these classes, $322$ have connected collapsed phase kernels, one has kernel torsion of order $2$, and 2 have kernel torsion of order $3$.
\end{theorem}
\end{minipage}
\par\addvspace{\medskipamount}

\begin{proof}
Applying $U$ to the phase system $Ax=f/5+g/3$, the zero-row conditions separate modulo $5$ and $3$. Each nonzero $d_i$ is invertible modulo $5$, so a consistent quintic part has a solution $x_5\in(\frac15\mathbb Z/\mathbb Z)^4$. Subtracting it leaves $A(x-x_5)=g/3$. Quintic-label and ternary-phase consistency can therefore be tested independently, with arbitrary original phases.

\medskip
For the quintic test give each of $P,Q,R$ 7 labels in $\mathbb{F}_5$: 5 individual labels on its pentagon, and one shared label on each triangle. Denote the corresponding 11-column label vectors by $a,b,h$. Require
\[
a_j-b_j=h_j\quad(j=0,\ldots,10), \qquad\text{each pentagon's labels are a permutation of }\mathbb{F}_5.
\]
These are exactly the quintic conditions just derived. For feasibility we may shift the $u,v$ labels to put their first pentagon labels at zero, and use a common nonzero field scaling to set the second $u$ pentagon label to one. Scaling is an existence normalization, not a claimed Hadamard equivalence. The complete normal form allows every solution of $a_j-b_j=h_j$ with the stated pentagon permutations.

Of the $97020$ presentations, $52810$ fail by linear inconsistency or a forced equality of 2 pentagon labels. Certificates use the column equations, the gauge conditions and the necessary zero sums of the 5 labels on each pentagon. Another $100$ fail despite having no such immediate collision. For each of these the affine linear solution space has dimension 4 over $\mathbb{F}_5$; a separate verifier exhausts all $5^4=625$ points and finds no solution of the 3 permutation conditions. For all remaining $44110$ presentations, a concrete quintic-label witness is checked directly.

For the ternary test choose the 6 signs $\sigma=(\sigma_1,\ldots,\sigma_6)\in\{1,-1\}^6$. Define $e_P,e_Q,e_R$ to be zero on each pentagon and $(0,\sigma,2\sigma)$ on each triangle in ascending column order. On the row of $A$ indexed by $(S,j)$ put
$$g_{(S,j)}=(e_R-e_P+e_Q)_j-(e_R-e_P+e_Q)_{j_S}.$$
The sign choice is feasible precisely when $(Ug)_i=0\pmod3$ for every $i\ge r$. For every feasible sign choice, all solutions are
\[
x_3=Vy,\qquad y_i=\frac{(Ug)_i+3k_i}{3d_i}\ (i<r),\quad 0\le k_i<d_i,\qquad y_i=t_i\ (i\ge r),
\]
where the $t_i$ vary freely in $\mathbb R/\mathbb{Z}$. All indices here start at zero. Thus this is a calculation on the full continuous torus.

The ternary test removes another $24578$ quintic-feasible presentations, leaving $19532$. Their $325$ support classes are those in Theorem~\ref{n11_noA_CCC:main}. Before support deduplication the numbers of ternary connected components are
$$\begin{array}{c|rrrrr}
\text{dimension}&0&1&2&3&4\\\hline
\text{components}&105096&15412&2616&172&4.
\end{array}$$
These $123300$ components are counted with chosen support presentations and triangle signs; the counts do not enumerate matrices or all quintic labels.

\medskip
For any allowed labels $a,b,h$, signs, and branch in the displayed parametrization of $x_3$, set
\begin{align*}
u_j(t)&=\exp\!\left(2\pi i\left[\frac{a_j}{5} +\frac{e_{P,j}}3+(D_Px_3)_j\right]\right),\\
v_j(t)&=\exp\!\left(2\pi i\left[\frac{b_j}{5} +\frac{e_{Q,j}}3+(D_Qx_3)_j\right]\right).
\end{align*}
The identities $a_j-b_j=h_j$, with the stated pentagon permutations, supply the pentagons, the sign choices supply the triangles, and the phase equations ensure that every quotient block has the required common phase. Conversely, the separation argument applied to an arbitrary original triple supplies such quintic labels and a ternary solution, so the normal form is exhaustive.

At $t=0$ the ternary phases have denominators dividing $3,6,9$, giving seeds in $\mu_{15},\mu_{30},\mu_{45}$. The free columns of $V$ are primitive. The cycle identities hold for every allowed quintic assignment and every free parameter; all $123300$ presentations using stored witnesses are independently checked.

The 4-parameter bound is attained by the common-block family. Put $w=e^{2\pi i/15}$ and $\omega=w^5=e^{2\pi i/3}$. Its rows, written in blocks, are
$$\begin{array}{c|c|c}
j=0,\ldots,4&j=0,1,2&j=0,1,2\\\hline
1&1&1\\
w^{3j}&s\omega^j&r\omega^j\\
w^{6j}&t\omega^{2j}&b\omega^{2j}
\end{array}
\qquad s,r,t,b\in\mathbb{T}.$$
Cross ratios recover the 4 parameters, up to finite permutation ambiguities, proving that this is genuinely 4-dimensional modulo equivalence.

For a concrete level-$45$ example take first row all ones and other rows $u_j=w^{a_j},v_j=w^{b_j}$, where $w=e^{2\pi i/45}$ and
$$\begin{pmatrix}a\\b\end{pmatrix}=
\begin{pmatrix}
0&9&18&27&36&22&37&7&26&11&41\\
0&18&36&43&8&9&28&38&27&13&23
\end{pmatrix}.$$
Cycle certificates, with columns numbered from zero, are
$$\begin{array}{c|c|c|c}
&\text{pentagon}&\text{triangle}&\text{triangle}\\\hline
u&\{0,1,2,3,4\}&\{5,6,7\}&\{8,9,10\}\\
v&\{0,1,2,5,8\}&\{3,6,9\}&\{4,7,10\}\\
u/v&\{0,1,2,6,10\}&\{3,7,8\}&\{4,5,9\}
\end{array}$$
The cross ratio from the first 2 rows and columns $0,5$ is $w^{22}$ or its inverse, a primitive $45$th root. Thus this matrix cannot be equivalent to one over $\mu_m$ unless $45\mid m$.

A level-$30$ example, with the same convention and $w=e^{\pi i/15}$, has exponent rows
$$\begin{pmatrix}
0&6&12&18&24&15&5&25&7&17&27\\
0&24&19&9&11&12&6&1&18&29&21
\end{pmatrix}.$$
Its first nontrivial row has exponent $7$ at column $8$, giving a primitive $30$th-root cross ratio with column $0$. Full cycle certificates for both examples are included in \texttt{examples.json}. These examples have zero-dimensional chosen presentations; no global isolation claim is made.
\end{proof}

\medskip
We now determine the generic extensions of the common-block construction. Put $w=e^{2\pi i/15}$ and $\omega=w^5$, and write $v_h=(1,\omega^h,\omega^{2h})$ for $h=0,1,2$. For $a=2,3,4$ define
\[
S_a(s,r,t,b)=
\begin{array}{c|c|c}
\text{columns }j=0,\ldots,4&\text{next 3 columns}&\text{last 3 columns}\\\hline
1&v_0&v_0\\
w^{3j}&s v_1&r v_1\\
w^{3aj}&t v_2&b v_2
\end{array},\qquad s,r,t,b\in\mathbb T.
\]
Here the common block partition has sizes $5,3,3$; having only a common pentagon is a weaker condition.

\par\addvspace{\medskipamount}\noindent
\begin{minipage}{\linewidth}
\begin{lemma}\label{ccc_common:normal}
Every regular $\mathrm{CCC}$ triple whose 3 chosen cycle partitions coincide is equivalent to a matrix $S_a(s,r,t,b)$. Every $4$-dimensional chosen $\mathrm{CCC}$ phase presentation has coincident cycle partitions.
\end{lemma}
\end{minipage}
\par\addvspace{\medskipamount}

\begin{proof}
Dephase the first row and rotate the pentagons of the other 2 rows to $\mu_5$. Permute their common pentagon columns so that the second row is $(w^{3j})_{j=0}^4$. The third has labels $f(j)$, with $f(0)=0$ and both $f$ and $f-j$ permutations of $\mathbb F_5$. Checking the $4!=24$ normalized permutations gives exactly $f(j)=aj$, $a=2,3,4$. On each triangle, normalize and order the second row as a multiple of $v_1$. The third row must then be a multiple of $v_2$, since its quotient with the second row is a triangle. The 4 block phases are unrestricted, giving the stated matrix $S_a(s,r,t,b)$.

For the last assertion, use the 4 triangle-phase variables of the phase system $Ax=f/5+g/3$. Dimension $4$ means that its matrix $A$ is zero. Thus the vector $D_P(j)-D_Q(j)$ is constant on every block of $R$. Because the variables of $P$ and $Q$ occupy disjoint coordinates, this means that the pair of block memberships $(P(j),Q(j))$ is constant there. The $5$ columns of the pentagon of $R$ must lie in both pentagons, so all 3 pentagons coincide. Each triangle of $R$ must likewise equal a triangle of $P$ and of $Q$.
\end{proof}

$\mathrm{CCC}$ triples of this form can reach $5$ rows. The next result determines every generic fourth row and its compatibility with the other choices; the mixed construction in Section~\ref{sec:mixed} gives a further consequence.

\par\addvspace{\medskipamount}\noindent
\begin{minipage}{\linewidth}
\begin{theorem}\label{ccc_common:extension}
There is an explicit proper closed subset $E\subset\mathbb T^4$, given by a union of $3608$ monomial equations, such that the regular extensions of $S_a(s,r,t,b)$ have maximum size $5$ whenever $(s,r,t,b)\notin E$.

More precisely, put
$$K_a=\mathbb F_5\setminus\{0,1,a\},\qquad c_0=1,\quad c_1=r/s,\quad c_2=b/t.$$
Every normalized fourth row is one of the $6$ circles
\[
x_{kh}(z)=\bigl((w^{3kj})_{j=0}^4\mid z v_h\mid -z c_h v_h\bigr),
\qquad k\in K_a,\quad h=0,1,2,\quad z\in\mathbb T.
\]
Two such rows $x_{kh}(z),x_{lm}(z')$ can be added together precisely when $k\ne l$ and $h\ne m$, with arbitrary $z,z'$. Thus their compatibility graph is a $6$-cycle. Every generic $5$-row extension is obtained by choosing one of its $6$ edges and 2 independent unit phases.

All these fourth rows and $5$-row extensions are valid for every $s,r,t,b\in\mathbb T$. The exceptional-set restriction concerns their exhaustiveness and the upper bound.
\end{theorem}
\end{minipage}
\par\addvspace{\medskipamount}

\begin{proof}
For $h=0$, the last 6 entries of $x_{kh}(z)$ form 3 antipodal pairs, while its quotients with the second and third rows have 2 triangles. For $h=1$, those pairs occur in its quotient with the second row, and for $h=2$ in its quotient with the third. The other 2 quotients have 2 triangles. Their first 5 entries are pentagons because $k\notin\{0,1,a\}$. Hence every displayed row extends the triple regularly, with no parameter restriction.

To prove completeness, use the same phase-graph procedure as in Theorem~\ref{bbb_common:extension}, now with 4 old parameters. A new quotient cannot have profile $\mathrm{A}$, by Theorem~\ref{n11_A_branch:profiles}. Enumerate every decorated presentation of the new row with profile $\mathrm{B},\mathrm{C},\mathrm{D},\mathrm{E},\mathrm{F}$, then every prime-cycle cover against the second and third rows. Forest integration retains arbitrary unit phases. Closing edges either give identities, inconsistent constant equations, or nonconstant monomial equations in $s,r,t,b$. Record the latter before rejecting the branch.

Let $E$ be the union of the recorded nonconstant equations in the audit calculation. Every rejected branch is impossible outside $E$. There are $3608$ recorded equations, of the form
$$s^i r^j t^k b^l=e^{\pi i d/105},\qquad
0<\max(|i|,|j|,|k|,|l|)\le3,\quad d\in\mathbb Z/210\mathbb Z.$$
Their characters are nonzero, so the finite union is closed with empty interior. The \texttt{ccc\_common} component supplies the complete list in \texttt{exceptional\_union.txt}, with lines $i\ j\ k\ l\ d$. This exceptional set is sufficient and is not asserted to be minimal.

The 2 implementations agree on all retained pair and triple records after exhausting $1480490$ initial presentations each. The independent verifier checks literal prime-cycle covers, the $24$ pentagon-label possibilities, the zero-rank support assertion over all $4620$ relative partitions, and the containment of every final record in the circles $x_{kh}(z)$. The table gives the initial presentation count, the records against the first 2 old rows, and the final records summed over the $3$ values of $a$:
$$\begin{array}{c|r|r|r}
P&\text{initial presentations}&\text{pair records}&\text{triple records}\\\hline
\mathrm{B}&712800&0&0\\
\mathrm{C}&443520&4383&852\\
\mathrm{D}&166320&672&426\\
\mathrm{E}&123200&0&0\\
\mathrm{F}&34650&0&0
\end{array}$$
For each $a$, the $284$ records of type $\mathrm{C}$ consist of the $4$ circles with $h=1,2$ and $280$ fixed specializations already on those circles. The $142$ records of type $\mathrm{D}$ consist of the $2$ circles with $h=0$ and $140$ such specializations. These are counts of normalized records, not equivalence classes.

It remains to decide which candidates can be added simultaneously. The pentagon contribution to their inner product is $5$ if $k=l$ and zero otherwise. The last 6 columns contribute $6z\bar z'$ if $h=m$ and zero otherwise. Since $|z\bar z'|=1$ and $5\ne6$, the sum vanishes exactly when $k\ne l$ and $h\ne m$. The resulting graph has 2 parts of size $3$, with the 3 equal-mode edges removed: it is a $6$-cycle. There are compatible pairs but no compatible triples. This proves both the maximum $5$ and the description of all extensions attaining it.
\end{proof}

The theorem settles the generic extension problem for the common-block $\mathrm{CCC}$ families, which include every $4$-dimensional chosen phase presentation. The exceptional parameters and the other $\mathrm{CCC}$ support configurations remain open. In particular, the theorem does not impose a common-pentagon restriction on the new rows; their possible supports were exhausted.


\subsection{The homogeneous DDD case}\label{sec:DDD}

\medskip
The homogeneous type $\mathrm{DDD}$ has one pentagon and 3 antipodal pairs in each chosen quotient. As before, normalize the rows to $1,u,v$, put $z=u/v$, and denote their support partitions by $P,Q,R$. The 6 pair phases lead to a binary phase system coupled to a separate quintic label problem.

Fix
$$P=((0,1,2,3,4),(5,6),(7,8),(9,10)).$$
There are $\binom{11}{5}\,15=6930$ partitions of this shape. The stabilizer $S_5\times(C_2\wr S_3)$ of $P$ has $26$ orbits on $Q$. Testing every $R$ for these representatives gives $180180$ support presentations.

Canonicalize the intersection tensor under the $1296$ quotient and pair-block permutations. A separate exhaustive group calculation verifies all canonical tensors and accounts for the $180180$ presentations.

\medskip
Rotate the pentagons of $u,v$ to $\mu_5$. On a pair $(a,b)$, $a<b$, write the entries as $e^{2\pi i x}(1,-1)$. Let $x\in(\mathbb R/\mathbb{Z})^6$ contain the 3 pair phases of $u$ and the 3 of $v$. Define $D_P,D_Q$ to be the $11\times6$ matrices selecting these variables on the pairs and having zero rows on the pentagons. Set $W=D_P-D_Q$.

Let $\epsilon_P,\epsilon_Q,\epsilon_R$ be zero on pentagons and the first column of each pair, and one on the second column. For every block $S$ of $R$, with $j_S=\min S$, and $j\in S\setminus\{j_S\}$, put
$$A_{(S,j)}=W_j-W_{j_S},\qquad g_{(S,j)}=(\epsilon_R-\epsilon_P+\epsilon_Q)_j -(\epsilon_R-\epsilon_P+\epsilon_Q)_{j_S}.$$
There are 7 equations in 6 phase variables. With the pentagon root labels included, the full system has the form
\[
Ax=f/5+g/2\pmod1,
\]
where $f$ is an integer vector determined by the pentagon label differences.

\par\addvspace{\medskipamount}\noindent
\begin{minipage}{\linewidth}
\begin{lemma}\label{n11_noA_DDD:smith}
For every support presentation, all nonzero Smith invariants of $A$ belong to $\{1,2,3,4\}$. In particular they are coprime to $5$. Whenever $Ax=g/2$ is consistent, the invariants belong to $\{1,2,4\}$ and $\operatorname{rank}A\ge3$.
\end{lemma}
\end{minipage}
\par\addvspace{\medskipamount}

\begin{proof}
Integer reduction supplies unimodular matrices $U\in\mathrm{GL}_7(\mathbb{Z})$, $V\in\mathrm{GL}_6(\mathbb{Z})$ such that
$$UAV=\operatorname{diag}(d_0,\ldots,d_{r-1},0).$$
The verifier checks this identity, $|\det U|=|\det V|=1$, and positivity and divisibility of the $d_i$, for all $180180$ systems. A second algorithm uses gcds of exact integer minors to check the rank, torsion order and largest Smith invariant. Binary consistency is tested by $(Ug)_i=0\pmod2$ for all $i\ge r$. The claimed invariant and rank restrictions hold in every consistent case.
\end{proof}

We now combine the binary phase equations with the quintic label test. The result includes all continuous phases, not just a finite root grid.

\par\addvspace{\medskipamount}\noindent
\begin{minipage}{\linewidth}
\begin{theorem}\label{n11_noA_DDD:main}
Every regular $3\times11$ matrix of chosen type $\mathrm{DDD}$ is equivalent to a member of a finite union of affine families with at most 3 parameters and a seed over $\mu_{40}$. The dimension bound is attained, and some examples cannot be equivalent to a partial Butson matrix over $\mu_m$ unless $40\mid m$.

There are $3300$ classes of chosen support triples under column permutations, permutations of the 3 quotient partitions, and permutations of the 3 pair blocks in each partition. Exactly $452$ are feasible, with phase dimensions
$$\begin{array}{c|rrrr}
\text{dimension}&0&1&2&3\\\hline
\text{feasible support classes}&218&190&41&3.
\end{array}$$
All but one feasible support class have affine presentations with seeds over $\mu_{20}$. The remaining class has Smith invariant $4$ and requires $\mu_{40}$ for its seed presentations.
\end{theorem}
\end{minipage}
\par\addvspace{\medskipamount}

\begin{proof}
Applying $U$ to the phase system $Ax=f/5+g/2$, the zero rows require $(Uf)_i/5+(Ug)_i/2=0$ modulo $1$. Coprimality separates these into $(Uf)_i=0\pmod5$ and $(Ug)_i=0\pmod2$. Since each nonzero $d_i$ is invertible modulo $5$, any consistent quintic part has a solution with coordinates in $\tfrac15\mathbb{Z}/\mathbb{Z}$. Subtracting that solution leaves the full binary phase system $Ax_2=g/2$. This argument treats arbitrary initial phases; it does not impose a root-order cutoff on the original matrix.

\medskip
The binary test excludes $164212$ support presentations, leaving $15968$. For each survivor introduce 8 labels over $\mathbb{F}_5$ for each of $P,Q,R$: 5 individual labels on the pentagon and one shared label on each pair. Write the corresponding 11-column label vectors as $a,b,h$. The quintic conditions are exactly
\[
a_j-b_j=h_j\quad(j=0,\ldots,10),\qquad \text{each pentagon's labels form a permutation of }\mathbb{F}_5.
\]
For existence testing we set the first $u,v$ pentagon labels to zero and the second $u$ pentagon label to one, using row-label shifts and a common nonzero field scaling. Scaling is not taken as a Hadamard equivalence: the normal form permits every solution of $a_j-b_j=h_j$ with the stated pentagon permutations.

Of the $15968$ survivors, $1348$ fail by exact linear certificates over $\mathbb{F}_5$: the column equations, the gauge conditions and the zero sums of pentagon labels are inconsistent or force 2 pentagon labels to coincide. The other $14620$ have directly checked label witnesses. No search-based nonexistence assertion is needed for this step.

For a feasible support let $r=\operatorname{rank}A$. All binary solutions are
\[
x_2=Vy,\qquad y_i=\frac{(Ug)_i+2k_i}{2d_i}\quad(0\le i<r),\quad 0\le k_i<d_i,\qquad y_i=t_i\quad(r\le i<6),
\]
with free $t_i\in\mathbb R/\mathbb{Z}$. For any allowed quintic labels the original rows are
\begin{align*}
u_j(t)&=\exp\!\left(2\pi i\left[\frac{a_j}{5} +\frac{\epsilon_{P,j}}2+(D_Px_2)_j\right]\right),\\
v_j(t)&=\exp\!\left(2\pi i\left[\frac{b_j}{5} +\frac{\epsilon_{Q,j}}2+(D_Qx_2)_j\right]\right).
\end{align*}
The label equations supply the pentagons, the signs supply the pairs, and the phase system gives the required common phase on each quotient block. The separation argument proves that every original $\mathrm{DDD}$ triple is covered.

There are $6-r\le3$ parameters. At $t=0$ the binary phases have denominators dividing $8$, so the entries lie in $\mu_{40}$. If no $d_i=4$ occurs, the seed lies in $\mu_{20}$. Exactly one feasible support class has such an invariant. Before support deduplication, the numbers of connected binary components are
$$\begin{array}{c|rrrr}
\text{dimension}&0&1&2&3\\\hline
\text{components}&17288&10902&1538&34.
\end{array}$$
The verifier checks every one of these $29762$ affine presentations symbolically with the stored quintic witness: all required cycles occur, and every parameter exponent is constant on its cycle. The formulas imply the same for every allowed quintic assignment. These component counts are redundant and do not count inequivalent matrices or all quintic labels.

\medskip
Let $w=e^{2\pi i/5}$. A simple 3-parameter family is
$$\begin{array}{c|c}
\text{columns }j=0,\ldots,4&\text{last 6 columns}\\\hline
1&(1,1,1,1,1,1)\\
w^{j}&s(1,-1,i,-i,b,-b)\\
w^{2j}&t(1,i,-1,-i,-b,b)
\end{array}
\qquad s,t,b\in\mathbb{T}.$$
On the last block, $u/v=(s/t)(1,i,-i,1,-1,-1)$, which splits into 3 antipodal pairs. Each individual row does so as well. Cross ratios recover $s,t,b$, proving that the dimension bound is attained modulo the finite ambiguities of monomial equivalence.

For a genuine level-$40$ seed set $w=e^{\pi i/20}$, take first row all ones, and take the other rows to have exponent vectors
$$\begin{pmatrix}
0&8&16&24&32&10&30&26&6&7&27\\
0&8&30&6&37&16&10&24&26&32&17
\end{pmatrix}.$$
The cross ratio from the first 2 rows and columns $0,9$ is $w^{7}$ or its inverse, a primitive $40$th root. Hence no equivalent partial Butson matrix can have root order not divisible by $40$. Full cycle certificates are recorded in \texttt{examples.json}. Its chosen presentation is zero-dimensional; no global isolation claim is made.
\end{proof}

The seed examples also lead to a restriction on all Butson representatives of a $\mathrm{DDD}$ triple.

\par\addvspace{\medskipamount}\noindent
\begin{minipage}{\linewidth}
\begin{proposition}\label{n11_noA_DDD:rootobstruction}
If a $\mathrm{DDD}$ triple is equivalent to a partial Butson matrix over $\mu_m$, then $20\mid m$.
\end{proposition}
\end{minipage}
\par\addvspace{\medskipamount}

\begin{proof}
Its pentagon and pair quotients require $10\mid m$. The exact certificate check also finds that none of the binary-consistent systems admits a half-integral solution: in Smith coordinates some even $d_i$ always has $(Ug)_i$ odd. If $m=2n$ with $n$ odd, normalization of the pentagons keeps all entries in $\mu_m$, and the pair phases $x$ in the phase system $Ax=f/5+g/2$ have denominator dividing $m$. Since $5\mid n$, multiplication by $n$ gives $A(nx)=g/2$ modulo $1$, with $nx$ half-integral, a contradiction.
\end{proof}

\medskip\label{sec:DDD_five}
The following construction will be used in the extension problem.

\par\addvspace{\medskipamount}\noindent
\begin{minipage}{\linewidth}
\begin{proposition}\label{n11_DDD:fiveconstruction}
There is a $4$-parameter family $H(t)$ of regular $5\times11$ partial Hadamard matrices for which every row quotient has profile $\mathrm{D}$, with the same pentagon support.
\end{proposition}
\end{minipage}
\par\addvspace{\medskipamount}

\begin{proof}
Use the Paley conference construction~\cite{pal}: let $C_6$ be the real symmetric $6\times6$ matrix indexed by $\{\infty\}\cup\mathbb{F}_5$ with zero diagonal, $(C_6)_{\infty,a}=(C_6)_{a,\infty}=1$, and $(C_6)_{a,b}=\chi(a-b)$ for distinct finite indices, where $\chi$ is the quadratic character modulo $5$. Direct multiplication gives $C_6^2=5I_6$. Thus $Q_6=I_6+iC_6$ is a Hadamard matrix of order 6, since $Q_6Q_6^*=I_6+C_6^2=6I_6$.

For 2 distinct rows of $Q_6$, the quotient entries at their diagonal positions are opposite imaginary units. The other 4 entries are real signs whose sum is zero, so they split into 2 further antipodal pairs. Every quotient of $Q_6$ consequently has profile $2+2+2$.

Join the 5 rows of $F_5$ to any 5 distinct rows of $Q_6$. Orthogonality holds separately on each block, giving a regular $5\times11$ matrix with every quotient of type $\mathrm{D}$. Multiply the 6-column part of row $i$ by an arbitrary $t_i\in\mathbb{T}$. After setting $t_0=1$ by equivalence, the remaining 4 phases give a genuine 4-parameter family.

Write $H(t)$ for the family obtained from the following dephased seed by multiplying the last 6 entries of rows $1,2,3,4$ by $t_1,t_2,t_3,t_4$, respectively, with rows numbered from zero. The seed $H(1,1,1,1)$ is obtained by raising $w=e^{\pi i/10}$ entrywise to
$$\begin{pmatrix}
0&0&0&0&0&0&0&0&0&0&0\\
0&4&8&12&16&0&10&15&5&5&15\\
0&8&16&4&12&0&15&10&15&5&5\\
0&12&4&16&8&0&5&15&10&15&5\\
0&16&12&8&4&0&5&5&15&10&15
\end{pmatrix}.$$
All 10 quotient decompositions are checked exactly and retained in \texttt{extension\_5x11.json}.
\end{proof}

A sixth row cannot retain this same pentagon support against all 5 existing rows: its first 5 entries would have to be orthogonal to the 5 rows of $F_5$. This support argument alone leaves other decompositions open. Theorem~\ref{n11_full:main} removes this support restriction: no member of the full $4$-parameter family admits a regular sixth row, with any quotient profiles.

\section{The triangle and pair cases}\label{sec:triangles_pairs}
The remaining homogeneous profiles use only triangles and antipodal pairs. We treat $\mathrm{EEE}$ and $\mathrm{FFF}$ in turn, and then collect the six homogeneous branches in one theorem.

\subsection{The homogeneous EEE case}\label{sec:EEE}

\medskip
For type $\mathrm{EEE}$, each quotient has 3 triangles and one antipodal pair. There is no separate pentagon or seventh-root label problem: the triangle orientations and the pair signs enter directly into 7 equations for 6 continuous phases.

\par\addvspace{\medskipamount}\noindent
\begin{minipage}{\linewidth}
\begin{theorem}\label{n11_EEE:main}
Every regular $\mathrm{EEE}$ triple is equivalent to a member of a finite union of affine families with at most 4 parameters. Each family has a seed with root level in the following set:

$$\{6,12,18,24,30,36,42\}.$$

The 4-parameter bound is attained. An explicit seed with a primitive 42nd-root cross ratio is given below.
\end{theorem}
\end{minipage}
\par\addvspace{\medskipamount}

\begin{proof}
There are exactly 6,904 feasible classes of chosen support triples, under column permutations, permutations of the 3 row quotients, and permutations of the 3 triangle blocks in each quotient. Their phase dimensions are:

\begin{center}
\begin{tabular}{@{}lr@{}}
\toprule
Dimension & Feasible support classes \\
\midrule


0 & 5,230 \\
1 & 1,544 \\
2 & 121 \\
3 & 8 \\
4 & 1 \\
\end{tabular}
\end{center}

\medskip

Fix $P=((0,1),(2,3,4),(5,6,7),(8,9,10))$. There are $15400$ partitions of shape $(2,3,3,3)$ and $42$ orbits for the second partition under its stabilizer. Testing the third partition gives $646800$ presentations. Canonicalizing the intersection tensors under the $1296$ quotient and triangle-block permutations gives $6904$ feasible support classes.

\medskip
Rotate $u$ and $v$ so that the first entry on each row's chosen pair is 1; the other entry is then $-1$. Each row has 3 remaining continuous parameters, one phase for each triangle. Let $x$ collect these 6 parameters.

In increasing column order, an arbitrary triangle is a phase times one of the 2 orientations below. Its 3 cyclic rotations are absorbed in the phase:

$$(1,w^{\sigma},w^{2\sigma}),\qquad w=e^{2\pi i/3},\quad \sigma\in\{1,-1\}.$$

There are 9 triangle orientations altogether, 3 for each of $P,Q,R$. Let $\sigma$ be their vector. Let $D_P,D_Q$ select the 6 triangle phases on their respective triangles, with zero rows on the pairs. Write $W = D_P - D_Q$.

Let $\epsilon_P,\epsilon_Q,\epsilon_R$ be $0$ except at the second column of the corresponding pair, where they are $1$. Let $e_P,e_Q,e_R$ be $0$ on pairs and $(0,\sigma,2\sigma)$ on each triangle. For a block $S$ of $R$, let $j_S$ be its first column. For every other column $j$ in $S$, define:

$$\begin{aligned}
A_{(S,j)}&=W_j-W_{j_S},\\
f_{(S,j)}&=(\epsilon_R-\epsilon_P+\epsilon_Q)_j -(\epsilon_R-\epsilon_P+\epsilon_Q)_{j_S},\\
(B\sigma)_{(S,j)}&=(e_R-e_P+e_Q)_j-(e_R-e_P+e_Q)_{j_S}.
\end{aligned}$$

The complete phase system is:

$$Ax=\frac f2+\frac{B\sigma}{3}\pmod1, \qquad A\in M_{7\times6}(\mathbb Z),\quad B\in M_{7\times9}(\mathbb Z).$$

There is no further root-label problem: a sign gives all triangle orientations up to the phase already present, and the pair orientation is fixed.

\medskip
Integer reduction supplies unimodular matrices $U,V$ such that:

$$UAV=\operatorname{diag}(d_0,\ldots,d_{r-1},0), \qquad U\in\mathrm{GL}_7(\mathbb Z),\quad V\in\mathrm{GL}_6(\mathbb Z).$$

For every one of the 646,800 systems, the calculation verifies this identity, both determinants being $\pm1$, and positivity and divisibility of the nonzero diagonal entries. The zero rows impose exactly:

$$(Uf)_i=0\pmod2,\qquad (UB\sigma)_i=0\pmod3, \qquad r\leq i<7.$$

The separation follows from coprimality of 2 and 3. Binary consistency leaves 152,132 support presentations. Exhausting all 512 orientation vectors leaves 147,642 feasible presentations.

For every feasible orientation, all continuous solutions are given by:

$$\begin{aligned}
x&=Vy,\\
y_i&=\frac{3(Uf)_i+2(UB\sigma)_i+6k_i}{6d_i} && (0\leq i<r),\qquad 0\leq k_i<d_i,\\
y_i&=t_i &&(r\leq i<6),\qquad t_i\in\mathbb R/\mathbb Z.
\end{aligned}$$

Recover the 2 rows by:

$$\begin{aligned}
u_j(t)&=\exp\!\left(2\pi i\left[\epsilon_{P,j}/2+e_{P,j}/3+(D_Px)_j\right]\right),\\
v_j(t)&=\exp\!\left(2\pi i\left[\epsilon_{Q,j}/2+e_{Q,j}/3+(D_Qx)_j\right]\right).
\end{aligned}$$

These formulas ensure every required cycle, for all parameter values. Conversely, normalizing any original $\mathrm{EEE}$ triple gives its triangle orientations and a solution of this system, proving coverage.

All feasible systems have rank at least 2, giving at most 4 parameters. The largest Smith invariant is at most 8. Setting the free parameters to zero initially bounds the seed level by 6 times that invariant.

The only potential seed level above 42 is 48, from invariant 8. Every orientation and torsion branch was checked for all 7 such feasible presentations; their seed levels are only 6, 12 and 24. More generally, every orientation and branch of every feasible presentation with largest invariant at least 3 was audited --- 3,511,344 phase components across 7,839 presentations. The actual seed levels are precisely 6, 12, 18, 24, 30, 36 and 42. Invariants one and 2 already restrict the level to a divisor of 6 or 12.

There are 19,460,244 feasible orientation choices and 27,688,632 phase-component presentations before support reduction:

\begin{center}
\begin{tabular}{@{}lr@{}}
\toprule
Dimension & Phase-component presentations \\
\midrule


0 & 25,144,768 \\
1 & 2,468,752 \\
2 & 73,592 \\
3 & 1,504 \\
4 & 16 \\
\end{tabular}
\end{center}

These counts include repeated descriptions and are not counts of inequivalent matrices or global connected components.

\medskip
Let $w=e^{\pi i/6}$ and let $B$ be the 3-by-5 matrix obtained by raising $w$ entrywise to:

$$\begin{pmatrix}
0&0&0&0&0\\
0&6&3&7&11\\
0&3&6&11&7
\end{pmatrix}.$$

Every row quotient of $B$ splits as $2+3$. With the standard 3-by-3 Fourier matrix $F_3$, form:

\[
H(s,t,a,b)= \left[\ B\ \middle|\ \operatorname{diag}(1,s,t)F_3\ \middle|\ \operatorname{diag}(1,a,b)F_3\ \right], \qquad s,t,a,b\in\mathbb T.
\]

Each quotient now has profile $2+3+3+3$. Cross ratios recover the 4 parameters, so the dimension is genuinely 4 modulo the finite permutation ambiguities.

\medskip
Take the first row to be all ones. With $w=e^{\pi i/21}$, take the other rows to be $w$ raised entrywise to:

$$\begin{pmatrix}
0&21&38&24&10&18&4&32&41&13&27\\
0&32&21&18&24&4&10&5&38&33&19
\end{pmatrix}.$$

The chosen cycles, with columns numbered from zero, are:

\begin{center}
\begin{tabular}{@{}lllll@{}}
\toprule
Quotient & Pair & Triangle 1 & Triangle 2 & Triangle 3 \\
\midrule


$u$ & 0,1 & 2,3,4 & 5,6,7 & 8,9,10 \\
$v$ & 0,2 & 1,3,5 & 4,6,8 & 7,9,10 \\
$u/v$ & 3,7 & 0,4,5 & 1,2,8 & 6,9,10 \\
\end{tabular}
\end{center}

The cross ratio from the first 2 rows and columns $0,8$ is $w^{41}$ or its inverse. It is primitive of order 42, so this matrix has no equivalent Butson representative of root order not divisible by 42. Its chosen presentation is zero-dimensional; this does not assert global isolation.
\end{proof}

To use this construction throughout the maximal-dimensional $\mathrm{EEE}$ branch, we first identify every presentation of dimension four.

\par\addvspace{\medskipamount}\noindent
\begin{minipage}{\linewidth}
\begin{lemma}\label{eee_generic:normal}
Every member of a $4$-dimensional chosen $\mathrm{EEE}$ phase presentation is equivalent to a member of the family $H(s,t,a,b)$ of Subsection~\ref{sec:EEE} or its entrywise conjugate.
\end{lemma}
\end{minipage}
\par\addvspace{\medskipamount}

\begin{proof}
The unique $4$-dimensional support class in Theorem~\ref{n11_EEE:main} has representative
$$\begin{aligned}
P&=((0,1),(2,3,4),(5,6,7),(8,9,10)),\\
Q&=((0,2),(1,3,4),(5,6,7),(8,9,10)),\\
R&=((1,2),(0,3,4),(5,6,7),(8,9,10)).
\end{aligned}$$
Normalize the first row to $1$ and the other rows to $u_0=v_0=1$. Return to $w=e^{\pi i/6}$ and put $\omega=w^4$. Their first 5 entries have the form
$$u=(1,-1,x,x\omega^\sigma,x\omega^{-\sigma}),\qquad
v=(1,y,-1,y\omega^\tau,y\omega^{-\tau}),\qquad \sigma,\tau\in\{1,-1\}.$$
The pair of $R$ gives $xy=-1$. Its triangle gives $\tau=-\sigma$ and $x/y=1$, hence $x=y=\pm i$. Exchanging columns $3,4$ handles the choice of $\sigma$, and conjugation handles the sign of $i$.

On either common triangle, the orientations of $u$ and $v$ must be opposite. Order its columns so that they are multiples of the second and third rows of $F_3$. Their 4 phases are unrestricted. The 4 core choices and the 2 orientation choices on each common triangle account for all 16 maximal-dimensional phase-component presentations. This proves the claimed normal forms; conjugation is stated separately from ordinary equivalence.
\end{proof}

Before testing a new row by phase graphs, we exclude its possible $11$-cycle quotients. The fixed five-column core gives an obstruction at every parameter value.

\par\addvspace{\medskipamount}\noindent
\begin{minipage}{\linewidth}
\begin{lemma}\label{eee_generic:noA}
For every parameter value in the family $H(s,t,a,b)$ of Subsection~\ref{sec:EEE}, a regular added row has no $\mathrm{A}$ quotient with an old row.
\end{lemma}
\end{minipage}
\par\addvspace{\medskipamount}

\begin{proof}
In any regular triple with an $\mathrm{A}$ quotient and a chosen $\mathrm{E}$ quotient, every ratio of 2 entries of the $\mathrm{E}$ quotient belongs to $\mu_{66}$. Indeed, Theorem~\ref{n11_A_branch:profiles} leaves $\mathrm{AEE}$ and $\mathrm{AEF}$. Their simple bipartite support graph has 11 edges and minimum degree 2. If disconnected, it has 2 components, the smaller with 4 or 5 edges. Five edges are impossible: each side would have at most 2 vertices. Four edges give a square, requiring 2 degree-2 vertices on the $\mathrm{E}$ side, which has only one. The graph is therefore connected. The phase integration in that theorem puts all entries of either non-$\mathrm{A}$ quotient in one rotated copy of $\mu_{66}$.

Write the rows of the family $H(s,t,a,b)$ of Subsection~\ref{sec:EEE} as $1,u,v$. Both $u$ and $v$ contain the entry ratio $i$, whose order does not divide 66. Applying the preceding observation to $(1,u,x)$ excludes an $\mathrm{A}$ quotient between a new row $x$ and either $1$ or $u$; applying it to $(1,v,x)$ excludes one between $x$ and $v$.
\end{proof}

The remaining profiles can be tested with the four old phases left free. This gives a generic exclusion with an explicit exceptional set.

\par\addvspace{\medskipamount}\noindent
\begin{minipage}{\linewidth}
\begin{theorem}\label{eee_generic:extension}
There is an explicit proper finite union $Z\subset\mathbb T^4$ of monomial loci such that $H(s,t,a,b)$ has no regular fourth row whenever $(s,t,a,b)\notin Z$. The same holds for its entrywise conjugate. Thus every $4$-dimensional chosen $\mathrm{EEE}$ phase presentation has generic maximum extension exactly $3$.
\end{theorem}
\end{minipage}
\par\addvspace{\medskipamount}

\begin{proof}
Lemma~\ref{eee_generic:noA} excludes all new $\mathrm{A}$ quotients, for every parameter value. For the other profiles use the phase-graph procedure of Theorem~\ref{bbb_common:extension}, now with independent parameters $s,t,a,b$ and root denominator $420=\operatorname{lcm}(12,2,3,5,7)$. Enumerate every decorated presentation of a proposed new row $x$ with profile $\mathrm{B},\mathrm{C},\mathrm{D},\mathrm{E},\mathrm{F}$, then every prime-cycle cover of $x/u$ and $x/v$. Tree integration retains one arbitrary unit phase per component. A closing edge gives an identity, an inconsistent constant equation, or a nonconstant monomial equation in the 4 old parameters. Record each nonconstant equation before rejecting its branch.

Put $w=e^{\pi i/210}$. The audit records 5108 distinct equations of the form
$$s^it^ja^kb^l=w^d,\qquad
0<\max(|i|,|j|,|k|,|l|)\le3,\qquad d\in\mathbb Z/420\mathbb Z.$$
Let $Z$ be their union. Its complement is open and dense, and every rejected nonconstant branch is impossible there. The equations are retained in \texttt{exceptional\_union.txt} of the \texttt{eee\_generic} component, with one line $i\ j\ k\ l\ d$ per equation. This is a sufficient exceptional set, not a minimal one.

Two implementations, using forest potentials and explicit component potentials with opposite cover orders, each exhaust $1480490$ initial presentations. They agree on every retained pair and triple record:
$$\begin{array}{c|r|r|r}
\text{profile of }x&\text{presentations}&\text{against }1,u&\text{against }1,u,v\\\hline
\mathrm{B}&712800&0&0\\
\mathrm{C}&443520&0&0\\
\mathrm{D}&166320&0&0\\
\mathrm{E}&123200&4184&0\\
\mathrm{F}&34650&6418&0
\end{array}$$
An independent verifier checks literal prime-cycle covers for all $10602$ pair records, the agreement of the complete scans, the exceptional equations and the small normal-form calculations. The initial counts are $11!/(\prod p\prod m_p!)$, where the product of $p$ runs over the cycle sizes and $m_p$ is the number of cycles of size $p$. Thus all supports and root labels are included; no finite restriction is imposed on the new row phases.

There are no final records. Hence no regular fourth row exists outside $Z$. Conjugation preserves regularity and extension length, and Lemma~\ref{eee_generic:normal} covers every maximal-dimensional chosen $\mathrm{EEE}$ presentation.
\end{proof}

\par\addvspace{\medskipamount}\noindent
\begin{minipage}{\linewidth}
\begin{example}\label{eee_generic:example}
Put $w=e^{2\pi i/17}$. Then $H(w,w^2,w^4,w^7)$ is a regular $3\times11$ matrix of minimal Butson level $204$ with no regular fourth row. Exact substitution violates all 5108 exceptional equations, as checked in the certificate. Its entries lie in $\mu_{204}$; its dephased core contains a primitive 12th root and its entry in column $5$ of the second row is a primitive 17th root. These cross ratios force every Butson representative to have level divisible by $\operatorname{lcm}(12,17)=204$.
\end{example}
\end{minipage}
\par\addvspace{\medskipamount}

The exceptional parameters and the lower-dimensional $\mathrm{EEE}$ presentations are not settled by this theorem.

\subsection{The homogeneous FFF case}\label{sec:FFF}

\medskip
For type $\mathrm{FFF}$, each chosen quotient has one triangle and 4 antipodal pairs. Normalize the rows to $1,u,v$, put $z=u/v$, and denote their chosen partitions by $P,Q,R$. The 8 pair phases are constrained by 6 equations, explaining the positive-dimensional families in the result below.

The support equivalence includes column permutations, permutations of the 3 quotient partitions, and permutations of the 4 pair blocks within each quotient. These are not counts of inequivalent matrices. Different supports can describe the same matrix or overlapping affine families. In particular, the $\mathrm{E}$ and $\mathrm{F}$ profiles can overlap, through the alternative decompositions $3+3$ and $2+2+2$.

\medskip
Fix the columns of the first quotient so that
$$P=((0,1,2),(3,4),(5,6),(7,8),(9,10)).$$
The number of partitions of this shape is
$$\binom{11}{3}\frac{8!}{2^4\,4!}=17\,325.$$
The stabilizer of $P$, of order $3!\,2^4\,4!=2304$, has $44$ orbits on $Q$. For one representative of each orbit, test all possible $R$. Thus there are
$$44\cdot17\,325=762\,300$$
support presentations to examine. The generators used are the adjacent transpositions on the distinguished triangle, swaps within pairs, and adjacent swaps of pair blocks. Orbit sizes sum to $17\,325$.

For a support triple define
$$n_{ijk}=|P_i\cap Q_j\cap R_k|,\qquad 0\leq i,j,k\leq4,$$
where block zero is the triangle. The $5\times5\times5$ intersection tensor determines the triple up to column permutation: each cell records exactly the number of columns with the indicated 3 block memberships. The remaining group is
$$S_3\ltimes(S_4)^3,\qquad |G|=6\cdot24^3=82\,944.$$
Permuting the quotient partitions is legitimate because row permutations permute the 3 undirected row pairs; reversing a quotient preserves its cycle supports.

\medskip
Rotate $u$ and $v$ so that the first entry on each chosen triangle is $1$. Writing $w=e^{2\pi i/3}$, the triangle entries, in increasing column order, are then
$$(1,w^\sigma,w^{2\sigma}),\qquad \sigma\in\{1,-1\}.$$
For each pair $(a,b)$ with $a<b$, write its entries as $e^{2\pi ix}(1,-1)$. There are 4 pair phases in $u$ and 4 in $v$. Let $x\in(\mathbb R/\mathbb{Z})^8$ collect them.

Let $D_P,D_Q$ be the $11\times8$ integer matrices selecting the appropriate pair variable at each column, with zero rows on their triangles, and put $W=D_P-D_Q$. Let $\epsilon_P,\epsilon_Q,\epsilon_R$ be $1$ at the second entry of each pair and $0$ elsewhere. For $X=P,Q,R$, let $e_X$ be zero on the pairs and $(0,\sigma_X,2\sigma_X)$ on the triangle. There are only 8 choices of $\sigma=(\sigma_P,\sigma_Q,\sigma_R)\in\{1,-1\}^3$.

For every block $S$ of $R$, put $j_S=\min S$. For $j\in S\setminus\{j_S\}$ define
\begin{align*}
A_{(S,j)}&=W_j-W_{j_S},\\
f_{(S,j)}&=(\epsilon_R-\epsilon_P+\epsilon_Q)_j -(\epsilon_R-\epsilon_P+\epsilon_Q)_{j_S},\\
(B\sigma)_{(S,j)}&=(e_R-e_P+e_Q)_j-(e_R-e_P+e_Q)_{j_S}.
\end{align*}
The last expression defines an integer matrix $B$ with 3 columns. The complete phase system is
\[
Ax=\frac f2+\frac{B\sigma}{3}\pmod1, \qquad A\in M_{6\times8}(\mathbb{Z}),\quad B\in M_{6\times3}(\mathbb{Z}).
\]
There are 2 equations from the triangle and one from each pair. The phase of each block of $z$ is eliminated by subtraction at its first column. No additional label constraints remain.

\par\addvspace{\medskipamount}\noindent
\begin{minipage}{\linewidth}
\begin{lemma}\label{n11_noA_FFF:dimension}
Whenever a chosen support and orientation are feasible, their connected phase components are tori of dimension $8-\operatorname{rank}A\geq2$.
\end{lemma}
\end{minipage}
\par\addvspace{\medskipamount}

\begin{proof}
An integer homomorphism from $(\mathbb R/\mathbb{Z})^8$ to $(\mathbb R/\mathbb{Z})^6$ has kernel components of dimension $8-\operatorname{rank}A$. A nonempty fibre is a translate of this kernel. The row normalization has removed the 2 row phases, and each pair phase is recoverable from its first entry. Thus these parameters describe actual deformations of the normalized matrix. Remaining permutation equivalences are finite and do not reduce dimension.
\end{proof}

\medskip
The ranks and torsion components of these systems determine the complete $\mathrm{FFF}$ cover. The following theorem records the census and its largest dimension.

\par\addvspace{\medskipamount}\noindent
\begin{minipage}{\linewidth}
\begin{theorem}\label{n11_noA_FFF:main}
Every regular $3\times11$ matrix of chosen type $\mathrm{FFF}$ is equivalent to a member of a finite union of affine families whose dimensions lie between 2 and 6. Each family in the normal form has a seed of root level $6$, $12$ or $24$. In particular, seeds over $\mu_{24}$ suffice.

The dimension bound 6 is attained. The feasible chosen support triples have dimension counts
$$\begin{array}{c|rrrrr}
\text{dimension}&2&3&4&5&6\\ \hline
\text{support classes}&17\,253&2\,394&175&17&2
\end{array}$$
for a total of $19\,841$ support classes.
\end{theorem}
\end{minipage}
\par\addvspace{\medskipamount}

\begin{proof}
For each presentation the calculation constructs a Smith normal form~\cite{KB} with unimodular matrices
$$U\in\operatorname{GL}_6(\mathbb{Z}),\qquad V\in\operatorname{GL}_8(\mathbb{Z}), \qquad UAV=\operatorname{diag}(d_0,\ldots,d_{r-1},0).$$
The rectangular diagonal matrix has 6 rows and 8 columns, and $r=\operatorname{rank}A$. For all $762\,300$ presentations we verify the identity, both determinants being $\pm1$, and positivity and successive divisibility of the nonzero Smith invariants.

The zero rows require exactly
\[
(Uf)_i=0\pmod2,\qquad (UB\sigma)_i=0\pmod3 \quad (r\leq i<6).
\]
Indeed, multiplication by 6 separates the 2 coprime denominators. The binary test leaves $607\,052$ presentations; checking all 8 orientations leaves $552\,102$.

There are $552\,102$ feasible presentations, yielding $320\,290$ distinct labelled intersection tensors. Their canonical forms give the $19\,841$ support classes in Theorem~\ref{n11_noA_FFF:main}. For efficiency, first minimize a 2-axis marginal under the group, and then minimize the full tensor among transformations attaining that marginal. For fixed first 2 axis permutations, sorting the 4 pair slices minimizes the last axis permutation. The verifier uses explicit last-axis permutations instead of this sorting step. It checks every canonical minimum and a recorded group-transformation witness for each of the $320\,290$ input tensors. All $552\,102$ feasible presentations are accounted for.

For every allowed orientation all solutions of the phase system $Ax=f/2+B\sigma/3$ are
\[
\begin{aligned}
x&=Vy,\\
y_i&=\frac{3(Uf)_i+2(UB\sigma)_i+6k_i}{6d_i}, &&0\leq i<r,\quad 0\leq k_i<d_i,\\
y_i&=t_i,&&r\leq i<8,\quad t_i\in\mathbb R/\mathbb{Z}.
\end{aligned}
\]
The rows are recovered by
\begin{align*}
u_j(t)&=\exp\!\left(2\pi i\left[ \frac{\epsilon_{P,j}}2+\frac{e_{P,j}}3+(D_Px)_j\right]\right),\\
v_j(t)&=\exp\!\left(2\pi i\left[ \frac{\epsilon_{Q,j}}2+\frac{e_{Q,j}}3+(D_Qx)_j\right]\right).\nonumber
\end{align*}
Every block remains a rotated cycle for all values of $t$. Conversely, an arbitrary original triple, after normalization, determines its orientation and a solution of the phase system $Ax=f/2+B\sigma/3$. The displayed Smith-coordinate solution therefore exhausts all triples with the chosen supports.

In every feasible case $2\leq r\leq6$. The possible largest Smith invariants and their support-class counts are
$$\begin{array}{c|rrrrr}
d_{r-1}&1&2&3&4&6\\ \hline
\text{support classes}&14\,757&4\,694&286&103&1.
\end{array}$$
Setting the free parameters to zero initially gives a seed of level dividing $6d_{r-1}$. We checked every allowed orientation and every torsion branch in every feasible presentation, comprising $4\,933\,076$ phase-component presentations. The actual seed levels are only $6,12,24$. In particular, the potential extra factor of 3 in the cases with largest invariant $3$ or $6$ always cancels. This is an exact denominator computation, not a bounded search for the original phases.

For one representative of each of the $19\,841$ support classes, the complete orientation and branch counts are:
$$\begin{array}{c|rrrrr}
\text{dimension}&2&3&4&5&6\\ \hline
\text{affine presentations}&180\,512&9\,306&470&36&4.
\end{array}$$
There are $146\,608$ orientation choices and $190\,328$ connected affine presentations for these representatives. Their seed levels are $6$ in $173\,112$ cases, $12$ in $15\,040$, and $24$ in $2\,176$. Every one of these affine presentations was checked directly: the seed differences form the required cycles, and each parameter exponent is constant on each quotient block. The component counts refer to separate chosen-support descriptions, not to inequivalent matrices or components of the union.

\medskip

A second implementation solves $A\ell=12(3f+2B\sigma)\pmod{72}$ over $\mathbb Z/8$ and $\mathbb Z/9$. The certified Smith invariants divide $12$, so every consistent torus fibre meets this grid. The test agrees with the Smith calculation for all $762300$ presentations and all $8$ orientations, including ranks and torsion. The reductions and symbolic seed checks use exact integers.

\medskip
Let
$$T_0(a)=
\begin{pmatrix}
1&1&1&1\\
1&-1&a&-a\\
1&1&-1&-1
\end{pmatrix},
\qquad T_1(a)=
\begin{pmatrix}
1&1&1&1\\
1&-1&a&-a\\
1&-1&-a&a
\end{pmatrix},\qquad a\in\mathbb{T}.$$
In the normalization of \cite{regular78}, $T_0(a)$ is $F_4^a$ restricted to rows $(1,3,2)$ and columns $(1,3,2,4)$, in that order, with one-based indices. Each row quotient of either matrix splits into 2 antipodal pairs. Consequently
\[
H_h(a,b,s,t,p,q)= \left[\ F_3\ \middle|\ \operatorname{diag}(1,s,t)T_0(a)\ \middle|\ \operatorname{diag}(1,p,q)T_h(b)\ \right], \quad h=0,1,
\]
is a regular $3\times11$ matrix of type $\mathrm{FFF}$ for all $a,b,s,t,p,q\in\mathbb T$.

The first row and first column are already normalized. The entries in the first columns of the 2 4-column blocks recover $s,t,p,q$; ratios within the second row recover $a,b$. Hence the 6 parameters are independent modulo finite permutation ambiguities. This proves attainment of the dimension bound.
\end{proof}

Every $\mathrm{FFF}$ triple lies in a positive-dimensional family of regular 3-row matrices. Additional rows may constrain these parameters; the next results address extensions of the maximal-dimensional families.

\par\addvspace{\medskipamount}\noindent
\begin{minipage}{\linewidth}
\begin{lemma}\label{fff_generic:normal}
Every member of a $6$-dimensional chosen $\mathrm{FFF}$ phase presentation is equivalent to a member of $H_0(a,b,s,t,p,q)$ or $H_1(a,b,s,t,p,q)$.
\end{lemma}
\end{minipage}
\par\addvspace{\medskipamount}

\begin{proof}
The 2 maximal support classes in Theorem~\ref{n11_noA_FFF:main} have a common triangle and 2 common 4-column blocks. In local coordinates on a 4-column block, put $I=((0,1),(2,3))$ and $J=((0,2),(1,3))$. After row and column permutations, their chosen pairings for $u,v,u/v$ are
$$\begin{array}{c|cc}
&\text{first 4-column block}&\text{second 4-column block}\\\hline
h=0&(I,J,I)&(I,J,I)\\
h=1&(I,J,I)&(I,I,J).
\end{array}$$
Normalize the common triangle to $F_3$. Its 2 possible orientations are exchanged by a column permutation. On a block of type $(I,J,I)$, write $u=(s,-s,sa,-sa)$ and $v=(t,c,-t,-c)$. The first pair of $u/v$ forces $c=t$, giving $T_0(a)$ with its row phases. On a block of type $(I,I,J)$, write $u=(p,-p,pb,-pb)$ and $v=(q,-q,c,-c)$. The first pair of $u/v$ now forces $c=-qb$, giving $T_1(b)$. All 6 phases are free. These normal forms account for the 4 maximal-dimensional affine presentations of the 2 support-class representatives.
\end{proof}

For the extension problem, the Fourier triangle gives a useful restriction on new $11$-cycle quotients. We first isolate the parameter values where such a quotient is possible.

\par\addvspace{\medskipamount}\noindent
\begin{minipage}{\linewidth}
\begin{lemma}\label{fff_generic:noA}
If none of $s,p,s/p,t,q,t/q$ belongs to $\mu_{66}$, then no regular added row of $H_h(a,b,s,t,p,q)$ has an $\mathrm{A}$ quotient with an old row.
\end{lemma}
\end{minipage}
\par\addvspace{\medskipamount}

\begin{proof}
In an $\mathrm{A}$-containing triple, each non-$\mathrm{A}$ quotient has its entries in at most 2 rotated copies of $\mu_{66}$. This is the graph argument of Subsection~\ref{sec:branch_A}: the simple bipartite support graph has at most 10 vertices, minimum degree 2, and thus at most 2 components; phase ratios in each component belong to $\mu_{66}$.

The row $u$ contains $1,s,p$. If an $\mathrm{A}$ quotient occurred between a new row $x$ and either $1$ or $u$, the triple $(1,u,x)$ would force 2 of these 3 entries into the same $\mu_{66}$ coset, contrary to the hypothesis. The row $v$ contains $1,t,q$, so the same argument applied to $(1,v,x)$ excludes an $\mathrm{A}$ quotient between $x$ and $v$.
\end{proof}

We can now test the other five profiles and combine their exceptional equations with this necessary condition for an $\mathrm{A}$ quotient.

\par\addvspace{\medskipamount}\noindent
\begin{minipage}{\linewidth}
\begin{theorem}\label{fff_generic:extension}
There is an explicit proper finite union $Z\subset\mathbb T^6$ of monomial loci such that neither $H_0(a,b,s,t,p,q)$ nor $H_1(a,b,s,t,p,q)$ has a regular fourth row outside $Z$. Thus every $6$-dimensional chosen $\mathrm{FFF}$ phase presentation has generic maximum extension exactly $3$.
\end{theorem}
\end{minipage}
\par\addvspace{\medskipamount}

\begin{proof}
First exclude the 396 equations saying that one of $s,p,s/p,t,q,t/q$ belongs to $\mu_{66}$. Lemma~\ref{fff_generic:noA} then excludes all new $\mathrm{A}$ quotients. The first 2 rows of $H_0$ and $H_1$ coincide; write them as $1,u$, and write the respective third rows as $v_0,v_1$.

Use the phase-graph enumeration of Theorem~\ref{bbb_common:extension}, with 6 old parameters and root denominator $210=\operatorname{lcm}(6,2,3,5,7)$. Enumerate every decorated presentation of a proposed row $x$ with profile $\mathrm{B},\mathrm{C},\mathrm{D},\mathrm{E},\mathrm{F}$, and every prime-cycle cover of $x/u$. For each retained pair family, test every cover of $x/v_0$ and $x/v_1$. Forest integration retains arbitrary component phases. At each closing edge, keep identities, discard inconsistent constant equations, and record each nonconstant old-parameter equation before rejecting its branch.

The audit records 21043 distinct nonconstant equations. Adjoining the 396 equations above gives 21403 distinct equations, after removing repetitions. With $w=e^{\pi i/1155}$, all have the form
$$a^i b^j s^k t^l p^m q^n=w^d,\qquad d\in\mathbb Z/2310\mathbb Z,$$
where $(i,j,k,l,m,n)\ne0$ and every coefficient has absolute value at most 4. Let $Z$ be their union. It is proper, with open dense complement. The list is \texttt{exceptional\_union.txt} in the \texttt{fff\_generic} component; each line $i\ j\ k\ l\ m\ n\ d$ represents the displayed equation. This sufficient exceptional set is not asserted to be minimal.

Two implementations, using different phase-component representations and opposite cover orders, each exhaust $1480490$ initial presentations. They agree on all retained records:
$$\begin{array}{c|r|r|r|r}
\text{profile of }x&\text{presentations}&\text{against }1,u&\text{against }1,u,v_0&\text{against }1,u,v_1\\\hline
\mathrm{B}&712800&0&0&0\\
\mathrm{C}&443520&0&0&0\\
\mathrm{D}&166320&0&0&0\\
\mathrm{E}&123200&0&0&0\\
\mathrm{F}&34650&61350&0&0
\end{array}$$
The initial count for a profile is $11!/(\prod p\prod m_p!)$, with $m_p$ its number of $p$-cycles. Thus all supports and root labels are covered. The independent verifier checks literal cycle covers for every retained pair record, the complete scan agreement, and the exceptional list. It also reconstructs the phase equations of the 2 maximal support representatives, checks their unimodular identities and all 4 affine components, and verifies the displayed normal forms.

Outside $Z$, every rejected nonconstant branch is impossible. No final records survive for either third row, proving the exclusion with unrestricted unit phases. Lemma~\ref{fff_generic:normal} covers every maximal-dimensional chosen $\mathrm{FFF}$ presentation.
\end{proof}

\par\addvspace{\medskipamount}\noindent
\begin{minipage}{\linewidth}
\begin{example}\label{fff_generic:examples}
Put $w=e^{2\pi i/83}$. Both matrices
$$H_h(w,w^7,w^{51},w^{15},w^{73},w^{20}),\qquad h=0,1,$$
have minimal Butson level $498$ and no regular fourth row. Exact substitution violates all 21403 exceptional equations. The certificate checks these congruences and the chosen cycles for both matrices. Their entries lie in $\mu_{498}$, while cross ratios have orders $2$, $3$ and $83$, forcing every Butson representative to have level divisible by $498$.
\end{example}
\end{minipage}
\par\addvspace{\medskipamount}

The exceptional parameters and the lower-dimensional $\mathrm{FFF}$ presentations remain open.

\medskip
Let $w=e^{\pi i/12}$. Take the first row to consist of ones and the other 2 rows to be $w$ raised entrywise to
$$\begin{pmatrix}
0&16&8&10&22&18&6&15&3&6&18\\
0&16&22&8&10&0&12&3&0&15&12
\end{pmatrix}.$$
The following chosen supports verify regularity; column indices start at zero:
$$\begin{array}{c|c|cccc}
&\text{triangle}&\multicolumn{4}{c}{\text{pairs}}\\ \hline
u&(0,1,2)&(3,4)&(5,6)&(7,8)&(9,10)\\
v&(0,1,3)&(2,4)&(5,6)&(7,9)&(8,10)\\
u/v&(2,3,5)&(0,4)&(1,7)&(6,10)&(8,9)
\end{array}$$
For example, the first 2 rows and columns $3,7$ give the cross ratio $w^{10-15}=w^{-5}$, primitive of order $24$. Since cross ratios are preserved as a multiset under equivalence, any equivalent partial Butson representative over $\mu_m$ must have $24\mid m$.

This chosen presentation has 2 parameters. Multiply columns $5,6$ of both displayed rows by $s\in\mathbb{T}$, and columns $7,8,9,10$ of both displayed rows by $t\in\mathbb{T}$, leaving the all-ones row fixed. The same cycles still work for every $s,t$. The example is therefore not an isolated triple. The cross-ratio observation concerns this matrix's Butson equivalence; it does not by itself prove minimality of the seed level among all alternative affine families through it.


We can now state the homogeneous result in a single form. A seed level in the next table means that each family has a seed at one of the indicated levels; it does not assert minimality of every seed or disjointness of the families.

\par\addvspace{\medskipamount}\noindent
\begin{minipage}{\linewidth}
\begin{theorem}\label{report:homogeneous}
The homogeneous regular $3\times11$ matrices admit finite affine covers, up to ordinary equivalence, with the following sharp maximum dimensions and sufficient seed levels:
\[
\begin{array}{c|c|l}
\text{chosen profile}&\text{maximum dimension}&\text{seed levels}\\\hline
\mathrm{AAA}&0&11\\
\mathrm{BBB}&3&28\\
\mathrm{CCC}&4&30,45\\
\mathrm{DDD}&3&40\\
\mathrm{EEE}&4&6,12,18,24,30,36,42\\
\mathrm{FFF}&6&6,12,24.
\end{array}
\]
Every listed dimension is attained. Every six-dimensional chosen homogeneous phase presentation is of type $\mathrm{FFF}$ and is ordinarily equivalent to one of the families
\[
H_h(a,b,s,t,p,q)=
\left[\ F_3\ \middle|\ \operatorname{diag}(1,s,t)T_0(a)
\ \middle|\ \operatorname{diag}(1,p,q)T_h(b)\ \right],\qquad h=0,1,
\]
with $F_3=(w^{ij})_{i,j=0}^2$, $w=e^{2\pi i/3}$, and
\[
T_0(a)=\begin{pmatrix}1&1&1&1\\1&-1&a&-a\\1&1&-1&-1\end{pmatrix},\qquad
T_1(a)=\begin{pmatrix}1&1&1&1\\1&-1&a&-a\\1&-1&-a&a\end{pmatrix}.
\]
Each family has generic maximum regular extension exactly three, outside the explicit finite union of proper monomial loci of Theorem~\ref{fff_generic:extension}.
\end{theorem}
\end{minipage}
\par\addvspace{\medskipamount}
\begin{proof}
For $\mathrm{AAA}$, normalize the first two rows to $1,(w^j)_j$ with $w=e^{2\pi i/11}$. The third row has entries in $\mu_{11}$ after a row rotation. Its exponent vector and its difference from $(j)_j$ are permutations, so the list is finite and the dimension is zero. Three distinct rows of $F_{11}$ give attainment.

For the other five profiles, the affine covers, upper bounds and their attainment are Theorems~\ref{n11_noA_BBB:main}, \ref{n11_noA_CCC:main}, \ref{n11_noA_DDD:main}, \ref{n11_EEE:main} and~\ref{n11_noA_FFF:main}. Each solves the complete phase system and retains every root-label choice needed for coverage. Dephasing removes the continuous equivalence operations; the remaining finite choices cannot lower the dimension of an exhibited family. Since the cover is finite, its largest component dimension is also the dimension of the corresponding homogeneous locus modulo ordinary equivalence.

Only $\mathrm{FFF}$ can have dimension six. Lemma~\ref{fff_generic:normal} gives its two normal forms, and Theorem~\ref{fff_generic:extension} excludes a regular fourth row on their stated open dense sets. These arguments prove all assertions, and also prove Theorem~\ref{intro:homogeneous}.
\end{proof}

The large dimension in the last row of the table is therefore compatible with generic failure at the very next extension step. It gives no indication that a six-dimensional family of regular square matrices exists. Conversely the exceptional loci in the generic theorems must be retained when considering square completion.

\section{Maximal five-row families}\label{sec:order11part2}
\subsection{The full 4-parameter DDD family}\label{sec:DDDfull}

We determine the full regular extension problem for the $4$-parameter family in Proposition~\ref{n11_DDD:fiveconstruction}.

Put $w=e^{\pi i/210}$ and $\omega=w^{21}=e^{\pi i/10}$. Index the rows by $r=0,1,2,3,4$ and the columns by $j=0,\ldots,10$. Put
$$d_j=\begin{cases}0&0\leq j<5,\\1&5\leq j\leq10,\end{cases}$$
and write the seed as $H^0_{rj}=\omega^{h_{rj}}$, where
\[
(h_{rj})=\begin{pmatrix}
0&0&0&0&0&0&0&0&0&0&0\\
0&4&8&12&16&0&10&15&5&5&15\\
0&8&16&4&12&0&15&10&15&5&5\\
0&12&4&16&8&0&5&15&10&15&5\\
0&16&12&8&4&0&5&5&15&10&15
\end{pmatrix}.
\]
Thus the family is
\[
H(t)_{rj}=t_r^{d_j}\omega^{h_{rj}},
\qquad t_0=1,\qquad (t_1,t_2,t_3,t_4)\in\mathbb T^4.
\]
The first 5 columns form $F_5$. The last 6 are the specified 5 rows of the quaternary matrix used in Subsection~\ref{sec:DDD}. Each old quotient is a pentagon on the first block and 3 antipodal pairs on the second block. Multiplication by $t_r\bar t_s$ on the latter block preserves these pairs. Consequently every member of the family $H(t)$ is a regular partial Hadamard matrix, and every old row pair has a chosen presentation of type $\mathrm{D}$.

A regular sixth row here means a row $x\in\mathbb T^{11}$ whose quotient with each of the 5 old rows is a union of prime cycles. Such a row would in particular be orthogonal to all 5 rows. The cycle condition is required against every old row. An orthogonal but nonregular sixth row is exhibited below.

\par\addvspace{\medskipamount}\noindent
\begin{minipage}{\linewidth}
\begin{theorem}\label{n11_full:main}
For every $t\in\mathbb T^4$, the $5\times11$ matrix $H(t)_{rj}=t_r^{d_j}\omega^{h_{rj}}$, with $t_0=1$ and the seed displayed above, admits no regular sixth row.
\end{theorem}
\end{minipage}
\par\addvspace{\medskipamount}

We first establish the finite phase-graph reduction and the exact exclusion tests, with all unit parameters retained. The proof of the theorem then assembles these reductions in an exhaustive enumeration: necessary norm equations remove the old multipliers, and further cycle decompositions or $2$-variable elimination exclude the remaining families.

\medskip
Let $x$ be a putative regular sixth row. Its quotient with row zero is $x$ itself. Profile $\mathrm{A}$ is impossible: an $11$-cycle quotient, together with any other old row, would give a chosen triple containing both $\mathrm{A}$ and $\mathrm{D}$, contrary to Theorem~\ref{n11_A_branch:profiles}. The same argument excludes $\mathrm{A}$ against each of rows $1,2,3,4$. We therefore use only
$$\mathrm{B}=7+2+2,\quad \mathrm{C}=5+3+3,\quad \mathrm{D}=5+2+2+2,\quad
\mathrm{E}=3+3+3+2,\quad \mathrm{F}=3+2+2+2+2.$$

Choose one of these presentations against row zero. Its blocks are $S_0,\ldots,S_{k-1}$, where $3\leq k\leq5$. On a block of length $p$, assign the labels $0,\ldots,p-1$ bijectively, with label zero at its least column. If $j\in S_b$ has label $\ell_j$, write
\[
x_j=u_bw^{e_j},\qquad e_j=\frac{420}{p}\ell_j,\qquad u_b\in\mathbb T.
\]

For block lengths $s_1,\ldots,s_k$, with $m_p$ blocks of length $p$, the number of normalized labelled presentations is
$$\frac{11!}{\prod_i s_i\prod_p m_p!}.$$
Indeed, the number of unordered supports is $11!/(\prod_i s_i!\prod_p m_p!)$, and the number of normalized label assignments is $\prod_i(s_i-1)!$. The resulting counts are
\[
\begin{array}{c|r}
\text{profile}&\text{normalized presentations}\\\hline
\mathrm{B}&712800\\\mathrm{C}&443520\\\mathrm{D}&166320\\\mathrm{E}&123200\\\mathrm{F}&34650\\\hline
\text{total}&1480490
\end{array}
\]
These are presentation counts. A row can have more than one presentation, and no matrix-equivalence claim is attached to these numbers.

Put $q=\bar t_1$. The entries of the next quotient are
$$y_j=x_j\bar H(t)_{1j}
=u_{b(j)}w^{e_j-21h_{1j}}q^{d_j}.$$
Consider a chosen $p$-cycle of this quotient, with anchor column $a$ assigned label zero. If column $j$ receives label $\ell$, the required equation $y_j/y_a=w^{420\ell/p}$ is
\[
\frac{u_{b(j)}}{u_{b(a)}}=q^dw^v,
\quad d=d_a-d_j,
\quad v=\frac{420}{p}\ell+e_a-e_j+21(h_{1j}-h_{1a})\pmod {420}.
\]
In particular, $d\in\{-1,0,1\}$. These are monomial equations on the full unit circle; $q$ is not assumed to have finite order.

The enumeration chooses the least unused column as the anchor of the next cycle, tries each permitted prime length, and assigns the other labels to distinct unused columns in every possible way. A completed cycle is followed by the same procedure on its complement. An edge equation is checked as soon as its label is assigned. Thus an inconsistent partial assignment can be discarded immediately. Conversely, every regular quotient has at least one cycle decomposition produced by this recursion. Anchoring at the least unused column loses no decomposition, because the common phase of a cycle absorbs its label rotation.

\medskip
For a fixed pair of cycle presentations, form a graph whose vertices are the $k$ phases in $x_j=u_bw^{e_j}$. Its edges are the displayed ratio constraints. Loops and multiple edges are allowed. In each connected component, choose a spanning tree and a root. Tree integration gives
\[
u_b=v_{g(b)}q^{a_b}w^{f_b},
\]
where $g(b)$ is the component containing $b$, $v_{g(b)}$ is its free phase, $a_b\in\mathbb Z$ and $f_b\in\mathbb Z/420\mathbb Z$. Every additional edge either checks a root identity or gives an equation involving $q$ alone.

\par\addvspace{\medskipamount}\noindent
\begin{minipage}{\linewidth}
\begin{lemma}\label{n11_full:graph}
For each consistent pair of chosen presentations, one of the following occurs.

\begin{enumerate}
\item The parameter $q$ remains free, and all phase solutions have the form $u_b=v_{g(b)}q^{a_b}w^{f_b}$.
\item The parameter $q$ belongs to an explicitly determined finite list of roots of unity. Every root in the list belongs to $\mu_{25200}$.
\end{enumerate}
If the graph is connected and $q$ remains free, then, after multiplication of the whole candidate row by one unit scalar,
$$x_j=r_jq^{a_j},\qquad r_j\in\mu_{420},\qquad
\max_j a_j-\min_j a_j\leq4.$$
If $q$ is fixed and the graph has 2 components, the analogous form is
$$x_j=r_jz^{b_j},\qquad r_j\in\mu_{25200},\qquad b_j\in\{0,1\},$$
with $z\in\mathbb T$ free.
\end{lemma}
\end{minipage}
\par\addvspace{\medskipamount}

\begin{proof}
A non-tree edge closes a fundamental cycle. Substituting the tree potentials into that edge gives
$$q^d=w^v.$$
The fundamental cycle has at most $k$ edges, and each original edge has charge $-1,0$ or $1$. Hence $|d|\leq k\leq5$. If $d=0$, the equation is either an identity or an inconsistency. If $d\ne0$, it has exactly $|d|$ solutions on $\mathbb T$, all of order dividing $420|d|$. The first such equation supplies a finite list; every subsequent equation is checked on the whole list. Since
$$\operatorname{lcm}(420,840,1260,1680,2100)=25200,$$
all retained values lie in $\mu_{25200}$. If no nonzero charge occurs, $q$ is free.

In a connected graph, the difference between any 2 tree potentials is the sum of charges along the tree path between their vertices. That path has length at most $k-1$. Therefore the range of the exponents in the candidate row is at most $k-1\leq4$. Removing the one common phase gives the first displayed form. In the fixed case, all root factors in $u_b=v_{g(b)}q^{a_b}w^{f_b}$ belong to $\mu_{25200}$. With 2 components there is, after removal of a common row phase, just one ratio of component phases. Calling it $z$ gives the second form.
\end{proof}

\medskip
For an arbitrary candidate row define
\[
P_r(x)=\sum_{j=0}^4x_j\omega^{-h_{rj}},\qquad
Q_r(x)=\sum_{j=5}^{10}x_j\omega^{-h_{rj}},\qquad 1\leq r\leq4.
\]
Orthogonality to row $r$ says
$$P_r(x)+\bar t_rQ_r(x)=0.$$
Taking absolute values gives the 4 necessary equations
\[
N_r(x):=|P_r(x)|^2-|Q_r(x)|^2=0,\qquad r=1,2,3,4.
\]
They contain none of the old multipliers. We use them only as necessary equations, so no division by $P_r$ or $Q_r$ is involved. They remain valid when either sum vanishes.

Suppose now that a retained row has the form $x_j=r_jz^{a_j}$, with $|z|=1$, and put $R=\max a_j-\min a_j$. Since $\bar z=z^{-1}$ and $\bar r_j=r_j^{-1}$,
\[
z^RN_r(x)=
\sum_{j,k<5}r_j\bar r_k\omega^{-h_{rj}+h_{rk}}z^{a_j-a_k+R}
-
\sum_{j,k\geq5}r_j\bar r_k\omega^{-h_{rj}+h_{rk}}z^{a_j-a_k+R}.
\]
This is a polynomial of degree at most $2R$. In the connected free-$q$ case, it has coefficients in $\mathbb Q(\mu_{420})$ and degree at most 8. In the fixed-$q$, 2-component case, it has coefficients in $\mathbb Q(\mu_{25200})$ and degree at most 2.

The equations are particularly useful because they avoid treating 4 independent old parameters as 4 independent elimination variables. They constrain the candidate row directly. However, a nonzero polynomial on the circle does not by itself imply that its roots are roots of unity. The next step again uses full regularity.

\par\addvspace{\medskipamount}\noindent
\begin{minipage}{\linewidth}
\begin{lemma}\label{n11_full:torsion}
Fix a chosen cycle system in the present full-family calculation, retaining the old parameters as variables. Let $z$ be a remaining phase coordinate, or a monomial in those coordinates with a root-of-unity factor. If a nonzero necessary polynomial in $z$ exists, then every possible value of $z$ in a regular extension is a root of unity. More precisely:
\begin{enumerate}
\item for coefficients in $\mathbb Q(\mu_{420})$ and degree at most 8, its order divides $352800$;
\item for coefficients in $\mathbb Q(\mu_{25200})$ and degree at most 2, its order divides $50400$.
\end{enumerate}
Thus both cases are covered by the common root order $L=352800$.
\end{lemma}
\end{minipage}
\par\addvspace{\medskipamount}

\begin{proof}
Include the 4 old parameters, the row phases, and chosen cycle decompositions for every new quotient. Every equation is a monomial equation with a root-of-unity constant. The phase-system argument of Section~\ref{sec:conventions} gives a finite union of tori translated by torsion points. Adding a monomial phase coordinate, or projecting to it, preserves this description. A nonzero polynomial leaves a finite image in that coordinate. Every connected torus in the cover therefore projects to a point, and its translated image is a torsion point.

If the order is $m$, the degree of $z$ over $\mathbb Q(\mu_n)$ is $\varphi(\operatorname{lcm}(m,n))/\varphi(n)$, since adjoining $z$ gives $\mathbb Q(\mu_{\operatorname{lcm}(m,n)})$. With $n=420$, the first degree bound gives
$$\frac{\varphi(\operatorname{lcm}(m,420))}{\varphi(420)}\leq8,$$
A new prime would contribute at least $10$ to this ratio. The exponents of $2,3,5,7$ in $m$ are therefore at most $5,2,2,2$, respectively, giving $m\mid2^5 3^2 5^2 7^2=L$. For the second bound the same argument gives
$$\frac{\varphi(\operatorname{lcm}(m,25200))}{\varphi(25200)}\leq2.$$
Here $25200=2^4 3^2 5^2 7$. A new prime would contribute at least 10, since the primes $2,3,5,7$ are already present. Increasing the exponent of $3,5$ or $7$ would contribute at least 3. Only one additional factor of 2 is possible. Thus $m\mid2^5 3^2 5^2 7=50400$.
\end{proof}

\medskip
Use the splitting prime $p=1764001=5L+1$, with the specified primitive $L$th root in the code. Reduction sends each characteristic-zero $L$th root to the same power of this finite-field root, injectively on $\mu_L$. Nonzero reductions of the norm polynomials certify nonzero characteristic-zero polynomials. Their common regular roots must lie in $\mu_L$ by Lemma~\ref{n11_full:torsion}. If $G$ is the gcd of the nonzero reduced polynomials, compute $\gcd(G,Z^L-1)$. A constant gcd excludes the family. Otherwise enumerate its roots, which are distinct because $p\nmid L$, and convert them to exponents of the chosen primitive root. Every genuine regular row survives this necessary test. If all reduced norm polynomials vanish identically, this test gives no exclusion and the family is retained. Each surviving fixed candidate is then tested by exact cycle covers, as follows.

\par\addvspace{\medskipamount}\noindent
\begin{minipage}{\linewidth}
\begin{lemma}\label{n11_full:pointtest}
For this test put $w=e^{2\pi i/L}=e^{\pi i/176400}$, and let $x_j=w^{a_j}$ be a fixed row. For any old row $r\in\{1,2,3,4\}$, whether $x/H(t)_r$ can be regular for some $t_r\in\mathbb T$ can be decided by at most $391$ choices of an exponent modulo $L$, followed by finite cycle-cover tests.
\end{lemma}
\end{minipage}
\par\addvspace{\medskipamount}

\begin{proof}
An $11$-cycle is impossible: its first 5 entries would have pairwise ratios in $\mu_L\cap\mu_{11}=\{1\}$, whereas the entries of an $11$-cycle are distinct. Set $b_j=a_j-(L/20)h_{rj}$. The quotient has entries
$$w^{b_j}q_r^{d_j},\qquad q_r=\bar t_r.$$
Suppose a cycle in a regular cover uses a column $j<5$ and a column $k\geq5$. Its length is $\ell\in\{2,3,5,7\}$. The 2 entries have different labels in that cycle, so for some $1\leq h<\ell$ their ratio is $w^{hL/\ell}$. Therefore
$$q_r=w^{b_j-b_k-hL/\ell}.$$
There are at most
$$5\cdot6\bigl((2-1)+(3-1)+(5-1)+(7-1)\bigr)=390$$
values on these lists, before repetitions are removed.

If no cycle crosses between the 2 column blocks, each cycle lies entirely within one block. Changing $q_r$ then only rotates complete cycles. In that case $q_r=1$ works, and it suffices to add exponent zero to the list. This also deals with the possibility that an initially given multiplier was not a root of unity.

For each candidate exponent, the 11 quotient entries are specified by residues modulo $L$. A set of $\ell$ columns is a prime cycle precisely when, relative to any one of its entries, its residues are $0,L/\ell,\ldots, (\ell-1)L/\ell$ without repetition. To cover all columns, choose the least uncovered column and try every cycle containing it, then recurse on the remaining columns. There are only $2^{11}$ subsets. This exact finite procedure decides the existence of a cover.
\end{proof}

Return now to $w=e^{\pi i/210}$ and $\omega=w^{21}$.

For exclusion it is enough that one of rows $1,2,3,4$ fails this test. When testing a fixed candidate row, the old multipliers may be chosen independently. Allowing a multiplier different from the one initially used to generate the candidate enlarges the test and is therefore safe for exclusion. A row that fails even this enlarged test cannot extend any member of the family $H(t)$.

\medskip
We next describe the cases not yet covered by the one-variable analysis. If $q$ is fixed and there are $c\geq3$ connected components, the row has the form
$$x_j=r_jv_{g(j)},\qquad g(j)\in\{0,\ldots,c-1\}.$$
For every such retained case in this calculation, the normalized coefficients $r_j$ actually lie in $\mu_{420}$. This divisibility condition is checked on the saved exponents; it is not inferred from the looser a priori bound $\mu_{25200}$. We can therefore impose a complete cycle decomposition against row 2, using its free multiplier and the same graph construction. A single connected component or a fixed multiplier with 2 components returns to the preceding finite or univariate tests. Any other result is retained and handled by the rules below.

If $q$ remains free and the graph has $c\geq2$ components, normalize each component at its first column. This absorbs a common power of $q$ and a common root factor into that component's phase. The projected row becomes
\[
x_j=r_jv_{g(j)}q^{a_j},\qquad r_j\in\mu_{420}.
\]
If all $a_j$ vanish, $q$ has disappeared from the row. With 2 components this is a one-variable family, and with more components we impose the next old-row cycle decomposition as just described. The initial old-row constraints need not be retained in this new parametrization: their omission enlarges the set being excluded.

There is one further normalization available when the original presentation against row zero respects the 5-column and 6-column blocks.

\par\addvspace{\medskipamount}\noindent
\begin{minipage}{\linewidth}
\begin{lemma}\label{n11_full:gauge}
Suppose every cycle in the chosen presentation of $x$ against row zero lies entirely within one of the 2 column blocks. For $\lambda\in\mathbb T$, multiply the last 6 entries of $x$ and all 4 parameters $t_1,t_2,t_3,t_4$ by $\lambda$. This preserves the existence of a regular extension with that chosen row-zero presentation.

In particular, if in $x_j=r_jv_{g(j)}q^{a_j}$ there are integers $c_g$ and $h$ such that
$$a_j=c_{g(j)}+hd_j,$$
then the powers of $q$ can be removed from the projected row by component-phase changes and this common block multiplication.
\end{lemma}
\end{minipage}
\par\addvspace{\medskipamount}

\begin{proof}
Against old rows $1,2,3,4$, both numerator and denominator of the quotient acquire the same factor $\lambda$ on the last 6 columns. These quotients are unchanged. Against row zero, the cycles in the first block are unchanged and the cycles in the second block are rotated as whole cycles by $\lambda$. Thus regularity is preserved. Applying $\lambda^{-1}$ proves the converse.

For the last assertion absorb $q^{c_g}$ into $v_g$ and take $\lambda=q^{-h}$. The candidate row then has entries $r_jv_{g(j)}$ after renaming the component phases. The old parameters are still arbitrary points of $\mathbb T^4$, so their simultaneous change introduces no restriction.
\end{proof}

This normalization is applied only after checking that no original row-zero block crosses the 2 column blocks, and after checking the displayed equality of exponents at all 11 columns. A possible equality is tested for each $h$ between $-4$ and $4$. Failure to find one does not discard the family. In the actual calculation this rule applies to $360$ retained records of profile $\mathrm{D}$.

After these reductions, the only families requiring a 2-variable calculation have 2 phase components and a genuine dependence on $q$. Up to a common row phase they are
\[
x_j=r_jq^{a_j}s^{b_j},\qquad r_j\in\mu_{420},\quad b_j\in\{0,1\},\quad
\max a_j-\min a_j=1.
\]
There are $240$ such records, all arising from profile $\mathrm{D}$. Their coefficients and exponents are saved separately. The asserted exponent range is checked for each record, rather than assumed for all phase graphs.

\medskip
Substituting $x_j=r_jq^{a_j}s^{b_j}$ in the norm equations and multiplying by $qs$ gives $4$ polynomials of bidegree at most $(2,2)$ over $\mathbb Z[w]$. The declared-degree quadratic Sylvester determinants give necessary equations of degree at most $8$ in either coordinate. This remains valid at vanishing leading coefficients; an identically zero determinant is ignored. For each nonzero projection, Lemma~\ref{n11_full:torsion} bounds every regular root by $L$. The reduced eliminant gcd with $Z^L-1$ is constant for every one of the $240$ families, excluding them all. The audit expands the Sylvester determinant directly instead of using the quadratic formula.

\begin{proof}[Proof of Theorem~\ref{n11_full:main}]
We apply the preceding reductions to every possible regular sixth row. The first recursion exhausts the $1480490$ normalized presentations in the presentation-count table in Subsection~\ref{sec:DDDfull}, and for each one exhausts the compatible chosen cycle covers against row one. Its complete-cover totals are
\[
\begin{array}{c|r|r|r|r}
&\multicolumn{4}{c}{\text{dimension before projecting to the candidate row}}\\
\text{profile}&0&1&2&3\\\hline
\mathrm{B}&0&0&0&0\\
\mathrm{C}&204560&4320&48&0\\
\mathrm{D}&7346080&173420&2022&12\\
\mathrm{E}&1619200&7040&0&0\\
\mathrm{F}&21095040&417280&2400&0
\end{array}
\]
Here the recorded dimension is $c-1$ when $q$ is fixed and $c$ when $q$ is free, where $c$ is the number of phase-graph components. A fixed state can carry more than one allowed value of $q$; the table counts that state once. These states need not be disjoint, and the dimensions are those of their parametrizations before any variable disappearing from the row is removed. There are no states of higher dimension.

Within each initial presentation, repeated fixed-row or one-variable records are identified by their exact coefficient and exponent data. The primary arithmetic calculation then gives
$$\begin{array}{c|r|r|r|r}
\text{profile}&\text{finite row tests}&\text{modular norm survivors}&
\text{retained families}&\text{regular rows}\\\hline
\mathrm{B}&0&0&0&0\\
\mathrm{C}&125600&0&48&0\\
\mathrm{D}&2849120&840&2064&0\\
\mathrm{E}&1221920&0&0&0\\
\mathrm{F}&16365120&0&2440&0
\end{array}$$
A finite row test in this table is a row passed to the norm equations $N_r(x)=0$, including rows produced by the univariate reduction. Repetitions between different initial presentations are allowed. The $840$ modular norm survivors in profile $\mathrm{D}$ all fail the exact regularity test of Lemma~\ref{n11_full:pointtest}. Thus they are exclusions, not exceptional orthogonal extensions.

The retained-family calculation has the following totals:
$$\begin{array}{c|r|r|r|r|r}
\text{profile}&\begin{gathered}\text{input}\\\text{records}\end{gathered}&\begin{gathered}\text{next-row}\\\text{refinements}\end{gathered}&
\begin{gathered}\text{finite row}\\\text{tests}\end{gathered}&\begin{gathered}\text{norm}\\\text{survivors}\end{gathered}&\begin{gathered}\text{2-variable}\\\text{records}\end{gathered}\\\hline
\mathrm{C}&48&0&0&0&0\\
\mathrm{D}&2064&912&13620&240&240\\
\mathrm{F}&2440&2440&240&0&0
\end{array}$$

Finally, the 2-variable calculation has
$$\begin{array}{l|r}
\text{records}&240\\
\text{exclusions by the homogeneous span test}&0\\
\text{exclusions by the torsion gcd test}&240\\
\text{remaining root pairs}&0\\
\text{unresolved records}&0.
\end{array}$$

For a hypothetical extension, choose its covers against rows $0,1$. The exhaustive recursion includes them and the phase equations cannot reject them. The resulting row is fixed, univariate, or one of the retained families. Norm tests are necessary; every finite regular projection is included by the torsion bound. Component normalization preserves the row family, and discarded old constraints only enlarge it. Further cycle covers include the actual covers of the extension. Every branch reaches an exact rejection or the final two-variable exclusion, with no unresolved record. This proves Theorem~\ref{n11_full:main}.

The complete calculation is repeated with splitting prime $2116801=6L+1$, another phase-graph implementation, coefficient convolution for norm polynomials, and direct Sylvester determinants. The structural counts agree; both runs leave no regular row or unresolved family. The separate integer verifier checks the seed, presentation counts, family accounting and every saved finite-row rejection. This proves the uniform assertion for $H(t)$, and hence Theorem~\ref{intro:five}, whose displayed six-column block is the same seed written in fourth roots.
\end{proof}

\medskip
At the fourth-root parameter points, and on an open dense parameter set, even regularity against the first row alone suffices for exclusion.

\par\addvspace{\medskipamount}\noindent
\begin{minipage}{\linewidth}
\begin{proposition}\label{n11_full:onesided}
For every $t\in\mu_4^4$, there is no unimodular row orthogonal to all $5$ rows of $H(t)$ and regular against its all-ones row. The same assertion holds on an open dense subset of $\mathbb T^4$ whose complement has Haar measure zero.
\end{proposition}
\end{minipage}
\par\addvspace{\medskipamount}

\begin{proof}
We begin at the seed $t=(1,1,1,1)$.
\medskip
A candidate sixth row $x$ that is regular against the all-ones row is a union of rotated prime cycles. Its profile is one of $\mathrm{A},\ldots,\mathrm{F}$ from Section~\ref{sec:conventions}. Composite cycles can be refined into prime cycles, so they add no cases.

For each support partition $S_1,\ldots,S_k$, enumerate all cycle labellings. The least column of each block receives label zero; its arbitrary cycle phase absorbs this normalization. A $p$-cycle then has $(p-1)!$ labellings. Blocks of equal size are ordered by their least column to remove only the irrelevant block-permutation duplication. No equivalence-class claim is made for these support and label counts.

Write $x_j=u_br_j$ for $j\in S_b$, where $|u_b|=1$ and $r_j$ is the prescribed root on its block. Orthogonality to the first row is automatic. The remaining 4 equations are $A(t)u=0$, with
$$A(t)_{r,b}=\sum_{j\in S_b}r_j\bar H(t)_{rj},\qquad r=1,\ldots,4.$$
The $u_b$ are arbitrary continuous phases throughout. There is no restriction to a finite root grid for the candidate sixth row.

\medskip
If $k\le4$ and $A$ has column rank $k$, there is no nonzero kernel vector. For any choice of $k-1$ rows, form its signed maximal-minor vector $c$. If $Au=0$ and every $|u_b|=1$, then either all these minors vanish or their squared absolute values are all equal. Indeed, a nonzero $c$ spans the kernel of this $(k-1)\times k$ submatrix. Thus unequal values of $c_b\bar c_b$ give an exact obstruction, even if the full matrix has larger rank.

At the seed, coefficients for profiles $\mathrm{B}$--$\mathrm{F}$ lie in $\mathbb{Z}[w]$. Reduce at primes $421,3361,4621$, using respectively primitive $420$th roots $2,3091,2048$. Complex conjugation is implemented by replacing the root by its inverse, not by treating conjugation as the identity in the finite field. A nonzero reduced determinant or norm difference proves that the exact cyclotomic integer is nonzero. Reduction is used only for exclusions; a vanishing reduction is never treated as a positive solution.

Every accepted modular exclusion is separately checked by direct Leibniz determinant expansion, independently of Gaussian elimination. All computations in the certificate path use exact integers or rationals. The counts are:
$$\begin{array}{c|rrrrr}
\text{profile}&\text{supports}&\text{labelled cases}&\text{rank excluded}&\text{norm excluded}&\text{left}\\\hline
\mathrm{B}&990&712800&686876&25924&0\\
\mathrm{C}&4620&443520&443520&0&0\\
\mathrm{D}&6930&166320&156604&9716&0\\
\mathrm{E}&15400&123200&120656&2544&0\\
\mathrm{F}&17325&34650&0&34636&14
\end{array}$$
There are $1\,480\,490$ labelled presentations in total, of which $1\,480\,476$ are excluded by these modular tests.

Profile $\mathrm{A}$ is excluded without enumeration. Over $\mathbb Q(\mu_{20})$, the minimal polynomial of a primitive eleventh root is $1+X+\cdots+X^{10}$: the extension degree is $\varphi(220)/\varphi(20)=80/8=10$. A putative 11-cycle orthogonal to the second seed row would therefore force its 11 coefficients, the permuted conjugate entries of that row, to be equal. They are not. This covers all $10!$ labellings of $\mathrm{A}$.

\medskip
In all $14$ remaining cases, exact arithmetic in $\mathbb Q(\mu_{60})$ gives rank 3. The computation uses
$$\Phi_{60}(X)=X^{16}+X^{14}-X^{10}-X^8-X^6+X^2+1.$$
A verified 2-column basis describes the complete kernel. In 10 cases every kernel vector has zero triangle-phase coordinate, which excludes a unimodular vector.

In the other 4 cases, use the 2 free coordinates $s,t$, which themselves must have absolute value one, and put $z=t/s$. For a kernel coordinate $as+bt$, the condition $|as+bt|=1$ gives
$$\bar abz+a\bar b\,\bar z=1-|a|^2-|b|^2.$$
These linear equations in $z,\bar z$ are inconsistent. The stored certificates give exact linear combinations yielding $0=c$ with $c\ne0$. The verifier checks the kernel identities, the rank via a separate modular lower bound, and the inconsistency identities directly. This proves the assertion at the seed.

For the generic assertion, each seed rank or norm exclusion supplies a necessary Laurent polynomial nonzero at the seed. For the $14$ degenerate presentations, the cofactor norms are unequal at $(1,1,1,i)$ modulo $421$. Profile $\mathrm{A}$ is handled by its nonzero seed orthogonality polynomial. The union of these finitely many proper zero sets has Haar measure zero. For the fourth-root grid the complete scan excludes $240$ points immediately; at the other $16$ points, its $56$ remaining presentations fail by a forced zero coordinate ($48$) or inconsistent unit norms ($8$). These grid certificates are retained in \texttt{n11\_DDD\_exceptional\_circles}.
\end{proof}

\medskip
\par\addvspace{\medskipamount}\noindent
\begin{minipage}{\linewidth}
\begin{example}
The stronger conclusion in Proposition~\ref{n11_full:onesided} does not hold at every parameter. Put $w=e^{2\pi i/5}$, $d=\sqrt{11}$, and
$$q=\frac{-6+3i+(1+2i)d}{10},\qquad a=\frac{-7+5id}{18}.$$
Both have modulus one: the squared moduli of their numerators are $100$ and $324$, respectively. Set
$$t=(\bar q,1,-i,-1)$$
and append to $H(t)$ the row
$$x=(1,aw,aw^{2},aw^{3},aw^{4},a,-ia,i,ia,-i,-1).$$
This is an exact unimodular sixth row orthogonal to the 5 old rows. The quotients with old rows $0,2,3,4$ have the following $\mathrm{D}$ decompositions, with columns numbered from zero:
$$\begin{array}{c|c|ccc}
\text{old row}&\text{pentagon}&\multicolumn{3}{c}{\text{antipodal pairs}}\\\hline
0&\{1,2,3,4,5\}&\{0,10\}&\{6,8\}&\{7,9\}\\
2&\{1,2,3,4,5\}&\{0,9\}&\{6,8\}&\{7,10\}\\
3&\{1,2,3,4,8\}&\{0,10\}&\{5,6\}&\{7,9\}\\
4&\{1,2,3,4,6\}&\{0,7\}&\{5,8\}&\{9,10\}.
\end{array}$$
The quotient with old row $1$ is
$$v=(1,a,a,a,a,aq,iaq,-q,aq,-q,-iq).$$
Its sum is zero, but none of its entries form an antipodal pair and its entries belong to $L=\mathbb Q(i,\sqrt{11})$. The quadratic subfields of $L$ are $\mathbb Q(i),\mathbb Q(\sqrt{11}),\mathbb Q(\sqrt{-11})$, excluding $\mu_3$; its biquadratic Galois group excludes $\mathbb Q(\mu_5)$; degrees exclude $\mu_7$ and $\mu_{11}$. Hence this quotient contains no prime cycle.

These assertions can also be read directly from the formulas. The 6 distinct values occurring in $v$ are $1,a,aq,iaq,-q,-iq$. Substitution of the displayed $a,q$ into the $21$ sums obtained by choosing 2 of these 6 values with repetition shows that none is zero; coefficients are compared in the basis $1,i,d,id$. Orthogonality to row $1$ follows from
$$1+4a+(2+i)aq-(2+i)q=0,$$
which is verified in the same basis. The 4 displayed cycle decompositions prove the other scalar-product identities. Thus orthogonality, even together with regularity against $4$ of the old rows, does not imply regular extension.
\end{example}
\end{minipage}
\par\addvspace{\medskipamount}
\subsection{A 3-parameter 5-row family from the DDD census}\label{sec:DDDpositivefive}

A particular three-dimensional $\mathrm{DDD}$ presentation has 3 different pentagon supports. Adding 2 fixed rows gives the family below. We determine its full sixth-row extension problem; extensions of its first 3 rows are treated in Subsection~\ref{sec:DDDgenerictriple}.

\medskip
In this subsection rows and columns are numbered from zero. Put $w=e^{\pi i/30}$ and let $a,b,c\in\mathbb T$. Define
\[
G(a,b,c)_{ij}=w^{E_{ij}}T_{ij}(a,b,c),\qquad
E=\begin{pmatrix}
0&0&0&0&0&0&0&0&0&0&0\\
0&12&24&36&48&3&33&30&0&0&30\\
0&24&48&12&51&36&21&0&30&0&30\\
0&36&12&48&24&24&24&44&44&4&4\\
0&48&36&24&12&12&12&52&52&32&32
\end{pmatrix}.
\]
The parameter factors are one in columns $0,\ldots,6$. Their remaining values are
\[
\begin{array}{c|cccc}
 &7&8&9&10\\\hline
T_{0j}&1&1&1&1\\
T_{1j}&ab/c&ab/c&a&a\\
T_{2j}&b&b&c&c\\
T_{3j}&1&1&1&1\\
T_{4j}&1&1&1&1.
\end{array}
\]
Write the rows as $1,u,v,s,t$, in this order. The rows $s,t$ are fixed, belong to $\mu_{15}^{11}$, and satisfy $s=t^2$ entrywise. Thus $(1;s;t)$ is a fixed 3-row submatrix throughout the family.

\par\addvspace{\medskipamount}\noindent
\begin{minipage}{\linewidth}
\begin{proposition}\label{n11_positive:regular}
For every $(a,b,c)\in\mathbb T^3$, the matrix $G(a,b,c)$ is a regular $5\times11$ partial Hadamard matrix. It has a chosen decomposition with 7 $\mathrm{D}$ quotients and 3 $\mathrm{C}$ quotients. Its first 3 rows have chosen type $\mathrm{DDD}$, and its rows $1,s,t$ have chosen type $\mathrm{CCC}$.

The 3 parameters describe a 3-dimensional family of dephased matrices. The seed $G(1,1,1)$ has minimal Butson level $60$ under ordinary equivalence.
\end{proposition}
\end{minipage}
\par\addvspace{\medskipamount}

\begin{proof}
For a pair $i<k$, use the quotient of row $i$ by row $k$. Each ordered block $(j_0,\ldots,j_{p-1})$ in the following table means that the corresponding quotient entries are
$$q_{j_h}=q_{j_0}w^{60h/p},\qquad 0\le h<p.$$
The order within a block therefore specifies the root labels, as well as the support. Every line partitions all 11 columns.
\begin{center}
\begin{tabular}{@{}ccl@{}}
\toprule
Rows & Profile & Ordered cycle blocks\\
\midrule
$0,1$ & $\mathrm{D}$ & $(0,4,3,2,1),\ (5,6),\ (7,8),\ (9,10)$\\
$0,2$ & $\mathrm{D}$ & $(0,2,5,1,3),\ (4,6),\ (7,8),\ (9,10)$\\
$0,3$ & $\mathrm{C}$ & $(0,3,1,4,2),\ (5,9,7),\ (6,10,8)$\\
$0,4$ & $\mathrm{C}$ & $(0,1,2,3,4),\ (5,7,9),\ (6,8,10)$\\
$1,2$ & $\mathrm{D}$ & $(0,6,3,2,1),\ (4,5),\ (7,9),\ (8,10)$\\
$1,3$ & $\mathrm{D}$ & $(0,2,4,1,3),\ (5,6),\ (7,8),\ (9,10)$\\
$1,4$ & $\mathrm{D}$ & $(0,3,1,4,2),\ (5,6),\ (7,8),\ (9,10)$\\
$2,3$ & $\mathrm{D}$ & $(0,5,3,2,1),\ (4,6),\ (7,8),\ (9,10)$\\
$2,4$ & $\mathrm{D}$ & $(0,2,5,1,3),\ (4,6),\ (7,8),\ (9,10)$\\
$3,4$ & $\mathrm{C}$ & $(0,4,3,2,1),\ (5,9,7),\ (6,10,8)$\\
\bottomrule
\end{tabular}
\end{center}
To check a block, subtract its anchor exponent from the other quotient exponents in the defining matrix for $G(a,b,c)$. The differences are $0,60/p,\ldots,(p-1)60/p$ modulo $60$, in the displayed order. The parameter monomial in the quotient is constant on that block. For example, on the pairs $(7,9)$ and $(8,10)$ in $u/v$, the respective parameter quotients are both $a/c$ at both positions. Every listed cycle consequently sums to zero for all 3 parameters. This proves all 10 orthogonality and regularity assertions.

The matrix is dephased. Its entries in positions $(1,9),(2,7),(2,9)$ are respectively $a,b,c$, so the displayed parameter map is injective. The remaining equivalences between dephased matrices are induced by finitely many row and column permutations. They do not lower the dimension of the family.

At the seed all entries lie in $\mu_{60}$. The entries in positions $(1,5)$ and $(3,7)$ are $w^{3}$ and $w^{44}$, of orders $20$ and $15$. Each dephased entry is a row-and-column cross ratio, using the first row and first column. The set of such cross ratios over all row and column pairs is preserved by ordinary equivalence. If an equivalent representative had all entries in $\mu_m$, both these orders would divide $m$. Hence $60\mid m$, proving minimality.
\end{proof}

The pentagon supports of $u,v,u/v$ are respectively
$$\{0,1,2,3,4\},\qquad\{0,1,2,3,5\},\qquad\{0,1,2,3,6\}.$$
Together with the pairs in the table, they give the chosen support recorded as class $3286$ in the $\mathrm{DDD}$ registry. The formulas for $u,v$ are one labelled 3-dimensional presentation from its complete expansion. The fixed rows $s,t$ provide a 5-row extension over the entire parameter torus. No assertion about all label choices on that support class, or about all possible fourth and fifth rows over this presentation, follows from this construction alone.

To exclude a sixth row, we first enumerate all rows regular against the fixed $\mathrm{CCC}$ submatrix $(1;s;t)$. We then impose regularity against $u(a,b,c)$ and $v(a,b,c)$ using integer character equations. In the second part the old parameters are variables, not sampled values. In particular, the calculation retains parameter subtori, their translates, and all their torsion components.

\medskip

For the cycle enumeration put $w=e^{\pi i/420}$. A new row cannot have an $\mathrm{A}$ quotient with a fixed core row, since that would give a triple containing both $\mathrm{A}$ and $\mathrm{C}$. Enumerate the $1480490$ first presentations of profiles $\mathrm{B}$--$\mathrm{F}$ as in the presentation-count table in Subsection~\ref{sec:DDDfull}. Write $x_j=z_{b(j)}w^{e_j}$. For the quotient by a fixed row $s$, a cycle anchored at $h$ imposes
\[
\frac{z_{b(j)}}{z_{b(h)}}=w^{840\ell/p+e_h-e_j}\frac{s_j}{s_h}.
\]
Tree integration gives every phase solution; closing edges check constant roots. Normalize each component at its first column and merge equal records. Imposing regularity against $1,s$ gives $859958$ records. Repeating against $t$ gives
\[
\begin{array}{c|rrrr|r}
\text{free row phases}&0&1&2&3&\text{total}\\\hline
\text{records}&87480&12780&522&6&100788.
\end{array}
\]
All phase groups are retained. Necessary quotient-sum and 2-group norm filters are used only for nonzero modular exclusions; the complete cycle recursion supplies the final lists. Thus every possible row over the fixed $\mathrm{CCC}$ core is included.

\medskip
Write $q=(a,b,c)$ and $q^n=a^{n_1}b^{n_2}c^{n_3}$ for $n\in\mathbb Z^3$. During the remaining refinement a candidate is recorded in the form
\[
x_j=z_{b(j)}w^{e_j}q^{n_j},\qquad
q^{r_k}=w^{d_k}\quad(1\le k\le m).
\]
The vectors $n_j,r_k$ are integral. Initially $n_j=0$ and there are no equations on $q$, because the list was obtained over the fixed core. The variables $z_b$ are again arbitrary unit phases. One overall phase is immaterial for regular extension.

A cycle equation for $x/u$ or $x/v$ is an equation between 2 such free phases with a root-of-unity coefficient and a monomial in $a,b,c$. Forest integration now carries 4 integers on each edge: 3 character exponents and one root residue. An edge joining 2 different components expresses one component phase in terms of the other and these known factors. An edge within a component instead gives a character equation involving $q$ alone. These equations must be kept, since they can define exceptional parameters where a previously unavailable cycle becomes possible.

For clarity, suppose the old row in the current test has entries
$$y_j=w^{f_j}q^{m_j}.$$
For a candidate block anchored at $h$ with label $\ell$ at $j$, the equation is
\[
\frac{z_{b(j)}}{z_{b(h)}}=
w^{\,840\ell/p+e_h-e_j+f_j-f_h}
q^{n_h-n_j+m_j-m_h}.
\]
This is both necessary and sufficient for the prescribed ratio of the 2 quotient entries. After all labels of all blocks have been assigned, the entire chosen cycle cover is satisfied.

\par\addvspace{\medskipamount}\noindent
\begin{minipage}{\linewidth}
\begin{lemma}\label{n11_positive:integerconsistency}
A finite system $q^{r_k}=w^{d_k}$ is soluble in $\mathbb T^3$ if and only if every integral relation $\sum_k c_kr_k=0$ satisfies $\sum_k c_kd_k=0\pmod{840}$.

Consistency can be decided by integer row operations on the augmented rows $(r_k,d_k)$, reducing only the last coordinate modulo $840$. The nonzero coefficient rows may be put in row Hermite form. A zero coefficient row with nonzero last coordinate is inconsistent; if no such row occurs, the remaining equations are soluble. No saturation of the coefficient lattice is permitted in this procedure.
\end{lemma}
\end{minipage}
\par\addvspace{\medskipamount}

\begin{proof}
Any integral relation among the left-hand characters forces the stated relation among the constants, proving necessity. To prove sufficiency, perform unimodular integer row operations on the coefficient matrix, carrying the same operations on the constants. Euclidean elimination produces independent nonzero rows and zero rows. If all the zero rows have zero constant residue, the system is equivalent to its independent rows.

An integral matrix with $m$ independent rows induces a surjection $\mathbb T^3\to\mathbb T^m$. For instance, Smith normal form reduces the map, by invertible integer changes of coordinates, to multiplication by positive integers on $m$ coordinates. Every unit complex number has a root of each positive integral order, so those maps are surjective. The retained independent equations are therefore soluble.

The positive integers in Smith normal form also explain why division or rational row reduction is inappropriate here. An equation such as $a^2=w^d$ has both square roots as solutions. Replacing it by a single equation for $a$ would discard a torsion component. Integer row operations retain the original equation and hence all its roots. The consistency calculation needs only the row lattice and constants, without selecting or enumerating individual roots.
\end{proof}

At each successful cycle cover, normalize the phase components as in the fixed-core stage, now absorbing a root times a monomial into the component phase at its first column. Put the old-parameter equations in row Hermite form, and reduce the monomial exponents using these equations. If $q^r=w^d$, replacing $n$ by $n-kr$ adds $kd$ to the root exponent. Thus the reduction preserves the same family, including its exceptional parameter locus.

\medskip
Starting from all $100788$ fixed-core records, impose every cycle cover against $u(a,b,c)$. A quotient of profile $\mathrm{A}$ is again impossible, now because $u/1$ has profile $\mathrm{D}$ and Theorem~\ref{n11_A_branch:profiles} permits no chosen triple containing both $\mathrm{A}$ and $\mathrm{D}$. The same argument applies later to $v$.

The cover search has a finite, explicit recursion. Select the least remaining column as the next anchor, choose a prime length from $2,3,5,7$, and successively assign the other labels to distinct remaining columns. Reject a branch only when its phase equations are inconsistent. Continue until every column has been used. Taking the largest remaining column instead gives the alternative traversal used in the audit. Every cycle cover has an anchor and root-label order in either traversal, so the recursion loses no possible presentation. Repeated entries occupy separate column positions throughout.

The first refinement yields $118030$ normalized records. Their distribution by the rank $r$ of the old-parameter constraints and the number $d$ of additional row phases is
\[
\begin{array}{cc|r}
r&d&\text{records regular against }1,s,t,u\\\hline
0&0&342\\
0&1&4\\
1&0&11436\\
1&1&48\\
2&0&106128\\
2&1&72\\\hline
&&118030.
\end{array}
\]
Here the old-parameter locus has dimension $3-r$; $d$ counts only the additional phases in $x$ after its row phase has been removed. 

Impose every cycle cover against $v(a,b,c)$ on every one of these records. The same integer consistency test applies, and there is no surviving record. The complete sequence of retained-list sizes is therefore
\[
\begin{array}{l|r}
\text{old rows against which regularity is imposed}&\text{records}\\\hline
1,s&859958\\
1,s,t&100788\\
1,s,t,u&118030\\
1,s,t,u,v&0.
\end{array}
\]
The final step is an exhaustive character-equation calculation. It uses neither numerical phase sampling nor a restriction of $a,b,c$ or the candidate row phases to roots of bounded order.

\par\addvspace{\medskipamount}\noindent
\begin{minipage}{\linewidth}
\begin{theorem}\label{n11_positive:main}
For every $(a,b,c)\in\mathbb T^3$, the matrix $G(a,b,c)$ admits no regular sixth row. Thus every matrix $G(a,b,c)$ is maximal for regular row extension and cannot occur as 5 rows of a regular square matrix of order 11.
\end{theorem}
\end{minipage}
\par\addvspace{\medskipamount}

\begin{proof}
Suppose that a regular sixth row exists at some parameter value. Its quotients against the fixed core have one of the profiles $\mathrm{B}$ through $\mathrm{F}$, and the fixed-core enumeration places it in one of the $100788$ families of the fixed-core dimension table. Its regular quotient against $u$ supplies a cycle cover in the next recursion. The corresponding phase equations are satisfied by the given parameters and row phases, so the consistency test cannot reject that branch. After normalization the row belongs to one of the $118030$ records of the refinement-count table.

Its regular quotient against $v$ similarly supplies a final consistent cycle cover. This contradicts the empty last list in the displayed successive record counts. Proposition~\ref{n11_positive:regular} supplies the existing 5 rows at every parameter, so 5 is the maximum number of rows in any regular extension. A regular square containing them would provide a regular sixth row after deletion of its other rows, which is impossible.
\end{proof}

\medskip

The fixed-core and parameter refinements are repeated with explicit component potentials, reversed cycle traversal and an independent integer diagonal consistency test. The complete normalized lists agree. A separate verifier checks the old family, the minimal level, and every retained row symbolically on its full parameter locus. Exhaustiveness uses the complete recursions above. The conclusion concerns these $5$ prescribed rows; deleting $s,t$ gives the larger extension problem treated in Subsection~\ref{sec:DDDgenerictriple}.

\medskip
There is also a direct way to distinguish generic members of the present construction from every member of the earlier 4-parameter family. Say that a 5-row matrix has a common pentagon if some 5-column subset supports a $5$-cycle in every quotient of 2 distinct rows. This property is invariant under ordinary equivalence: column factors cancel in each quotient, row factors rotate its cycles, and row and column permutations merely permute the conditions and the common subset. Every matrix $H(t)$ has this property, since its first 5 columns form $F_5$.

\par\addvspace{\medskipamount}\noindent
\begin{minipage}{\linewidth}
\begin{proposition}\label{n11_positive:separation}
Let
$$V=\{(a,b,c)\in\mathbb T^3:
 a,b,c,ab/c,b/c\notin\mu_{60}\}.$$
The set $V$ is open and dense, with complement of Haar measure zero. For every $(a,b,c)\in V$, the matrix $G(a,b,c)$ has no common pentagon and is not ordinarily equivalent to any matrix in the family $H(t)$.
\end{proposition}
\end{minipage}
\par\addvspace{\medskipamount}

\begin{proof}
Each forbidden condition is a nonconstant character of $\mathbb T^3$ taking one of finitely many values. Its solution set is a finite union of proper translated subtori, which is closed and has Haar measure zero. The same holds for the finite union of all 5 forbidden conditions. Its complement is open and dense.

For this comparison put $w=e^{\pi i/30}$, as in the defining matrix for $G(a,b,c)$. For the row $u$, the parameter factors on the 3 column groups $\{0,\ldots,6\}$, $\{7,8\}$ and $\{9,10\}$ are $1,ab/c,a$. The ratios between different groups are, up to inversion, $ab/c,a,b/c$. On $V$ none lies in $\mu_{60}$. Since the fixed coefficients lie in $\mu_{60}$, entries from different groups cannot have ratio in $\mu_5$. A pentagon in $u$ must therefore lie in one group. The 2 groups of size 2 are too small, leaving only the first 7 columns.

In those columns the root exponents of $u$, reduced modulo $12$, are
$$0,0,0,0,0,3,9.$$
A rotated $5$-cycle in $\mu_{60}$ has constant exponent modulo $12$, and its 5 values must be distinct. Hence the only possible pentagon support in $u$ is $\{0,1,2,3,4\}$; the displayed entries on it indeed form a pentagon.

For $v$, the 3 parameter factors are $1,b,c$. The corresponding cross-group ratios are $b,c,b/c$, also excluded from $\mu_{60}$ on $V$. Its first 7 root exponents modulo $12$ are
$$0,0,0,0,3,0,9.$$
Its only pentagon support is consequently $\{0,1,2,3,5\}$. These 2 supports differ. Already the quotients of $u$ and $v$ with $1$ cannot share a pentagon, so the 5-row matrix has no common pentagon. The invariant property of $H(t)$ proves the inequivalence assertion.
\end{proof}

\section{Extensions of DDD triples}\label{sec:extensionsDDD}\label{sec:continuousDDD}
We now consider extensions throughout the $\mathrm{DDD}$ support census. The zero-dimensional presentations give a finite classification; the subsequent subsections treat selected continuous families, with the domains of the extension results stated explicitly.

\subsection{Finite triples}\label{sec:finiteDDD}\label{sec:DDDfourth}

We next determine all regular extensions of the triples arising from the zero-dimensional chosen $\mathrm{DDD}$ presentations.

The feasibility census stored one quintic witness per support. Extensions require every label assignment, so first expand all labels on representatives of the $452$ feasible support classes.

Of the $452$ feasible support classes, $218$ have dimension zero, by Theorem~\ref{n11_noA_DDD:main}. Let $Z_{\mathrm D}$ denote the set of regular $3\times11$ matrices admitting a chosen $\mathrm{DDD}$ presentation in one of these $218$ classes, taken up to ordinary equivalence. 

\par\addvspace{\medskipamount}\noindent
\begin{minipage}{\linewidth}
\begin{theorem}\label{n11_fourth:main}
The set $Z_{\mathrm D}$ consists of $54836$ ordinary equivalence classes. Every representative has entries in $\mu_{40}$ after equivalence.

No member of $Z_{\mathrm D}$ can occur as 3 rows of a regular $6\times11$ partial Hadamard matrix. More precisely, the maximum number of rows in a regular extension is distributed as follows:
$$\begin{array}{c|r}
\text{maximum number of rows}&\text{classes in }Z_{\mathrm D}\\\hline
3&4000\\
4&44292\\
5&6544\\\hline
\text{total}&54836.
\end{array}$$
Thus none of these triples can occur in a regular square matrix of order 11.
\end{theorem}
\end{minipage}
\par\addvspace{\medskipamount}

We first expand all root labels and identify the resulting triples up to equivalence. We then enumerate their fourth-row families and compare those families. The lemmas below justify these steps; the proof concludes with the compatibility graphs and their extension witnesses.

\medskip
Fix a representative $(P,Q,R)$ of one feasible support class. The chosen row quotients are $u$, $v$ and $u/v$, with first row $1$. The pentagon in each partition has a distinguished anchor, its least column. Normalize the pentagon phases of $u$ and $v$ by row multiplication. Assign labels $0,1,2,3,4$ to each pentagon, with label zero at its anchor. There are $4!=24$ choices for each of the 3 pentagons, hence
$$24^3=13824$$
label triples to examine on each support representative. The label origin on the pentagon of $u/v$ is absorbed in its own cycle phase. Only differences of its labels enter the phase equations.

In particular, the first nonzero label on the pentagon of $u$ is not fixed to one in this expansion. That normalization was useful for the earlier existence test, but field scaling is not one of our ordinary Hadamard equivalences. Here all 4 possible nonzero labels are included.

Extend the 3 pentagon label functions to 11 columns by setting them equal to zero outside their pentagons. Write them as $a_j,b_j,c_j$. On every block $S$ of $R$, with anchor $j_S$, define
\[
f_{S,j}=(c_j-a_j+b_j)-(c_{j_S}-a_{j_S}+b_{j_S}),
\qquad j\in S\setminus\{j_S\}.
\]
The pair phases of $u,v$ remain arbitrary torus variables; assigning zero labels outside the pentagons does not constrain them. The 7 equations are precisely
$$Ax=f/5+g/2\pmod1,$$
with $A,g$ and $UAV=\operatorname{diag}(d_0,\ldots,d_{r-1},0)$ as in Subsection~\ref{sec:DDD}.

Every nonzero $d_i$ is invertible modulo 5. Therefore the quintic labels are feasible exactly when
\[
(Uf)_i=0\pmod5\qquad(i\geq r).
\]
For a feasible assignment choose $\widehat y_i\in\{0,1,2,3,4\}$ satisfying
$$d_i\widehat y_i=(Uf)_i\pmod5\qquad(i<r),$$
and put $\widehat y_i=0$ for $i\geq r$. All full phase solutions are then
\[
x=V y,\qquad
y_i=\frac{\widehat y_i}{5}+\frac{(Ug)_i+2k_i}{2d_i}\quad(i<r),
\qquad y_i=t_i\quad(i\geq r),
\]
where $0\leq k_i<d_i$ and $t_i\in\mathbb R/\mathbb Z$. The fixed quintic contributions on the pentagons are supplied by $a,b$, and the antipodal signs by the original pair offsets. The displayed Smith-coordinate parametrization therefore reconstructs the entire affine family for each allowed label triple.

The complete expansion has the following counts:
\[
\begin{array}{c|r|r|r}
\text{dimension}&\text{support classes}&\text{feasible label triples}&\text{affine components}\\\hline
0&218&576960&1227520\\
1&190&94768&194336\\
2&41&3168&6480\\
3&3&216&432\\\hline
\text{total}&452&675112&1428768.
\end{array}
\]
The initial feasibility scan has $452\cdot13824=6248448$ decisions. The last column counts the binary choices in the displayed Smith-coordinate parametrization separately. It is a count of labelled affine presentations, with repetitions allowed between support representatives and with no assertion that intersecting families are disjoint.

\medskip
For dimension zero, the displayed Smith-coordinate parametrization gives a fixed triple over $\mu_{40}$ for every finite component. To compare these triples under ordinary equivalence, write their exponent rows as
$$e_0=0,\qquad e_1,e_2\in(\mathbb Z/40\mathbb Z)^{11}.$$
Choose an ordering $(a,b,c)$ of the 3 rows and a base column $j_0$. Dephasing with row $a$ and column $j_0$ gives the 11 pairs
\[
\left(
 e_{b,j}-e_{a,j}-e_{b,j_0}+e_{a,j_0},\
 e_{c,j}-e_{a,j}-e_{c,j_0}+e_{a,j_0}
\right)\pmod {40}.
\]
Sort these pairs lexicographically. Take the least resulting word over the $6\cdot11=66$ row-order and base-column choices.

\par\addvspace{\medskipamount}\noindent
\begin{minipage}{\linewidth}
\begin{lemma}\label{n11_fourth:equivalence}
2 triples over $\mu_{40}$ have the same canonical word if and only if they are ordinarily equivalent as partial Hadamard matrices.
\end{lemma}
\end{minipage}
\par\addvspace{\medskipamount}

\begin{proof}
Each operation used to form the word is an ordinary equivalence: row permutation, column permutation, row multiplication and column multiplication. Equality of words therefore gives an equivalence.

Conversely, once a first row and a base column are chosen, dephasing removes all row and column multipliers. For a fixed row ordering, column equivalence is then exactly equality of the multiset of the 11 ordered pairs. Sorting computes that multiset. An arbitrary ordinary equivalence chooses some row ordering and some base column, both among the $66$ possibilities. Hence equivalent triples give the same least word. Neither complex conjugation nor a Galois action is added to this procedure.
\end{proof}

Applying this test to the $1227520$ fixed labelled components gives $54836$ distinct words. The archive records which chosen support classes produce each word. This is a classification of the finite set $Z_{\mathrm D}$ under ordinary equivalence, rather than a classification of all $\mathrm{DDD}$ triples. The possible overlap with positive-dimensional presentations does not change equality or inequality between 2 members of this finite set.

For the extension computation it is convenient to choose another representative of the same triple: minimize the entire first nonconstant exponent row before minimizing the second one. This groups the $54836$ triples into $36$ lists with a common row pair $(1,u)$. The grouping is only a way to reuse the 2-row calculation. Every distinct third row in each list is subsequently tested.

\medskip
For each of the $36$ fixed row pairs over $\mu_{40}$, take $M=840$, put $w=e^{\pi i/420}$, and enumerate all $1480490$ first presentations $x_j=z_{b(j)}w^{a_j}$. A quotient of type $\mathrm{A}$ is impossible, since the old pairs admit type $\mathrm{D}$. If $u_j=w^{h_j}$, a cycle anchored at $k$ imposes
\[
\frac{z_{b(j)}}{z_{b(k)}}=w^{(M/p)\ell+a_k-a_j+h_j-h_k}.
\]
Enumerate every cycle cover and solve these constant-ratio equations by forest integration. Closing edges test consistency. Repeat against each third row in the group.

\par\addvspace{\medskipamount}\noindent
\begin{minipage}{\linewidth}
\begin{lemma}\label{n11_fourth:rowfamilies}
For one of the fixed row pairs $(1,u)$ above, whose quotient has profile $\mathrm{D}$ and entries in $\mu_{40}$, the complete set of rows regular against both old rows is a finite union of families
\[
x_j=c_j z_{g(j)},\qquad c_j\in\mu_{840},
\]
where each group has an independent unit phase. After removing the common row phase, a family with $c$ groups has $c-1$ parameters. Imposing regularity against the third row of the given $\mathrm{DDD}$ triple gives another finite union of families of this same form.
\end{lemma}
\end{minipage}
\par\addvspace{\medskipamount}

\begin{proof}
The old $\mathrm{D}$ quotient excludes profile $\mathrm{A}$ for either new quotient by Theorem~\ref{n11_A_branch:profiles}. All remaining cycle lengths divide $840$. The first assertion now follows by tree integration of the displayed constant-ratio equations, with all remaining edges checked for consistency. A root for each connected component gives its independent phase. There is no root-order assumption on these free phases.

For the second assertion use the component phases as the new graph vertices and enumerate a complete cycle cover against the third old row. Its edges again have constants in $\mu_{840}$ and have the same form as the displayed constant-ratio equations. Further connected components are joined or their consistency is checked. The remaining component phases are still independent, and the coefficients still belong to $\mu_{840}$.
\end{proof}

Normalize each phase group at its first column and number groups in order of appearance. Identical coefficient and group vectors are merged. Intersecting families remain in the union. The $36$ pair scans are reused for every associated third row; every one of the $54836$ triples is tested.

The first-presentation scan makes
$$36\cdot1480490=53297640$$
presentation decisions in total. Its output for each common row pair is reused for every associated third row. The subsequent calculation is exhaustive for each of the $54836$ triples; grouping the first 2 rows does not identify or omit third-row tests.

\medskip

For fixed rows, a nonzero reduction of the quotient sum excludes regularity. For a $2$-group family with quotient sum $A+zB$, the necessary condition is $A\bar A=B\bar B$. The primary and audit use splitting primes $2521$ and $3361$; every survivor is subjected to a complete cycle-cover test. Families with more groups pass directly to that test.

The resulting fourth-row lists have $479580$ point records and $9192$ circle records; $1560$ old triples have a circle. No family with $2$ or more free row phases survives. The union is complete, although records may overlap.

\medskip
The odd order provides a useful restriction which does not depend on the particular coefficients of a circle.

\par\addvspace{\medskipamount}\noindent
\begin{minipage}{\linewidth}
\begin{lemma}\label{n11_fourth:circle}
A family $x_j=c_jz_{g(j)}$ with 2 nonempty phase groups at order 11 contains no pair of orthogonal rows, even when the rows have different parameter values.
\end{lemma}
\end{minipage}
\par\addvspace{\medskipamount}

\begin{proof}
Remove a common row phase and write $x(z)=A+zB$, where $A$ and $B$ have disjoint supports $S_0,S_1$, and every nonzero entry has modulus one. Then
$$\langle x(z),x(t)\rangle=|S_0|+z\bar t\,|S_1|.$$
If this vanished, absolute values would give $|S_0|=|S_1|$. Their sum is 11, so equality is impossible.
\end{proof}

\begin{proof}[Proof of Theorem~\ref{n11_fourth:main}]
The complete label expansion and Lemma~\ref{n11_fourth:equivalence} give the $54836$ classes and their representatives over $\mu_{40}$. Lemma~\ref{n11_fourth:rowfamilies} and the exhaustive scans above give every possible fourth row. By Lemma~\ref{n11_fourth:circle}, an orthogonal collection of added rows can use each point or circle record at most once. This statement remains valid when records overlap: 2 orthogonal rows cannot belong to the same circle, so choosing a containing record for each added row necessarily chooses distinct records.

\medskip
For a fixed old triple, form a graph whose vertices are its retained point and circle records. Join 2 vertices when some row from the first and some row from the second have a regular quotient. The parameters in this definition are arbitrary unit numbers. Every regular extension of the old triple selects a clique in this graph, by Lemma~\ref{n11_fourth:circle}.

To decide an edge, write the 2 families as
$$x_j=c_j q^{g_j},\qquad y_j=d_j s^{h_j},
\qquad g_j,h_j\in\{0,1\}.$$
For a point record use the constant group index zero. Their quotient is
\[
\frac{x_j}{y_j}=c_j\bar d_j\,q^{g_j}s^{-h_j}.
\]
There are at most 4 distinct pairs $(g_j,h_j)$, called cells in this calculation. Inside one cell all ratios of entries belong to $\mu_{840}$.

No $11$-cycle is possible in $x_j/y_j=c_j\bar d_jq^{g_j}s^{-h_j}$. Such a cycle could use at most one column from each cell, because the ratio of 2 different entries in an $11$-cycle has order 11, whereas a within-cell ratio has order dividing $840$. 4 cells cannot supply 11 columns. It is therefore enough to enumerate cycles of lengths $2,3,5,7$.

First temporarily attach an independent phase $U_{ab}$ to each cell present. For every chosen cycle, constant-ratio graph equations determine relations between these cell phases. There is one constraint to restore if all 4 cells occur:
\[
U_{00}U_{11}=U_{01}U_{10}.
\]
Indeed, up to a common phase the 4 required values are $1,s^{-1},q,qs^{-1}$, which satisfy this identity. Conversely, any 4 unit values satisfying $U_{00}U_{11}=U_{01}U_{10}$ have this form after division by $U_{00}$.

If fewer than 4 cells occur, arbitrary phases on the present cells can be represented, up to a common phase, by some $q,s\in\mathbb T$. For 3 cells this follows by solving 2 independent equations whose exponent determinant is $\pm1$. For 2 cells their exponent difference is primitive, so the corresponding character maps onto the circle. The one-cell case has no phase condition.

After a complete cycle cover has been chosen, let $t_k$ be the independent component phases and let $k(a,b)$ index the component containing cell $(a,b)$. Write the solved cell phases as
$$U_{ab}=t_{k(a,b)}w^{e_{ab}}.$$
Substitution into $U_{00}U_{11}=U_{01}U_{10}$ gives
$$\prod_k t_k^{m_k}=w^{e}.$$
If some $m_k$ is nonzero, this equation has a solution on the torus, because a nonzero character maps onto $\mathbb T$. If all $m_k$ vanish, the required root must be one. This decides exactly whether the complete cycle cover is feasible for the original 2 parameters. Trying all cycle covers therefore decides the graph edge.

For an independent calculation, keep $q,s$ from the beginning. Every cycle edge gives an equation
\[
q^a s^b=w^{v},\qquad a,b\in\{-1,0,1\}.
\]
The first nontrivial equation has primitive coefficient vector. A parallel equation is a direct root-consistency check. 2 independent equations have determinant $\pm1$ or $\pm2$, so they have an explicit finite solution list in $\mu_{1680}^2$. Every subsequent equation is checked on that list. If the rank stays below 2, the remaining variable is free. This gives a second exact decision procedure which does not use cell phases or the rectangle identity.

The second solver saves an actual compatible pair whenever an edge exists: take free phases equal to $1$ in ranks $0,1$, and a retained root solution in rank $2$. The resulting rows lie in $\mu_{1680}$ and are checked against the old triple and each other.

Every compatibility graph in this calculation is triangle-free. This is checked independently from the saved edge lists: for each edge, its endpoints have no common neighbor. In particular, it is not a numerical estimate of clique size.

If a regular extension had 6 rows, the 3 rows added to the original triple would require 3 distinct vertices by Lemma~\ref{n11_fourth:circle}. Their 3 quotients would be regular, so the vertices would form a triangle. The checked triangle exclusion proves that 6 rows are impossible.

The same graphs determine the exact maximum number of rows in a regular extension. An empty vertex list means that no fourth row exists. A nonempty graph with no edges supplies a fourth row but no fifth. A graph with an edge supplies 2 compatible rows and hence a regular 5-row extension; the triangle exclusion prevents a sixth. The saved point or circle records give the fourth-row witnesses, and the saved parameter values from $q^as^b=w^v$ give the fifth-row witnesses.

The totals are therefore exactly the 3 rows of the table in Theorem~\ref{n11_fourth:main}. This completes its proof. No compatibility between 3 independently chosen parameter pairs is assumed: the triangle-free upper bound needs only necessary pairwise compatibility, while existence of a 5-row extension needs just one pair, for which the saved parameters are simultaneous witnesses.
\end{proof}

\medskip
The one support class with Smith invariant 4 is included in the theorem. For the chosen representative its pentagons and pairs are
$$\begin{array}{c|c|c|c|c}
&\text{pentagon}&\text{pair}&\text{pair}&\text{pair}\\\hline
P&\{0,1,2,3,4\}&\{5,6\}&\{7,8\}&\{9,10\}\\
Q&\{0,1,5,7,9\}&\{2,6\}&\{3,8\}&\{4,10\}\\
R&\{2,3,5,7,10\}&\{0,6\}&\{1,8\}&\{4,9\}.
\end{array}$$
Its Smith diagonal is $(1,1,1,1,1,4)$. Exactly $2816$ normalized pentagon-label triples are feasible; their 4 binary components give $11264$ fixed labelled matrices. Ordinary equivalence reduces them to $632$ triples. The maximum numbers of rows in their regular extensions are
$$\begin{array}{c|r}
\text{maximum number of rows}&\text{exceptional-class triples}\\\hline
3&568\\4&56\\5&8\\\hline
\text{total}&632.
\end{array}$$
Thus the exceptional class does have fourth-row and fifth-row survivors. Its exclusion from square matrices requires the compatibility argument, not merely the absence of a fourth row.

Here is one fifth-row survivor, with rows numbered from zero. Put $w=e^{\pi i/60}$ and raise it entrywise to
\[
\begin{pmatrix}
0&0&0&0&0&0&0&0&0&0&0\\
0&0&6&15&30&54&60&60&75&78&102\\
0&0&60&45&105&114&18&66&90&60&42\\
0&30&71&105&90&65&43&11&60&103&25\\
0&90&1&75&60&115&53&61&30&113&35
\end{pmatrix}.
\]
The first 3 rows belong to the exceptional $\mathrm{DDD}$ support class. Their displayed coefficients lie in $\mu_{40}$. The new quotients admit the following profiles:
$$\begin{array}{c|cccccccccc}
\text{row pair}&01&02&12&03&04&13&14&23&24&34\\\hline
\text{chosen profile}&D&D&D&F&F&F&F&F&F&E.
\end{array}$$
Complete cycle supports, in the same column numbering, are
$$\begin{array}{c|l}
01&(0,6),\ (1,7),\ (2,10,9,5,4),\ (3,8)\\
02&(0,2),\ (1,9),\ (3,4),\ (5,8,7,10,6)\\
12&(0,5),\ (1,10),\ (2,3,7,9,6),\ (4,8)\\
03&(0,8),\ (1,4),\ (2,7),\ (3,5,10),\ (6,9)\\
04&(0,4),\ (1,8),\ (2,7),\ (3,10,5),\ (6,9)\\
13&(0,4),\ (1,3),\ (2,9,8),\ (5,7),\ (6,10)\\
14&(0,3),\ (1,4),\ (2,8,9),\ (5,7),\ (6,10)\\
23&(0,3),\ (1,8),\ (2,5),\ (4,7,6),\ (9,10)\\
24&(0,8),\ (1,3),\ (2,5),\ (4,6,7),\ (9,10)\\
34&(0,1),\ (2,6,3),\ (4,5,9),\ (7,10,8).
\end{array}$$
For each quotient, the entries along a listed block differ successively by a primitive root of that block's length. These 10 certificates prove regularity and orthogonality exactly.

The cross ratio from rows zero and 3 and columns zero and 2 is $w^{71}$ or its inverse. Since $\gcd(71,120)=1$, it has order $120$. Cross-ratio multisets are invariant under ordinary equivalence, up to permutation. Hence any equivalent partial Butson representative over $\mu_m$ must have $120\mid m$. The displayed matrix already has level $120$, so its minimal Butson level is exactly $120$.

Theorem~\ref{n11_fourth:main} excludes any regular sixth row for this example, because its first 3 rows lie in $Z_{\mathrm D}$. This illustrates why new quotient types must be allowed when extending a $\mathrm{DDD}$ triple, even though the old 3 quotients remain of type $\mathrm{D}$.

\medskip

The quintic-label expansion is independently checked by Gaussian elimination over $\mathbb F_5$. Both expansions agree on all $6248448$ decisions. The $36$ pair enumerations and third-row refinements are repeated with a different phase graph and splitting prime; the retained lists agree. The 2 compatibility solvers give the same graphs. The final verifier checks the Smith identities, all symbolic row families, fifth-row witnesses and triangle exclusions.

The preceding theorem and its proof establish Theorem~\ref{intro:finite}. They concern all regular extensions of the listed triples, with arbitrary unit phases in the added rows.

\medskip
The finite extension theorem does not settle the positive-dimensional $\mathrm{DDD}$ presentations. The following subsections retain the results obtained for whole support classes, a distinguished three-parameter triple and several of its exceptional loci. The hypotheses distinguish a uniform statement on a full torus from a generic statement or a calculation with bounded root level.

\subsection{Two one-parameter support classes}\label{sec:DDDOne}
The positive-dimensional census gives $201248$ labelled affine presentations in $190+41+3$ support classes. We now treat 2 whole one-parameter classes. For the first, number the columns from zero and take the following supports for $u,v,u/v$:
$$\begin{aligned}
P&=((0,1,2,3,4),(5,6),(7,8),(9,10)),\\
Q&=((0,1,2,3,5),(4,7),(6,9),(8,10)),\\
R&=((0,1,2,4,10),(3,7),(5,8),(6,9)).
\end{aligned}$$
Both $v$ and $u/v$ have an antipodal pair in columns $6,9$, so $u_6=u_9$. The free phase rotates these 2 entries of $v$ together. This is class $3203$ in the census.

\par\addvspace{\medskipamount}\noindent
\begin{minipage}{\linewidth}
\begin{theorem}\label{ddd_one:bound}
Every regular extension of a triple admitting this chosen support, up to support equivalence, has at most 6 rows. In particular the class admits no regular square completion.
\end{theorem}
\end{minipage}
\par\addvspace{\medskipamount}

\begin{proof}
The matrix in the phase system $Ax=f/5+g/2$ has Smith diagonal $(1,1,1,1,2)$. Of the $(4!)^3$ normalized pentagon labellings, 480 satisfy the zero-row equations. Each gives 2 binary components and one free phase. Applying the Smith-coordinate parametrization of Subsection~\ref{sec:finiteDDD} to all 960 presentations gives the following normalization. With $w=e^{\pi i/10}$, the first row is $1$ and the second row is $u=(w^{b_j})$ with
$$b=(0,0,1,5,5,9,10,10,13,15,17).$$
The third row has entries $w^{c_j}z^{d_j}$, where $z\in\mathbb T$ and exactly 2 of the $d_j$ are 1, the others being zero. To obtain the list, rotate the fixed quotient to $u$, allowing reversal by a row permutation. Use a row phase and a rotation of $z$ to normalize one coefficient on each of the 2 parts of $v$. Sort columns within equal entries of $u$ and take the least resulting list. This gives 32 descriptions. For instance, one has
$$c=(0,5,1,9,15,17,0,10,5,10,13),\qquad d_6=d_7=1.$$

The fixed pair $(1,u)$ is group 21 of the exhaustive pair calculation in Subsection~\ref{sec:finiteDDD}. It admits $339008$ normalized added-row records with at most 3 phase groups. By Theorem~\ref{n11_A_branch:profiles}, a new quotient of type $\mathrm{A}$ against any old row is excluded by a triple containing a chosen $\mathrm{D}$ quotient. For each pair record $x$, enumerate all covers of $x/v$ by cycles of lengths $2,3,5,7$. Within each proposed cycle the ratios impose equations on the phase groups of $x$. Each edge prescribes a ratio in $\mu_{840}$ times $z^d$ with $d=0,1,-1$. Integrate the edges along a forest. A closing edge gives $z^k=c$, with $c\in\mu_{840}$ and $|k|\le3$, since its cycle has length at most 3. For $k=0$ test the constant equation. For $k\ne0$ retain every solution, all lying in $\mu_{5040}$. Thus the enumeration covers the full circle and retains every free group phase.

Over the 32 descriptions, the 2 solvers agree on $4112$ records valid for free $z$ and $48640$ records with specified $z$. The specified values all lie in $\mu_{60}$. The verifier checks literal prime-cycle covers against the old rows for every record. Each remaining record has at most 2 phase groups. At any given $z$, two rows from one record cannot be orthogonal: for 2 groups their scalar product is a sum of 2 terms whose moduli are the group sizes, unequal since their sum is 11.

Make one graph of the free records for each description. At each specified value of $z$ make a graph of the specialized free records together with all records at that value. This gives 32 free graphs and 1376 specialized graphs. To bound simultaneous extensions, join 2 records whenever the following necessary test permits orthogonality. In their quotient, collect columns with the same powers of the free phases, including $z$ if it is free. A single coefficient sum must vanish. For 2 sums their squared moduli must agree. With more sums retain the edge. The tests use roots in $\mu_{840}$ and exclude an edge only by a nonzero reduction at a splitting prime.

Any simultaneous additions give a clique in the appropriate graph. The calculation modulo 2521 bounds every clique by 3. A separate calculation modulo 3361 checks directly that no edge has 2 adjacent common neighbours, excluding a 4-clique in every graph. The component \texttt{ddd\_one} contains both calculations and the labelled records. At most 3 rows can therefore be added to the original triple.
\end{proof}

We next consider class $3217$. Its first support $P$ is as above, and its other supports are
$$\begin{aligned}
Q&=((0,1,2,3,5),(4,7),(6,8),(9,10)),\\
R&=((0,1,2,4,5),(3,8),(6,9),(7,10)).
\end{aligned}$$
Here $P$ and $Q$ share the pair $(9,10)$. The free phase rotates that pair in both rows, so their quotient is fixed. After a row permutation it supplies the same fixed pair $(1,u)$ as in Theorem~\ref{ddd_one:bound}. Normalize the third row as $v_j=w^{c_j}z^{d_j}$, with $w=e^{\pi i/10}$, $d_6=d_7=1$ and the other $d_j=0$. The 16 resulting descriptions all have $(c_0,c_1,c_6,c_7,c_9)=(0,5,0,10,10)$. Four of them are specified by the following table:
$$\begin{array}{c|rrrrrr}
r&c_2&c_3&c_4&c_5&c_8&c_{10}\\\hline
1&11&3&19&15&15&7\\
2&15&3&11&19&15&7\\
3&11&7&19&15&3&15\\
4&15&7&11&19&3&15
\end{array}$$

\par\addvspace{\medskipamount}\noindent
\begin{minipage}{\linewidth}
\begin{theorem}\label{ddd_one:3217}
In class $3217$, each of the 12 descriptions outside the table has maximum regular extension exactly 5 for every $z\in\mathbb T$. For the description indexed by $r$ in the table, the maximum is 4 precisely at the 40 values
$$z=e^{\pi i m/30},\quad m\in\mathbb Z/60\mathbb Z,\quad
5\nmid m,\quad (m\not\equiv r\pmod5\ \text{or}\ 3\mid m).$$
At all other phases of that description the maximum is 3. In particular every regular extension in the class has at most 5 rows.
\end{theorem}
\end{minipage}
\par\addvspace{\medskipamount}

\begin{proof}
The Smith diagonal is $(1,1,1,1,2)$. The label equations give 480 pentagon labellings, each with 2 binary components. Normalizing all 960 presentations as in Theorem~\ref{ddd_one:bound} gives the 16 descriptions above. Against each third row we test all $339008$ records of the fixed pair. The two cycle solvers agree on $2976$ records valid for free $z$ and $32688$ records at specified phases in $\mu_{60}$. Each record has at most 2 phase groups. As before, at a fixed $z$ it contributes at most one row to a simultaneous extension.

In each of the 4 tabulated descriptions every record specifies $z$. The literal covers give a fourth row at exactly the displayed values. At those values the compatibility graphs have vertices and no edges. This proves the two maxima for these descriptions.

In each of the other 12 descriptions there are 248 records valid for free $z$. Six are circles, each carrying one extra phase on the same block of columns. Of the 242 individual records, 240 are specializations of those circles: on each of the 2 blocks the ratio to the circle coefficients is a single Laurent monomial. This leaves 6 circles and 2 individual rows to consider at arbitrary $z$.

For a pair involving a circle, the scalar product has the form $A(z)+sB(z)$, with $s\in\mathbb T$ and $A,B$ Laurent polynomials over $\mathbb Q(w)$. Orthogonality requires $A\bar A-B\bar B=0$, with conjugation sending $z$ to $z^{-1}$. For two individual rows use their scalar product directly. Remove monomial factors from each resulting polynomial. Three records can contribute mutually orthogonal rows only if their 3 polynomials have a common root. Among the $12\binom83=672$ triples, 624 contain a nonzero constant polynomial. In each of the remaining 48, two quadratics have nonzero resultant. The certificate checks these resultants and all common divisors in $\mathbb Q(w)$ using $w^8-w^6+w^4-w^2+1=0$. Thus at most 2 free records can contribute simultaneously. At every specified phase, the graphs include both the specialized free records and all additional records. Every such graph in both splitting-prime calculations has clique number at most 2.

For the matching lower bound, set the extra circle phases to 1. In each of the 12 descriptions, two of the resulting rows have a mutual prime-cycle cover identically in $z$. For example, take $c=(0,5,3,11,15,19,0,10,7,10,15)$. Use the rows given by $w$ raised to
$$\begin{aligned}
X&=(0,15,0,12,10,4,5,15,16,5,8),\\
Y&=(0,10,0,16,5,12,15,0,8,10,4),
\end{aligned}$$
multiplying columns $6,7$ of the first and $4,6,7,9$ of the second by $z$. The verifier checks all 10 quotients of each displayed or saved 5-row construction by literal cycles with $z$ free. This gives 5 rows throughout the 12 descriptions and proves the result.
\end{proof}

\subsection{All generic fourth rows of the 3-parameter DDD triple}\label{sec:DDDgenerictriple}

Write $R=(1;u;v)$ for the first $3$ rows of $G$. We determine every fourth row on an explicit open dense parameter set, allowing all subsequent row choices.

Put $w=e^{\pi i/10}$ and write
\[
R(a,b,c)=\begin{pmatrix}
1&1&1&1&1&1&1&1&1&1&1\\
1&w^{4}&w^{8}&w^{12}&w^{16}&w&w^{11}&-ab/c&ab/c&a&-a\\
1&w^{8}&w^{16}&w^{4}&w^{17}&w^{12}&w^{7}&b&-b&c&-c
\end{pmatrix}.
\]
Thus $R=(1;u;v)$, with the parameters and column order of Subsection~\ref{sec:DDDpositivefive}. In particular it is regular, of chosen type $\mathrm{DDD}$, for every $(a,b,c)\in\mathbb T^3$.

\medskip
The exceptional conditions will be stated using 2 finite lists of characters. Define
\[
\begin{split}
M_{420}={}&\left\{c,b,\frac ab,\frac{ac}b,\frac ac,
 \frac{ab}{c^2},ab,\frac{ab^2}{c^2},\frac{a^2b}{c^2},
 \frac{a^3b}c,\frac{a^3b^2}{c^2}\right\},\\
M_{840}={}&\left\{\frac bc,a,\frac{ab}c,\frac{a^2b}c\right\}.
\end{split}
\]
Each member of either list denotes the indicated function of $(a,b,c)$. Let
\[
U=\left\{(a,b,c)\in\mathbb T^3:
 \begin{array}{ll}
 m(a,b,c)\notin\mu_{420}&\text{for every }m\in M_{420},\\
 m(a,b,c)\notin\mu_{840}&\text{for every }m\in M_{840}
 \end{array}\right\}.
\]
This is an explicit nonresonance condition. Its complement is a finite union of proper translated subtori. Consequently $U$ is open and dense, and its complement has Haar measure zero. In particular, the condition is not an assumption that the parameters are roots of unity, nor an appeal to a numerical meaning of genericity.

Normalize each new row by making its first entry equal to one. For the following row display put $w=e^{\pi i/60}$ and define 8 fixed rows by $x_{ij}=w^{L_{ij}}$, where
\[
L=\begin{pmatrix}
0&72&24&96&48&48&48&8&8&88&88\\
0&72&24&96&48&48&48&88&88&8&8\\
0&72&24&96&78&78&108&18&18&48&48\\
0&72&24&96&78&78&108&48&48&18&18\\
0&96&72&48&24&24&24&64&64&104&104\\
0&96&72&48&24&24&24&104&104&64&64\\
0&96&72&48&54&54&84&24&24&114&114\\
0&96&72&48&54&54&84&114&114&24&24
\end{pmatrix}.
\]
The row indices in this display are $0,\ldots,7$. They label possible additional rows; they are not the row indices of $R$.

\par\addvspace{\medskipamount}\noindent
\begin{minipage}{\linewidth}
\begin{theorem}\label{n11_generic:main}
For every $(a,b,c)\in U$, the normalized rows which extend $R(a,b,c)$ to a regular 4-row partial Hadamard matrix are exactly $x_0,\ldots,x_7$ in the displayed exponent matrix $L$. None of these rows has an additional free parameter, and all 8 are independent of $a,b,c$.

The only unordered pairs of these rows which are orthogonal are
\[
\{x_0,x_4\},\qquad \{x_1,x_5\}.
\]
Both quotients in $\{x_0,x_4\}$ and $\{x_1,x_5\}$ are regular. Thus the largest number of rows in a regular extension of $R(a,b,c)$ is exactly 5. Up to the order and phases of the added rows, its 5-row extensions are exactly
\[
(R;x_0;x_4),\qquad (R;x_1;x_5).
\]
In particular, $R(a,b,c)$ has no regular 6-row extension and cannot occur in a regular square matrix of order 11.
\end{theorem}
\end{minipage}
\par\addvspace{\medskipamount}

\begin{proof}
For the enumeration put $w=e^{\pi i/420}$. Profile $\mathrm{A}$ is excluded by the $\mathrm{A}$-branch theorem, since the old quotients admit $\mathrm{D}$ decompositions. For each of the $1480490$ presentations write $x_j=z_{b(j)}w^{e_j}$ and impose all cycle covers against $u,v$. If $h_j=w^{d_j}a^{p_j}b^{q_j}c^{r_j}$, an edge of a $p$-cycle gives
\[
\frac{z_{b(j)}}{z_{b(q)}}=w^{840\ell/p+e_q-e_j}\frac{h_j}{h_q}.
\]
Forest integration solves the component phases. A closing edge gives
\[
a^r b^s c^t=w^{e}.
\]
The generic scan accepts identities, rejects inconsistent constants, and records every nonzero character equation before rejecting that branch. The graph has at most $5$ vertices, and each old parameter has exponent range at most $1$ in either old row. Fundamental cycles therefore give $|r|,|s|,|t|\le5$. This also holds for the original graph combining the equations from both old rows.

\medskip
For each first presentation, the recursion chooses an uncovered column as anchor, tries each permitted prime, assigns its remaining labels bijectively to unused columns, and checks each phase-graph edge relation as soon as it is imposed. Every complete cycle cover occurs. Repeating this first against $u$ and then against $v$ therefore exhausts all formal generic fourth-row presentations.

After a complete cover against $u$, normalize the resulting family by numbering its connected components on first occurrence and dividing each component by the monomial at its first column. Record the component of each column, its 3 integer exponents and its root residue. Identical records define exactly the same set of rows for all old parameter values and are retained once. The subsequent calculation against $v$ depends only on this record. No ordinary equivalence between distinct rows is used to discard a candidate.

The complete counts are
\[
\begin{array}{c|r|r|r}
\text{first profile}&\text{first presentations}&
 \text{records after }u&\text{records after }v\\\hline
\mathrm{B}&712800&0&0\\
\mathrm{C}&443520&276&4\\
\mathrm{D}&166320&1800&4\\
\mathrm{E}&123200&0&0\\
\mathrm{F}&34650&11264&0\\\hline
\text{total}&1480490&13340&8
\end{array}
\]

Every final record has a single connected phase component and zero exponents of $a,b,c$. Its common phase is removed by normalizing the first entry to one. The resulting rows are exactly $x_i=(w^{L_{ij}})_j$, $0\le i\le7$, with the exponent matrix $L$ displayed above. Rows $x_0,x_1,x_4,x_5$ come from profile $\mathrm{C}$ and belong to $\mu_{15}^{11}$. Rows $x_2,x_3,x_6,x_7$ come from profile $\mathrm{D}$ and belong to $\mu_{20}^{11}$. The primitive $120$th root used in the displayed exponent matrix $L$ permits the $8$ rows to be displayed in one integer matrix; it is not a claim that each row requires level $120$.

A separate symbolic checker verifies all $24$ quotients of these 8 rows against $1,u,v$. It groups entries with the same parameter character and checks the ordered prime-cycle root labels inside each group. Thus these are regularity identities valid for every $(a,b,c)\in\mathbb T^3$, not tests at selected parameter values.

\medskip
The alternative traversal described below records $5001$ distinct nonzero character equations $a^rb^sc^t=w^e$. A relation and its negative are identified by requiring the first nonzero character coefficient to be positive. Root residues are reduced modulo $840$. The complete list is retained in the companion file \texttt{exceptional\_equations.txt}.

Every actual fourth row has cycle covers followed by the generic recursion. Its first rejected equation cannot be an inconsistent constant, so either it reaches one of the $8$ final rows or its old parameters satisfy a recorded nonzero character equation. Merging identical intermediate records preserves their subsequent extension problems.

It remains to explain the simple character lists used in the theorem. All fixed old coefficients belong to $\mu_{20}$, and all new prime labels belong to $\mu_{420}$. The constants in every graph equation and every recorded relation therefore belong to $\mu_{420}$. In the modulus-$840$ representation this is checked by the evenness of every saved root residue.

For a recorded nonzero vector $(r,s,t)$, let
$$d=\gcd(|r|,|s|,|t|),\qquad
m=a^{r/d}b^{s/d}c^{t/d}.$$
Equation $a^rb^sc^t=w^e$ implies
\[
m\in\mu_{420d}.
\]
Group the $5001$ equations by their primitive vector and take the least common multiple of the resulting orders $420d$ within each group. There are exactly fifteen primitive vectors. 11 give order $420$, and 4 give order $840$. They are precisely the 2 lists in $M_{420},M_{840}$.

For completeness, the numbers of distinct recorded equations within the separate first-profile scans are
$$\begin{array}{c|rrrrr}
\text{first profile}&B&C&D&E&F\\\hline
\text{recorded equations}&3553&2454&2078&3044&4533.
\end{array}$$
These sets overlap, and their union has $5001$ elements. Their actual character coefficients have absolute value at most 4, consistent with the a priori bound $|r|,|s|,|t|\le5$. The exceptional list depends on the traversal: it records a sufficient obstruction to following a generic branch, not a canonical description of all special extensions.

\medskip
It remains to compare the 8 possible fourth rows with each other. Since all their entries lie in $\mu_{120}$, orthogonality is an exact polynomial calculation. For $i<k$, form
\[
P_{ik}(X)=\sum_{j=0}^{10}X^{(L_{ij}-L_{kj})\bmod120}.
\]
Then $x_i$ and $x_k$ are orthogonal if and only if the remainder of $P_{ik}$ modulo the cyclotomic polynomial $\Phi_{120}$ is zero. Integer polynomial division for all $\binom82=28$ pairs gives exactly the 2 pairs $\{x_0,x_4\}$ and $\{x_1,x_5\}$. A second check searches prime-cycle covers of the root differences. It finds the same 2 pairs, proving their regularity as well as their orthogonality.

The compatibility graph therefore consists of 2 disjoint edges and 4 isolated vertices. A regular 5-row extension must choose one of the 2 edges. A regular 6-row extension would require 3 distinct added rows which are pairwise adjacent. There is no such triple. Repetition of one normalized row cannot avoid this obstruction, because 2 proportional unit rows are not orthogonal. Each edge gives an actual regular 5-row extension, so the upper bound of 5 rows is attained.
\end{proof}

There is also a direct identification with the family in Subsection~\ref{sec:DDDpositivefive}. Its fixed rows $s,t$ are exactly $x_1,x_5$. Let $\pi$ exchange the 2 last column pairs, sending $(7,8,9,10)$ to $(9,10,7,8)$ and fixing the other columns. Put
\[
(a',b',c')=(-ab/c,c,b).
\]
Entrywise substitution gives
$$R(a,b,c)\pi=R(a',b',c'),\qquad
 x_0\pi=x_1,\qquad x_4\pi=x_5.$$
The parameter transformation $(a,b,c)\mapsto(-ab/c,c,b)$ is an involution of $\mathbb T^3$. Thus the first 5-row matrix $(R;x_0;x_4)$ is ordinarily equivalent to $G(a',b',c')$, and the second is $G(a,b,c)$ itself. Both specified 5-row constructions consequently have no regular sixth row at any parameter value, by Theorem~\ref{n11_positive:main}. The new point of Theorem~\ref{n11_generic:main} is the exhaustive statement: on $U$, every possible fifth-row choice has now been accounted for.

\medskip

The primary and audit scans exhaust the same first presentations, using different phase graphs and opposite cycle anchors. Their final lists agree. The independent verifier checks the triple, all $8$ rows, all $28$ pair comparisons, the parameter involution and the exceptional character cover.

Any sixth-row extension of this labelled triple must satisfy one of the $5001$ saved equations. These are necessary conditions only. Other labelled triples in support class $3286$, the other $2$ 3-dimensional support classes, and the one- and 2-dimensional families remain outside this result.

\subsection{2 complete exceptional surfaces of the DDD triple}\label{sec:DDDexceptionalsurfaces}

The generic calculation leaves a finite union of $2$-dimensional parameter sets. We now treat the $2$ surfaces $b=c$ and $b=-c$, including their special curves and points.

This choice is natural from the displayed triple. On $b=c$, the 2 parameter factors $ab/c$ and $a$ in $u$ coincide, as do the factors $b$ and $c$ in $v$. The last 4 columns become 3 rows of a real Hadamard matrix of order 4, with independent row multipliers. This coincidence produces additional fourth rows. We determine their full parameter loci and test all their possible compatibilities.

\medskip
Put $w=e^{\pi i/420}$. We first make the exceptional input precise. The $5001$ equations from Subsection~\ref{sec:DDDgenerictriple} have the form
$$a^r b^s c^t=w^{e},\qquad (r,s,t)\ne(0,0,0).$$
Put $d=\gcd(|r|,|s|,|t|)$ and $v=(r,s,t)/d$. The solutions are the union of the connected tori
\[
a^{v_1}b^{v_2}c^{v_3}=w^{f},\qquad
 df\equiv e\pmod{840}.
\]
For the saved equations there are exactly $d$ residues $f$ satisfying this congruence. Thus all components are represented in the same group $\mu_{840}$. This last assertion is checked on the actual input; it is not a general assertion about extracting roots without enlarging the root group.

The character $v$ is primitive, so its kernel in $\mathbb T^3$ is connected. 2 of the tori in $a^{v_1}b^{v_2}c^{v_3}=w^f$ coincide only if their primitive normals differ by sign and their constants change accordingly. We fix the first nonzero coordinate of $v$ to be positive. Deduplicating the resulting pairs $(v,f)$ therefore gives distinct connected tori, rather than distinct equations for possibly identical sets.

The exact result is
\[
\begin{array}{c|r@{\qquad}c|r}
\text{primitive character}&\text{tori}&\text{primitive character}&\text{tori}\\\hline
c&124&b/c&830\\
b&124&a/b&120\\
ac/b&414&a/c&124\\
a&840&ab/c^2&120\\
ab/c&780&ab&120\\
ab^2/c^2&380&a^2b/c^2&120\\
a^2b/c&480&a^3b/c&380\\
a^3b^2/c^2&60&&\\\hline
\multicolumn{3}{r|}{\text{total distinct connected tori}}&5016
\end{array}
\]
Different tori in this table can intersect. Such intersections are not removed from either torus. The table is an exact decomposition of the saved sufficient cover; it does not assert that every listed torus actually supports a new fourth row.

\medskip
Use the matrix $R(a,b,c)$ of Subsection~\ref{sec:DDDgenerictriple}, and put $R_+(a,b)=R(a,b,b)$. With $\omega=w^{42}=e^{\pi i/10}$, it has the block form
\[
R_+(a,b)=\left(\begin{array}{c|c} K&Q(a,b)\end{array}\right),
\]
where
\begin{equation*}
K=\begin{pmatrix}
1&1&1&1&1&1&1\\
1&\omega^{4}&\omega^{8}&\omega^{12}&\omega^{16}&\omega&\omega^{11}\\
1&\omega^{8}&\omega^{16}&\omega^{4}&\omega^{17}&\omega^{12}&\omega^{7}
\end{pmatrix},
\quad
Q(a,b)=\begin{pmatrix}
1&1&1&1\\
-a&a&a&-a\\
b&-b&b&-b
\end{pmatrix}.
\end{equation*}
The rows are orthogonal separately in both blocks. This observation explains why the surface is special, but does not by itself obstruct an arbitrary additional row: a new row need not split into 2 orthogonal blocks.

\par\addvspace{\medskipamount}\noindent
\begin{minipage}{\linewidth}
\begin{theorem}\label{n11_surface:main}
For every $a,b\in\mathbb T$, the largest number of rows in a regular extension of $R(a,b,b)$ is exactly 5. The same assertion holds for $R(a,b,-b)$. Neither triple has a regular 6-row extension, and neither can occur in a regular square matrix of order 11.
\end{theorem}
\end{minipage}
\par\addvspace{\medskipamount}

\begin{proof}
The lower bound is supplied by the already displayed family $G$: its first 3 rows are $R$, and its last 2 rows provide a regular 5-row extension for every parameter value. The substance of the theorem is the upper bound for arbitrary choices of the added rows. We prove that bound first on $b=c$.

The passage to $b=-c$ is exact and elementary. Number the columns from $0$ to $10$. Let $\sigma$ exchange columns $7$ and $8$ and fix the other columns. Direct substitution in the matrix defining $R(a,b,c)$ in Subsection~\ref{sec:DDDgenerictriple} gives
\[
R(a,b,c)\sigma=R(a,-b,c).
\]
In particular, $R(a,b,-b)\sigma=R(a,-b,-b)$. Column permutation preserves regularity and every row-extension problem, so the result for $R_+$ implies the result on the second surface.

\medskip
Use the first cycle presentations counted in the presentation-count table in Subsection~\ref{sec:DDDfull} and the phase-graph equations of Subsection~\ref{sec:DDDgenerictriple}, now after substituting $c=b$. The old entries are monomials in $a,b$ with fixed coefficients in $\mu_{20}$. There is no generic rejection rule in this calculation. An equation
\[
a^r b^s=w^{e}
\]
is retained as a condition on the old parameters, even when $(r,s)\ne(0,0)$.

The integer relations are kept as an unsaturated lattice with their root constants. This is necessary. For example, replacing $a^2=q$ by one chosen equation $a=\sqrt q$ would discard a component. The incremental integer row operations preserve the entire solution set of the equations. A zero character with a nonzero root constant rejects a branch; all consistent nonzero character relations remain in the record.

The phase graph still solves the relative phases of the new cycle blocks. After normalization, each record consists of the old-parameter lattice and a row family
\[
x_j=z_{g(j)}a^{p_j}b^{q_j}w^{e_j}.
\]
The first column in every independent phase component has monomial coefficient one. As before, one common row phase is irrelevant. Identical normalized records are retained once. Distinct records can intersect or describe specializations of one another; no ordinary-equivalence classification is implicit in this normalization.

The 2 new quotients are imposed successively against $u$ and $v$, with every support, every root labelling and every prime-cycle cover included. Profile $\mathrm{A}$ is impossible by Theorem~\ref{n11_A_branch:profiles}, since each old row pair has a chosen $\mathrm{D}$ decomposition. The complete counts are
\[
\begin{array}{c|r|r|r}
\text{first profile}&\text{first presentations}&
 \text{records after }u&\text{records after }v\\\hline
\mathrm{B}&712800&0&0\\
\mathrm{C}&443520&61608&3604\\
\mathrm{D}&166320&973758&19620\\
\mathrm{E}&123200&806720&0\\
\mathrm{F}&34650&10014968&256\\\hline
\text{total first presentations}&1480490&&
\end{array}
\]
Thus profiles $\mathrm{B}$ and $\mathrm{E}$ supply no fourth row anywhere on the surface. Profile $\mathrm{F}$ does supply rows and must be retained. Removing duplicate final records across the 5 scans gives exactly $23416$ records.

Let $r$ denote the rank of the old-parameter character lattice, and let $d$ be the number of additional row phases after removing the common phase. The final records have distribution
\[
\begin{array}{c|c|r}
r&d&\text{records}\\\hline
0&0&24\\
1&0&1296\\
2&0&22000\\
2&1&96\\\hline
&&23416.
\end{array}
\]
In particular, no 2-parameter row family survives, and a row with one additional free phase occurs only at a fixed old-parameter point. These statements follow from the complete retained lists, not from a dimension estimate.

\medskip
The $24$ rank-zero records have a particularly simple form. For these row lists and the constants below put $\omega=w^{14}=e^{\pi i/30}$. For $d\in\{1,a,b\}$, define 8 rows
\[
x^{(d)}_{ij}=
\begin{cases}
\omega^{E^{(d)}_{ij}},&0\le j\le6,\\
d\omega^{E^{(d)}_{ij}},&7\le j\le10,
\end{cases}
\qquad 0\le i\le7.
\]
The 3 coefficient matrices are the following.
\[
E^{(1)}=\begin{pmatrix}
0&36&12&48&24&24&24&4&4&44&44\\
0&36&12&48&24&24&24&44&44&4&4\\
0&36&12&48&39&39&54&9&9&24&24\\
0&36&12&48&39&39&54&24&24&9&9\\
0&48&36&24&12&12&12&32&32&52&52\\
0&48&36&24&12&12&12&52&52&32&32\\
0&48&36&24&27&27&42&12&12&57&57\\
0&48&36&24&27&27&42&57&57&12&12
\end{pmatrix}.
\]

\[
E^{(a)}=\begin{pmatrix}
0&36&12&48&24&39&9&26&16&56&46\\
0&36&12&48&24&39&9&46&56&16&26\\
0&36&12&48&39&9&24&6&21&36&51\\
0&36&12&48&39&9&24&51&36&21&6\\
0&48&36&24&12&27&57&14&4&44&34\\
0&48&36&24&12&27&57&34&44&4&14\\
0&48&36&24&27&57&12&39&24&9&54\\
0&48&36&24&27&57&12&54&9&24&39
\end{pmatrix}.
\]

\[
E^{(b)}=\begin{pmatrix}
0&36&12&48&9&39&24&33&18&48&3\\
0&36&12&48&9&39&24&48&3&33&18\\
0&36&12&48&39&24&9&8&58&28&38\\
0&36&12&48&39&24&9&28&38&8&58\\
0&48&36&24&27&12&57&16&26&56&46\\
0&48&36&24&27&12&57&56&46&16&26\\
0&48&36&24&57&27&12&21&6&36&51\\
0&48&36&24&57&27&12&36&51&21&6
\end{pmatrix}.
\]

All $24$ rows are regular against $R_+(a,b)$ for every $a,b\in\mathbb T$. Their coefficients are checked by exact prime-cycle decompositions of the 3 old/new quotients. For generic parameters the rows are distinct. At a special parameter, coincidences are permitted and are removed when the corresponding point list is formed.

The 6 compatible pairs as identities in $a,b$ are
\[
\begin{array}{c|cc}
d&\text{first pair of row indices}&\text{second pair of row indices}\\\hline
1&\{0,4\}&\{1,5\}\\
a&\{1,4\}&\{0,5\}\\
b&\{2,4\}&\{3,5\}.
\end{array}
\]
No identity compatibility joins rows from different values of $d$. Every one of the 6 resulting 5-row families is ordinarily equivalent to a specialization of the family $G(A,B,C)$ in Subsection~\ref{sec:DDDpositivefive}. The companion certificate records, for each pair, a row permutation, a column permutation, the dephasing factors and the explicit monomial substitutions for $A,B,C$. In the order of the 6 pairs in the preceding pair table, the substitutions are
\[
\begin{array}{c|c|c}
A&B&C\\\hline
-a&b&b\\
a&b&b\\
\omega^{26}/a&\omega^{58}b/a&\omega^{38}b/a\\
\omega^{56}/a&\omega^{58}b/a&\omega^{38}b/a\\
\omega^{6}a/b&\omega^{42}/b&\omega^{12}/b\\
\omega^{6}a/b&\omega^{12}/b&\omega^{42}/b.
\end{array}
\]
The file \texttt{5\_equivalences.json} records the corresponding row and column permutations. After applying them, dephase by dividing columnwise by the first row and then dividing each row by its first entry. Each of the $55$ resulting entry identities is checked, so this assertion does not rely on a numerical equivalence test. It also explains why these additional generic 5-row constructions have no sixth row: Theorem~\ref{n11_positive:main} applies to each of them at every parameter value.

For the full surface theorem we still need all non-generic compatibilities and all the other records in the rank-and-dimension table above. They are treated below without assuming that their 5-row extensions have this same form.

\medskip
Every rank-one record is supported on one of 3 kinds of circle:
\[
a=q,\qquad b=q,\qquad a/b=q.
\]
All their constants belong to $\mu_{60}$. More precisely, put $S=\mu_{20}\cup\mu_{30}$. Since $S$ has $20+30-10=40$ elements, the following sets each have $44$ elements:
\[
\begin{split}
S_a&=S\cup\{\omega^{13},\omega^{23},\omega^{43},\omega^{53}\},\\
S_b&=S\cup\{\omega,\omega^{11},\omega^{31},\omega^{41}\},\\
S_{a/b}&=S\cup\{\omega^{7},\omega^{17},\omega^{37},\omega^{47}\}.
\end{split}
\]
The rank-one records occur on exactly the $132$ circles obtained from $a=q$, $b=q$, $a/b=q$ using the respective sets $S_a,S_b,S_{a/b}$. Every such record has no further free row phase.

Every rank-2 record is supported on points of
\[
\mu_{60}^2.
\]
The union of their supports has $1936$ distinct points. This statement retains all torsion components. Choose a basis of the full rank-2 character lattice, carrying its root constants along. One way to solve its 2 equations is triangular integer reduction with rational angles. For an independent verification, apply Cramer's rule, allowing every integer shift of the right-hand sides modulo the absolute determinant. After removing repetitions, the resulting set has exactly that determinant many points. Every resulting coordinate has denominator dividing $60$.

The possible triangular diagonals and their numbers of rank-2 records are
$$\begin{array}{c|rrrrr}
\text{diagonal}&(1,1)&(1,2)&(1,4)&(2,1)&(2,2)\\\hline
\text{records}&20408&1040&64&520&64.
\end{array}$$
The $96$ records with an additional row phase become 6 distinct circles of fourth rows at each of the $16$ parameter points
\[
(a,b)\in\omega^{3}\mu_4\ \times\ \omega^{6}\mu_4.
\]
No free row circle occurs at another parameter point.

For the compatibility calculation we use the slightly larger collection of all $180$ circles in $a=q$, $b=q$, $a/b=q$ with $q\in\mu_{60}$, and all $3600$ points in $\mu_{60}^2$. The extra circles are useful because a pair of generic fourth rows can acquire a new compatibility even where no new fourth-row record appears. Keeping the whole collection also makes the partition of the parameter space explicit.

Outside these $180$ circles, only the $24$ generic rows occur. At a point of one circle with its remaining parameter outside $\mu_{60}$, no other circle in the collection passes through that point. At points of $\mu_{60}^2$, all applicable generic, circle-supported and fixed-point records are combined. Thus every intersection is included in the finite-point calculation.

\medskip
A quotient of 2 generic rows is either constant or has 2 character groups, with coefficients in $\mu_{60}$ and relative character $t\in\{a,b,a/b\}$, up to inversion. In the latter case, if a prime cycle mixes the groups, then $t\in\mu_{420}$. An $11$-cycle is impossible, since it would use at most one column from each group. If no cycle mixes the groups, regularity is an identity in $t$. Testing all $276$ pairs shows that nonidentity compatibilities occur only in $\mu_{60}$. The identity graph has $6$ disjoint edges and $12$ isolated vertices, excluding $3$ added rows away from the $180$ circles.

Fix one of the $180$ circles. Use a remaining parameter $t$ as follows:
\[
\begin{array}{c|cc}
\text{circle}&a&b\\\hline
a=q&q&t\\
b=q&t&q\\
a/b=q&qt&t.
\end{array}
\]
Substitute these expressions into every row record whose old-parameter equations hold identically on the circle. Normalize and remove duplicates. This gives a finite list of rows
\[
x_j(t)=w^{e_j}t^{k_j},\qquad e_j\equiv0\pmod{14}.
\]
The congruence means that the fixed coefficients are still in $\mu_{60}$. There are $4512$ row records across the $180$ circle lists, with at most $44$ rows on any one circle. These counts keep separate appearances of a row on different circles.

For each pair of rows in one list, the character exponents of their quotient have range at most 2:
\[
\max_j(k_j-k'_j)-\min_j(k_j-k'_j)\le2.
\]
This is checked for every pair, and independently after reconstructing the circle lists directly from all old-parameter equations.

Suppose a regular cycle mixes 2 distinct character exponents. Their difference is $d$ with $|d|=1$ or $2$. The entry ratio then gives
$$t^d\in\mu_{420},\qquad\text{hence }t\in\mu_{840}.$$
An $11$-cycle is again impossible. By the bound $\max_j(k_j-k'_j)-\min_j(k_j-k'_j)\le2$ there are at most 3 character groups, and it could contain at most one entry from each group.

It follows that for $t\notin\mu_{840}$ every regular cycle lies within a single character group. The exact compatibility graph on this whole complement is obtained by checking prime-cycle covers separately in those groups. For $t\in\mu_{840}\setminus\mu_{60}$, substitute the root and compute the ordinary fixed-row compatibility graph. The cases $t\in\mu_{60}$ are reserved for the full finite-point lists, which may contain additional rows.

There are $840-60=780$ root values to test on each circle, and one further graph for its entire complement of $\mu_{840}$. The complete calculation therefore gives
\[
180(780+1)=140580
\]
compatibility graphs. Every one is triangle-free. Their edge counts are
$$\begin{array}{c|rrrr}
\text{edges}&6&14&40&48\\\hline
\text{graphs}&131208&4686&4674&12.
\end{array}$$
Thus this step covers every parameter on the special circles outside $\mu_{60}^2$, including roots of arbitrarily high order and parameters which are not roots of unity.

\medskip
At each of the $3600$ points of $\mu_{60}^2$, specialize every applicable record and merge equal row families. There are $101768$ families in total. At the $16$ points $\omega^3\mu_4\times\omega^6\mu_4$, the list includes $6$ circles
\[
x_j(z)=w^{e_j}z^{g_j},\qquad g_j\in\{0,1\}.
\]
At every other point all rows are fixed. A circle supplies at most one orthogonal row by Lemma~\ref{n11_fourth:circle}. Between distinct families, the rectangle-phase and two-variable congruence procedures of Subsection~\ref{sec:DDDfourth} give the same compatibility graphs, with actual fifth-row witnesses. All $3600$ graphs are triangle-free and have an edge. Together with the open and circle calculations this proves the upper bound $5$ everywhere. The family $G(a,b,b)$ attains it; the sign symmetry gives $b=-c$.
\end{proof}

Both complete scans agree on their final $23416$ records. The audit checks every consistency decision using primitive integer relations. Direct substitution independently reconstructs every circle and point list, retaining all torsion components. The compatibility graphs are recomputed with alternative arithmetic and cycle traversal; at points with free row circles the 2-variable congruence solver gives the same edges. Symbolic row checks, fifth-row witnesses and triangle tests complete the verification.

\subsection{2 further exceptional surfaces and simultaneous phase consistency}\label{sec:DDDImaginary}

We now treat $b/c=\pm i$, including all special curves and points on these surfaces.

There is a new feature in this calculation. At 4 old-parameter points, the graph of possible pairwise compatibilities contains triangles. Each vertex of such a triangle represents a family of rows, and an edge allows its own choice of the 2 row phases. The 3 edges need not admit one common choice of the 3 phases. We therefore retain and intersect the exact phase equations for the edges. In the present calculation every triangle fails this simultaneous test.

Put $w=e^{\pi i/420}$, $\omega=w^{42}=e^{\pi i/10}$ and $R_i(a,b)=R(a,b,-ib)$. Explicitly,
\[
R_i(a,b)=
\begin{pmatrix}
1&1&1&1&1&1&1&1&1&1&1\\
1&\omega^{4}&\omega^{8}&\omega^{12}&\omega^{16}&\omega&\omega^{11}&-ia&ia&a&-a\\
1&\omega^{8}&\omega^{16}&\omega^{4}&\omega^{17}&\omega^{12}&\omega^{7}&b&-b&-ib&ib
\end{pmatrix}.
\]
The 2 parameters $a,b$ range independently over $\mathbb T$. Columns in this subsection are numbered from $0$ to $10$.

\par\addvspace{\medskipamount}\noindent
\begin{minipage}{\linewidth}
\begin{theorem}\label{n11_imaginary:main}
For every $a,b\in\mathbb T$, the largest number of rows in a regular extension of $R(a,b,-ib)$ is exactly 5. The same assertion holds for $R(a,b,ib)$. Consequently, for the specified triple $R(a,b,c)$, every parameter value with $b/c\in\mu_4$ has maximum regular extension equal to 5. None of these triples can occur in a regular square matrix of order 11.
\end{theorem}
\end{minipage}
\par\addvspace{\medskipamount}

\begin{proof}
The existing family $G(a,b,c)$ supplies the lower bound for all parameters. The issue is to exclude arbitrary choices of 3 added rows. We prove the upper bound on $R_i$; the sign symmetry $R(a,b,c)\sigma=R(a,-b,c)$ then transfers it to the surface $b/c=-i$. The last assertion about $\mu_4$ combines this result with Theorem~\ref{n11_surface:main}.

\medskip
The input entries in $R_i(a,b)$ are monomials in $a,b$ with coefficients in $\mu_{20}$. Starting with a cycle presentation of the quotient of a new row by the first row, impose regularity against the second and third rows successively. As in the preceding subsection, every support and every root labelling is included. There is no restriction to a numerical set of old parameters. Profile $\mathrm{A}$ is impossible by Theorem~\ref{n11_A_branch:profiles}, because the old pairs have chosen $\mathrm{D}$ decompositions. Thus the first profile ranges over $\mathrm{B},\mathrm{C},\mathrm{D},\mathrm{E},\mathrm{F}$. Write $u,v$ for the second and third rows of $R_i$.

We retain every consistent old-parameter relation
\[
a^r b^s=w^{e}.
\]
An equation with zero character is inconsistent precisely when its constant residue $e$ is nonzero modulo $840$. Consistent nonzero character equations are retained as conditions on $a,b$. Nonprimitive relations are retained without choosing a root. The resulting lattice is not saturated: for example, a relation on $a^2$ retains both possible values of $a$ whenever it is later solved. This convention is essential for the finite intersections in this calculation.

The normalized fourth-row records have the form
\[
x_j=z_{g(j)}a^{p_j}b^{q_j}w^{e_j},
\qquad 0\leq j\leq10,
\]
together with their old-parameter lattice. The first coefficient in each independent phase component is one. One common row phase is removed. Equality of normalized records is used to remove duplicates; no assertion about ordinary equivalence of different records is made.

The full enumeration gives the following counts.
\[
\begin{array}{c|r|r|r}
\text{first profile}&\text{first presentations}&\text{records after }u&\text{records after }v\\\hline
\mathrm{B}&712800&0&0\\
\mathrm{C}&443520&30748&2692\\
\mathrm{D}&166320&1134000&11076\\
\mathrm{E}&123200&518016&0\\
\mathrm{F}&34650&10393080&176\\
\hline
\text{total first presentations}&1480490&&
\end{array}
\]
Profiles $\mathrm{B}$ and $\mathrm{E}$ give no fourth row. Profile $\mathrm{F}$ must be retained. After removing duplicate final records across the 5 scans, the list has exactly $13944$ records.

Let $r$ be the rank of the old-parameter lattice, and let $d$ count additional row phases after removal of the common phase. The distribution is
\[
\begin{array}{c|c|r}
r&d&\text{records}\\\hline
0&0&8\\
1&0&552\\
2&0&13272\\
2&1&112\\
\end{array}
\]
In particular, free row circles occur only at the fixed old-parameter points supplied by rank-2 records. Each rank-one record has no additional free row phase; its row is determined by the old parameters, up to a common row phase.

\medskip
The rank-zero part consists of 8 constant rows. For the row lists and constants below put $\omega=w^{14}=e^{\pi i/30}$ and define $x^{(k)}_j=\omega^{E_{kj}}$ for $0\leq k\leq7$, where
\[
E=\begin{pmatrix}
0&36&12&48&24&24&24&4&4&44&44\\
0&36&12&48&24&24&24&44&44&4&4\\
0&36&12&48&39&39&54&9&9&24&24\\
0&36&12&48&39&39&54&24&24&9&9\\
0&48&36&24&12&12&12&32&32&52&52\\
0&48&36&24&12&12&12&52&52&32&32\\
0&48&36&24&27&27&42&12&12&57&57\\
0&48&36&24&27&27&42&57&57&12&12
\end{pmatrix}.
\]
These are the 8 constant rows already obtained in the generic 3-parameter calculation. Here they remain present on the whole surface, including its exceptional curves and points.
Each row is regular against all 3 old rows for every $a,b$. This is checked by prime-cycle decompositions with the parameter characters left formal. The fixed compatibility graph on these 8 rows has exactly 2 edges,
\[
\{0,4\},\qquad\{1,5\}.
\]
It has no triangle. Because these rows are constant, no new compatibility between them can appear by specializing $a$ or $b$. This distinguishes the present calculation from the $24$ generic rows on $b=c$, where additional pair compatibilities had to be included in the exceptional set.

\medskip
The rank-one records are supported on 5 kinds of connected curve:
\[
a=q,\qquad b=q,\qquad a/b=q,\qquad a/b^2=q,\qquad a^2/b=q.
\]
Every constant $q$ belongs to $\mu_{60}$. There are exactly $140$ distinct connected curves in their union. Their constants, written as exponents of $\omega$, are
\[
\begin{array}{c|l}
\text{character}&\{e:q=\omega^{e}\}\\\hline
a&0,1,\ldots,59\\
b&0,1,\ldots,59\\
a/b&0,6,12,18,24,27,30,36,42,48,54,57\\
a/b^2&6,21,36,51\\
a^2/b&0,15,30,45.
\end{array}
\]
These are the supports of the retained records. A record on a curve is allowed to coincide with a generic row, or with a specialization of another record. Such coincidences are removed when the actual row list on that stratum is formed.

All rank-2 records are supported in $\mu_{120}^2$. Their support union consists of exactly $1920$ points. The diagonal entries in their triangular integer presentations have the following distribution:
\[
\begin{array}{c|r}
\text{diagonal entries}&\text{records}\\\hline
(1, 1)&11520\\
(1, 2)&1232\\
(1, 4)&16\\
(2, 1)&600\\
(2, 2)&16\\
\end{array}
\]
These counts include all components of the nonsaturated equations. Choose 2 equations forming a basis of the full character lattice, with their corresponding root constants, and write their coefficient matrix with determinant $\Delta\ne0$. Cramer's rule, applied with every integer shift of the 2 root angles modulo $|\Delta|$, gives exactly $|\Delta|$ distinct points after repetitions are removed. All coordinate denominators in the retained list divide $120$.

The intersections of distinct nonparallel curves in the displayed table of characters and constants give $4080$ distinct points, again all in $\mu_{120}^2$. This is checked from the actual constants in the table, including the 2 nonlinear characters. Distinct parallel connected curves are disjoint. We therefore partition the parameter torus as follows: the complement of all listed curves and of $\mu_{120}^2$; each listed curve with its points of $\mu_{120}^2$ removed; and all $14400$ points of $\mu_{120}^2$. Outside the curves and this finite set, only the 8 constant rows occur. On the punctured curves, all applicable rank-zero and rank-one records are used. At a finite point, all applicable records of every rank are combined.

For clarity, the use of the larger finite set $\mu_{120}^2$ is an exact covering device. It is not a phase grid intended to approximate arbitrary old parameters. The remaining parameters on each punctured curve and on the open complement are treated symbolically.

\medskip
Use the following coordinates, with $t\in\mathbb T$:
\[
\begin{array}{c|c}
\text{curve}&(a,b)\\\hline
a=q&(q,t)\\
b=q&(t,q)\\
a/b=q&(qt,t)\\
a/b^2=q&(qt^2,t)\\
a^2/b=q&(t,q^{-1}t^2).
\end{array}
\]
In every case, $(a,b)\in\mu_{120}^2$ if and only if $t\in\mu_{120}$. After substitution, each retained row is
\[
x_j(t)=c_j t^{k_j},\qquad c_j\in\mu_{60}.
\]
The complete curve lists have $1376$ row presentations in total, and at most $36$ distinct rows on any one curve. For any 2 rows on a curve, the integer exponents in their quotient satisfy
\[
\max_j(k_j-\ell_j)-\min_j(k_j-\ell_j)\leq2.
\]
Both the coefficient group and this exponent bound are checked directly after substitution into every original record.

Suppose a prime cycle in such a quotient joins positions with different powers of $t$. For some $d\in\{1,2\}$, the cycle equations then imply
\[
t^d\in\mu_{\operatorname{lcm}(60,p)}.
\]
Here $p\in\{2,3,5,7\}$. It follows that $t\in\mu_{840}$. If no cycle mixes different powers of $t$, regularity is an identity in $t$ and is checked separately within each exponent group. Thus one formal calculation covers every $t\notin\mu_{840}$, while the remaining possibilities are tested exactly in $\mu_{840}$.

An $11$-cycle cannot have been omitted from this argument. 2 positions with the same exponent of $t$ have quotient in $\mu_{60}$, so they cannot be distinct positions in an $11$-cycle. By the bound $\max_j(k_j-\ell_j)-\min_j(k_j-\ell_j)\le2$ there are at most 3 exponent groups. An $11$-cycle would require 11 of them.

The points with $t\in\mu_{120}$ belong to the finite-point part of the partition. Each curve therefore requires $840-120=720$ exact root calculations and one formal complementary calculation. The total number of curve graphs is
\[
140\,(720+1)=100940.
\]
Their numbers of edges have distribution
\[
\begin{array}{c|r}
\text{edges}&\text{graphs}\\\hline
2&93730\\
6&5768\\
36&1442\\
\end{array}
\]
Every one of these graphs is triangle-free. Consequently no point of a punctured curve admits 3 pairwise regular added rows.

\medskip
At every $(a,b)\in\mu_{120}^2$, substitute the old parameters into every applicable lattice record and normalize each row family again. The combined lists have $149616$ distinct normalized family presentations, counted separately at each point. At $14384$ points every family is a fixed row. At the remaining $16$ points, some families contain one free phase and have the form
\[
x_j(s)=c_j s^{g_j},\qquad
c_j\in\mu_{840},\quad g_j\in\{0,1\},\quad s\in\mathbb T.
\]
The free families occur at precisely
\[
(a,b)\in\omega^{3}\mu_4\ \times\ \omega^{6}\mu_4.
\]
At the 4 points with $ab=\omega^{54}$ there are 8 distinct free row circles; at the other $12$ points there are 6. The fixed rows at those points are retained as well.

A graph vertex is one normalized family, fixed or free. 2 vertices are joined when there exist choices of their phases making their quotient regular. This gives an upper bound for actual compatible collections of rows. It need not identify a simultaneous phase choice for a clique.

There is also no missing possibility involving 2 rows from the same free family. If its 2 exponent groups in $x_j(s)=c_js^{g_j}$ have $m$ and $11-m$ entries, the inner product of 2 of its rows has the form
$$m+(11-m)s/t.$$
For $s,t\in\mathbb T$, this could vanish only if $m=11-m$, which is impossible. Thus 2 different rows from a single family cannot be orthogonal. A regular 6-row extension must use 3 distinct vertices of the family graph.

The maximum clique cardinalities in the point graphs are
\[
\begin{array}{c|r}
\text{maximum clique cardinality}&\text{old-parameter points}\\\hline
2&14396\\
3&4\\
\end{array}
\]
The 4 points with triangles are
\[
(a,b)=(\omega^{3},\omega^{21})\,(i^k,(-i)^k),\qquad k=0,1,2,3.
\]
They form one exact symmetry orbit. The column permutation with last 4 columns in the order $(9,10,8,7)$, fixing columns $0,\ldots,6$, sends $R_i(a,b)$ to $R_i(ia,-ib)$. Every entry of this identity is checked directly.

There are $176$ graph triangles in total, with $44$ at each of these points. Of them, $160$ use one fixed row and 2 free row circles, and $16$ use 3 free row circles. These triangles require an additional argument; a triangle-free claim at this stage would be false.

\medskip
For 2 point families $x(s)$ and $y(t)$, their quotient entries have the form
\[
c_jd_j^{-1}s^{g_j}t^{-h_j}.
\]
Exhaust every prime-cycle cover of these 11 entries. Each cover imposes integer character equations
\[
s^u t^v=w^{e}.
\]
We retain the full consistent lattice for that cover. Taking the union over covers gives an exact description of all compatible phase pairs. Rank-one loci remain circles, with every torsion component retained; rank-2 loci remain their complete finite sets. The phases themselves have not been assumed to belong to any prescribed finite root group.

Here too an $11$-cycle is impossible. The exponents $(g_j,h_j)$ give at most 4 cells. Within one cell, the ratio of 2 coefficients belongs to $\mu_{840}$, whose intersection with $\mu_{11}$ is trivial. An $11$-cycle would therefore use at most one position per cell, fewer than 11 positions.

For a graph triangle with row phases $s,t,u$, choose one retained lattice for each of its 3 edges. Embed the 3 systems in the common character group on $(s,t,u)$. A simultaneous solution exists if and only if their union of equations is consistent. Every choice of the 3 edge lattices is checked. The complete enumeration uses $312$ distinct edges belonging to triangles, $600$ normalized pair-locus records, and $3296$ combinations of 3 pair loci. None is consistent.

The exclusion can be checked without trusting a consistency verdict from the enumeration. Write the combined equations in angle form as
\[
A\begin{pmatrix}\arg(s)/(2\pi)\\\arg(t)/(2\pi)\\\arg(u)/(2\pi)\end{pmatrix}
\equiv\frac{e}{840}\pmod{\mathbb Z^m},
\]
where $A$ is an integer matrix and $e$ an integer column. For each combination, the certificate supplies an integer row vector $r$ satisfying
\[
rA=0,\qquad re\not\equiv0\pmod{840}.
\]
Multiplying the displayed system of three phase equations by $r$ gives an immediate contradiction. The independent verifier reconstructs all equations from the 3 saved pair loci and checks the identities $rA=0$ and $re\not\equiv0\pmod{840}$ with integer arithmetic. There is one such certificate for every one of the $3296$ combinations. It also verifies every pair-locus record symbolically by a prime-cycle cover of $c_jd_j^{-1}s^{g_j}t^{-h_j}$.

Thus the triangles in the family graphs do not yield any actual regular 6-row extension. Together with the open and curve parts of the parameter partition, this proves the upper bound in Theorem~\ref{n11_imaginary:main}. The lower bound is attained everywhere, and the sign symmetry completes the proof on both surfaces.
\end{proof}

\medskip
The obstruction just used can be seen in a small exact example. Here put $w=e^{\pi i/30}$ and $\omega=w^3=e^{\pi i/10}$. Take
$$(a,b,c)=(\omega,\omega^{7},\omega^{2}),$$
which is the first point in the four-point list above. Define 3 row families by
\[
x_j=\omega^{X_j},\qquad y_j(s)=\omega^{Y_j}s^{g_j},\qquad
z_j(t)=\omega^{Z_j}t^{h_j},
\]
with coefficient and exponent vectors
\[
\begin{array}{c|rrrrrrrrrrr}
j&0&1&2&3&4&5&6&7&8&9&10\\\hline
X_j&0&12&4&16&3&8&8&13&13&18&3\\
Y_j&0&16&12&8&0&15&4&19&10&9&5\\
Z_j&0&16&12&8&0&10&14&4&5&15&4\\
g_j&0&0&0&0&1&1&0&0&1&0&1\\
h_j&0&0&0&0&1&1&0&0&1&1&0.
\end{array}
\]
Every one of these rows is regular against the old triple; the assertions for $y(s)$ and $z(t)$ hold for every $s,t\in\mathbb T$. They follow by separating the 2 phase groups and checking their prime-cycle decompositions against each old row.

The inner product of $x$ with $y(s)$ has the form $A+s^{-1}B$, where $A\ne0$ and $B=-\omega^{19}A$. Hence orthogonality already forces $s=\omega^{19}$. At this value the quotient is indeed regular. The same calculation for $x$ and $z(t)$ forces $t=\omega^{19}$, again with a regular quotient. On the other hand, $y(s)$ and $z(t)$ are regular at either of the phase pairs
\[
(s,t)=(w^{17},w^{2})
\quad\hbox{or}\quad
(s,t)=(w^{37},w^{22}).
\]
Thus all 3 edges of the family triangle exist.

They cannot be realized together. The 2 edges incident to $x$ force the common choice $s=t=\omega^{19}$, and direct evaluation gives
\[
\langle y(\omega^{19}),z(\omega^{19})\rangle=5.
\]
All statements here are exact. The scalar identities are verified by polynomial reduction modulo
$$\Phi_{20}(X)=X^8-X^6+X^4-X^2+1,$$
and the asserted regular quotients are verified by their prime-cycle decompositions. The example explains why the family graph is an intermediate object: its edges answer separate existence questions, while a regular extension requires one common row choice at each vertex.

\medskip

The complete scans agree on all $13944$ row records. Independent substitution reconstructs the curve and point lists. The 2 compatibility calculations agree, including the phase loci for triangle edges. A separate integer verifier checks all $600$ pair loci and all $3296$ contradictions. These computations settle $b/c=\pm i$ with arbitrary remaining parameters; together with Theorem~\ref{n11_surface:main}, they settle $b/c\in\mu_4$.

\subsection{A level-120 Butson test on a further DDD slice}\label{sec:DDDButson120}

Put $w=e^{\pi i/60}$ and $\omega=w^{40}=e^{2\pi i/3}$. The preceding $2$ subsections treated complete exceptional surfaces with unrestricted unit parameters. We now consider a different, finite question. Fix the level of every matrix entry to divide $120$, and study the remaining surface $b/c=\omega$ for the same labelled triple $R(a,b,c)$. This fixes the coefficient group of every old and new row. It permits a direct enumeration of actual fourth rows and their compatibility graphs, including every allowed value of a remaining row phase.

The restriction on the added rows is part of the question. Even when an old triple has root-of-unity entries, an unrestricted regular extension may use phases outside that root group. Accordingly, the theorem below is a fixed-level result. It does not assert the corresponding upper bound for arbitrary unit entries.

The primitive $20$th root is now $w^6=e^{\pi i/10}$. We write
\[
R_\omega(a,b)=R(a,b,\omega^{-1}b)=
\begin{pmatrix}
1&1&1&1&1&1&1&1&1&1&1\\
1&w^{24}&w^{48}&w^{72}&w^{96}&w^{6}&w^{66}&-\omega a&\omega a&a&-a\\
1&w^{48}&w^{96}&w^{24}&w^{102}&w^{72}&w^{42}&b&-b&\omega^{-1}b&-\omega^{-1}b
\end{pmatrix}.
\]
In this subsection $a,b\in\mu_{120}$. There are exactly $120^2=14400$ labelled parameter pairs. Columns are numbered from $0$ to $10$, and the common phase of every added row is normalized by requiring its entry in column $0$ to be one.

\par\addvspace{\medskipamount}\noindent
\begin{minipage}{\linewidth}
\begin{theorem}\label{n11_bh120:main}
For every $a,b\in\mu_{120}$ and $q\in\{\omega,\omega^{-1},-\omega,-\omega^{-1}\}$, the largest number of rows in a regular extension of $R(a,b,q^{-1}b)$ with all entries in $\mu_{120}$ is exactly 5. None of these 4 level-120 slices of the specified labelled $\mathrm{DDD}$ triple admits a regular 6-row extension with entries in $\mu_{120}$. Combined with Theorems~\ref{n11_surface:main} and \ref{n11_imaginary:main}, the same fixed-level conclusion holds whenever $b/c\in\mu_4\cup\mu_6$. No such triple can occur as the displayed triple in a regular $\mathrm{BH}(11,120)$ matrix.
\end{theorem}
\end{minipage}
\par\addvspace{\medskipamount}

\begin{proof}
We prove the statement on $b/c=\omega$. The column transposition $(7\ 8)$ gives the exact identity
\[
R(a,b,c)^{(7\ 8)}=R(a,-b,c).
\]
It sends the ratio $b/c$ to its negative and preserves the entry group $\mu_{120}$. A second column permutation, which exchanges the final 2 column pairs and fixes the first 7 columns, gives
\[
R(a,b,c)^{(7\ 9)(8\ 10)}=R(-ab/c,c,b).
\]
This transformation inverts $b/c$. Its action on $(a,b,c)$ is an involution and preserves $\mu_{120}^3$. Both identities are checked entry by entry on the original symbolic triple. They transfer the calculation on $b/c=\omega$ to the 4 distinct ratios in the theorem, with all remaining parameters included. The 2 ratios $1,-1$ are already covered by the earlier unrestricted surface theorem, and $i,-i$ by the imaginary-surface theorem. Thus $\mu_4\cup\mu_6$, which has 8 elements, follows by combining the results. Only the new 4-ratio assertion is being computed in the present subsection.

\medskip

In $\mu_{120}$ a prime-cycle length must divide $120$, so only $2,3,5$ occur. The possible profiles are precisely $\mathrm{C},\mathrm{D},\mathrm{E},\mathrm{F}$, with $443520,166320,123200,34650$ first presentations. Anchoring each cycle at its least column and trying every remaining root order gives the complete input; its cycle multiplier stays free until the level condition is imposed.

\medskip
Write $u,v$ for the second and third rows of $R_\omega(a,b)$. Starting from a first-row cycle presentation of a possible new row $x$, impose regularity against $u$ and then against $v$. Each cycle equation identifies 2 row-phase components up to a monomial in $a,b$ and a power of $w$. Closing a chain of such identifications gives an old-parameter equation
\[
a^r b^s=w^{e},
\qquad r,s\in\mathbb Z,\quad e\in\mathbb Z/120\mathbb Z.
\]
All these equations are retained in an integer triangular lattice. A relation with zero character and nonzero constant is inconsistent. A nonprimitive relation is retained as written, with all of its solutions; the lattice is not saturated.

The fixed-level hypotheses enter as the 2 additional relations
\[
a^{120}=1,\qquad b^{120}=1.
\]
The first is present from the start. The second may be inserted after regularity against $u$ has been imposed, because $u$ contains no $b$ after the substitution $c=\omega^{-1}b$. No equation involving $b$ can arise before that stage. This ordering reduces the intermediate record while retaining exactly the same final solution set.

Each normalized final record consists of its old-parameter lattice and a row family
\[
x_j=z_{g(j)}a^{p_j}b^{q_j}w^{e_j},
\qquad 0\leq j\leq10.
\]
The coefficient at the first position in each phase component is normalized to one, and the phase of the component containing column $0$ is fixed to one. Every other $z_g$ is independent before the fixed-level restriction on the new row is applied. Since $a,b$ and the coefficients already belong to $\mu_{120}$, the row belongs to $\mu_{120}^{11}$ if and only if every remaining component phase belongs to $\mu_{120}$. Thus the later finite phase expansion is exact in both directions.

This gives completeness of the enumeration. Any regular added row has one of the first presentations in the profile list $\mathrm{C},\mathrm{D},\mathrm{E},\mathrm{F}$. Its quotients by $u$ and $v$ have prime-cycle covers using the same 3 prime lengths. The enumerator visits these covers and retains their consistent equations. Conversely, every retained record carries such covers against all 3 old rows, and every specialization satisfying its equations gives a regular added row. The period relations select exactly the prescribed old parameters.

\medskip

The full scans against $u,v$ leave respectively $2020,8852,0,412$ records from $\mathrm{C},\mathrm{D},\mathrm{E},\mathrm{F}$. Their union has $11252$ distinct records: $11188$ fixed rows and $64$ families with one extra row phase. The old-parameter equations include both period relations. Primary and alternative scans agree, and every record is checked symbolically against the old triple.

\medskip

Put $a=w^{r},b=w^{s}$. Solve every saved character equation modulo $120$ at all $14400$ pairs $(r,s)$, retaining every congruence solution. After specialization and normalization there are $132024$ point-family incidences. Free row phases occur at $16$ points, and each is expanded through all $120$ values. Removing duplicate actual rows gives $137784$ row incidences. An independent audit substitutes all parameter pairs directly into every original record and reconstructs these same family and row lists.

\medskip

At each $(r,s)$ use the actual normalized rows as vertices, joining $x,y$ exactly when $x\bar y$ has a cover by $2$-, $3$- and $5$-cycles. Since every vertex already extends the old triple, a $k$-clique is exactly a choice of $k$ simultaneous new rows. There are no remaining free phases in the vertices.

All $3142996$ row pairs are tested, giving $386136$ edges over $14400$ graphs. Every graph has an edge and no triangle; the largest has $758$ vertices. A second cycle-cover implementation gives the same edge lists. The independent verifier checks every edge and nonedge and excludes triangles by intersecting endpoint neighbourhoods. Consequently every maximum regular extension in the prescribed root group has exactly $5$ rows.

\medskip
The uniform 5-row lower bound can be displayed without the enumeration. In addition to the 3 old rows, use the 2 constant rows $x,y$ given by
\[
\begin{aligned}
x&=(w^{X_j})_{j=0}^{10},
&X&=(0,72,24,96,48,48,48,8,8,88,88),\\
y&=(w^{Y_j})_{j=0}^{10},
&Y&=(0,96,72,48,24,24,24,64,64,104,104).
\end{aligned}
\]
These are 2 of the constant rows from the generic calculation. Their quotients against each old row admit prime-cycle decompositions with $a,b$ left formal, and their mutual quotient admits such a decomposition as well. The certificate checks these identities symbolically. Since the 2 rows have entries in $\mu_{60}\subset\mu_{120}$, they supply a 5-row extension at every pair in the theorem.

The graph calculation excludes 3 simultaneous added rows for every parameter pair, while the displayed pair $x,y$ supplies 2. This proves the assertion of Theorem~\ref{n11_bh120:main} on $b/c=\omega$. The identities $R(a,b,c)^{(7\ 8)}=R(a,-b,c)$ and $R(a,b,c)^{(7\ 9)(8\ 10)}=R(-ab/c,c,b)$ supply the other 3 ratios. The earlier surface theorems supply the remaining ratios in $\mu_4\cup\mu_6$. A regular square extension with entries in $\mu_{120}$ would contain a regular 6-row extension, so the final assertion follows.
\end{proof}

\par\addvspace{\medskipamount}\noindent
\begin{minipage}{\linewidth}
\begin{example}
For a concrete instance whose minimal Butson level is $120$, take $a=w$, $b=1$ and $c=w^{80}=\omega^{-1}$. With the 2 rows above, the resulting partial matrix is $H_{ij}=w^{E_{ij}}$, where
\[
E=\begin{pmatrix}
0&0&0&0&0&0&0&0&0&0&0\\
0&24&48&72&96&6&66&101&41&1&61\\
0&48&96&24&102&72&42&0&60&80&20\\
0&72&24&96&48&48&48&8&8&88&88\\
0&96&72&48&24&24&24&64&64&104&104
\end{pmatrix}.
\]
Every row pair has an explicit prime-cycle partition, retained in the certificate. The first row and first column are zero in this exponent display, so $H$ is dephased. Its entry at row $1$, column $9$ is $w$. The associated 4-entry quotient with row $0$ and column $0$ is therefore a primitive $120$th root. Such 4-entry quotients are preserved up to permutation by ordinary equivalence, so every Butson representative has level divisible by $120$. The displayed matrix has level $120$, proving minimality.
\end{example}
\end{minipage}
\par\addvspace{\medskipamount}

\medskip

This is a complete result for the stated level-$120$ slices of $R$. The corresponding unrestricted surfaces, other labelled $\mathrm{DDD}$ triples, and the existence of regular $BH(11,120)$ matrices in other branches remain open.

\subsection{Two- and three-parameter DDD classes}\label{sec:ddd_pilots}

Class $3250$ provides a two-parameter case with a fixed row pair. Its chosen supports for $u,v,u/v$ are
$$\begin{aligned}
P&=((0,1,2,3,4),(5,6),(7,8),(9,10)),\\
Q&=((0,1,2,3,5),(4,6),(7,8),(9,10)),\\
R&=((0,1,2,3,7),(4,9),(5,10),(6,8)).
\end{aligned}$$
The pairs $(7,8)$ and $(9,10)$ rotate independently in the first two quotients. Their quotient is fixed.

\par\addvspace{\medskipamount}\noindent
\begin{minipage}{\linewidth}
\begin{theorem}\label{ddd_pilots:3250}
Every regular extension in support class $3250$ has at most 9 rows. In particular the class admits no regular square completion.
\end{theorem}
\end{minipage}
\par\addvspace{\medskipamount}

\begin{proof}
The Smith diagonal of the phase matrix is $(1,1,1,2)$. There are 72 normalized pentagon labellings and 2 binary components for each. Rotate the fixed quotient to the second row, as in Theorem~\ref{ddd_one:bound}. The verifier reconstructs all 144 presentations from their phase equations. Normalizing one coefficient in each of the 3 parts of the third row, interchanging the names of the two free phases and sorting columns within equal entries of the fixed row gives 3 descriptions.

With $w=e^{\pi i/10}$, the fixed row is again $u_j=w^{b_j}$ with $b=(0,0,1,5,5,9,10,10,13,15,17)$. The third row is $w^{c_j}$ multiplied by $z$ in columns $3,4$ and by $t$ in columns $6,7$. In the 3 descriptions
$$\begin{aligned}
c&=(0,5,c_2,0,10,c_5,0,10,c_8,10,c_{10}),\\
(c_2,c_5,c_8,c_{10})&\in\{(9,1,17,13),(13,17,9,1),(17,13,1,9)\}.
\end{aligned}$$

For each of the $339008$ records $x$ of the fixed pair, impose all prime-cycle covers of $x/v$. As before, the old $\mathrm{D}$ quotient excludes a new $\mathrm{A}$ quotient. Forest integration expresses each group phase of $x$ as a Laurent monomial in $z,t$, a constant in $\mu_{840}$, and one free phase per component. Every closing edge gives an integer character equation in $z,t$. Retain its complete solution set, including every torsion branch. Two solvers give identical normalized records. At ranks $0,1,2$ the counts over the 3 descriptions are $312,8004,67568$, respectively. The verifier tests a literal cycle cover against every old row for each record.

A rank-one domain splits into curves $z^a t^b=c$ with $(a,b)$ primitive. A rank-two domain gives finitely many points. For each description first test the records valid on the whole torus. On each curve adjoin all records valid along it. At all specified points and all intersections of nonparallel curves specialize every applicable record. This covers the full torus. Each description has 136 curves and 3280 points after intersections are included. The verifier independently reconstructs the domains by integer determinants and checks all graph memberships.

There are $10251$ graphs in total. On each domain remove repeated records and individual rows contained in a retained circle, testing ratios of monomials on its two blocks. Every record has at most 2 groups, so two orthogonal rows require distinct records. Use the coefficient-sum tests of Theorem~\ref{ddd_one:bound}, now allowing both old phases to be free. Two calculations, modulo 3361 and 20161, bound every clique by 6. Thus at most 6 rows can be added to the original triple.
\end{proof}

Class $3297$ gives a different three-parameter problem. All 3 pentagons occupy the first 5 columns. On the final 6 columns, the pairings for $u,v,u/v$ are respectively
$$((5,6),(7,8),(9,10)),\quad ((5,6),(7,9),(8,10)),\quad
((5,7),(6,10),(8,9)).$$
Here the quotient of the two nonconstant rows depends on a phase $s$. We keep $s$ free when computing extensions.

\par\addvspace{\medskipamount}\noindent
\begin{minipage}{\linewidth}
\begin{lemma}\label{ddd_pilots:normal}
Every presentation of class $3297$ is equivalent to one of the triples below, where $w=e^{2\pi i/5}$.
$$\begin{array}{c|c|c}
j=0,\ldots,4& j=5,6 & j=7,8,9,10\\\hline
1&(1,1)&(1,1,1,1)\\
w^j&s(1,1)&s(-1,-1,i,-i)\\
w^{aj}&p(1,-1)&q h
\end{array}$$
Here $a=2,3,4$, the vector $h$ is $(1,i,-1,-i)$ or $(1,-i,i,-1)$, and $s,p,q$ are independent unit phases.
\end{lemma}
\end{minipage}
\par\addvspace{\medskipamount}

\begin{proof}
After dephasing, write the last 6 entries of the rows as $(A,-A,B,-B,C,-C)$ and $(D,-D,E,F,-E,-F)$. The 3 pair conditions on their quotient give $B/E=C/F=-A/D$ and $BE=-CF$. Hence $F=\pm iE$. Take the quotient as second row by a row permutation and column factors. Put $s=A/D$, $p=1/D$, $q=1/E$. Order the columns of its pentagon by their labels. As in Lemma~\ref{ccc_common:normal}, the labels on the other pentagon are $aj$ with $a=2,3,4$. Ordering the last 4 columns to give the displayed second row yields the two choices of $h$. The Smith diagonal is $(1,1,2)$. The direct label check maps all 144 presentations to the 6 forms, 24 to each, and verifies that the 3 parameter coordinates form an invertible integer change of variables.
\end{proof}

These normal forms allow the generic extension calculation to be shared across the class. The result below also controls ten complete parameter slices.

\par\addvspace{\medskipamount}\noindent
\begin{minipage}{\linewidth}
\begin{theorem}\label{ddd_pilots:3297}
Each of the 6 forms in Lemma~\ref{ddd_pilots:normal} has a regular 5-row extension for every $(s,p,q)$. Its maximum is exactly 5 outside an explicit union of $1982$ monomial loci. For every $s$ satisfying $s^{10}=-1$, every regular extension has at most 9 rows for all $p,q$. In the entire class, the values of $s$ at which a 10-row extension exists form a finite set of roots of unity.
\end{theorem}
\end{minipage}
\par\addvspace{\medskipamount}

\begin{proof}
For the generic assertion retain $s,p,q$ as independent characters and test every presentation of a new row against $1,u,v$ as in Theorem~\ref{bbb_common:extension}. Each solver exhausts $1480490$ initial presentations. There are $552$ pair records from profile $\mathrm{C}$ and $19332$ from $\mathrm{D}$. Only $\mathrm{D}$ gives final records, 526 per form. The audit lists 1982 distinct nonconstant closing equations $s^i p^j q^k=e^{\pi i d/210}$ with $\max(|i|,|j|,|k|)\le5$. Off their union the retained records are exhaustive.

In each form there are 6 circles, and the 520 individual records are specializations of these circles. Their first 5 entries are the two unused Fourier rows, 3 circles for each. Two members of the same circle cannot be orthogonal because the blocks have sizes 5 and 6. Between distinct circles with the same Fourier row the scalar product on the last 6 columns vanishes identically because the fourth-root coefficients of each parameter monomial cancel in opposite pairs. The first block contributes 5. At most one row of each Fourier type can therefore be added. Conversely the verifier exhibits 2 circles of different Fourier types whose rows, with both circle phases set to one, have a mutual prime-cycle cover for all $s,p,q$. It checks every old and new quotient symbolically, giving the uniform lower bound.

If $s^{10}=-1$, the pair $(1,u)$ becomes group 21. Testing the rotations, reversals and column permutations for the 10 values of $s$ and the 6 forms gives 3 descriptions. For each, normalize the third-row coefficients separately on the pentagon, the 2-column block and the 4-column block. Apply the unrestricted two-parameter calculation of Theorem~\ref{ddd_pilots:3250}, with $p,q$ free. The two scans agree on $221196$ records. Each description has 84 curves and 1360 specified or intersection points, giving $4335$ graphs altogether.

The two graph tests bound every clique by 6 except on one curve per description. On each of these 3 curves the free records reduce to 18 circles. They fall into 3 groups of 6 by the support of the extra phase. Within a group compare the pair polynomials $A\bar A-B\bar B$ as in Theorem~\ref{ddd_one:3217}. In all $3\cdot3\binom63=180$ triples their common divisor over $\mathbb Q(e^{\pi i/10})$ is constant. Thus each group supplies at most 2 rows. This proves the uniform bound of 9 on all ten slices.

Finally fix the supports and root labels of every quotient of a proposed 10-row extension. Include $s,p,q$ among the phase variables. By Proposition~\ref{n11_phase:affine} the solutions form a finite union of tori translated by torsion points. Projection of each such torus to $s$ is either the full circle or a single torsion value. The full circle would meet the slices just excluded. There are only finitely many choices of decorated supports, proving the assertion. Determining the remaining possible values of $s$ and deciding square completion there remain open.
\end{proof}

\section{Mixed triples}\label{sec:mixed}

Normalize a 3-row matrix to rows $1,u,v$, and use the quotients $u,v,z=u/v$. The choices $\mathrm{CCD}$ and $\mathrm{CDF}$ mean that these 3 quotients admit the indicated decompositions. Different profiles can sometimes describe the same quotient, so we also check that the examples are genuinely mixed: at specified points the only profiles available are precisely the stated ones.

For the finite calculations we impose the following explicit restriction. In $\mathrm{CCD}$, all 3 chosen pentagons use the same 5 columns. In $\mathrm{CDF}$, the pentagons in $u$ and $v$ use the same 5 columns. We call this the \emph{common-pentagon case}. It leaves a substantial part of the unrestricted mixed-type problem open.

\medskip
Throughout this subsection put $w=e^{\pi i/15}$ and $\omega=w^{10}=e^{2\pi i/3}$. Write $r=(1,\omega,\omega^2)$; the pentagons use $w^6=e^{2\pi i/5}$. For $s,t,a\in\mathbb{T}$ consider
\[
\begin{array}{c|c|c}
\text{columns }j=0,\ldots,4&\text{columns }5,6,7&\text{columns }8,9,10\\ \hline
1&(1,1,1)&(1,1,1)\\
w^{6j}&s\,r&sa\,r\\
w^{12j}&t\,r&-ta\,r
\end{array}.
\]
Both $u$ and $v$ have a pentagon and 2 triangles. Their quotient has first block $w^{-6j}$ and last 6 entries
$$(s/t)(1,1,1,-1,-1,-1),$$
which split into 3 antipodal pairs. Thus the displayed $\mathrm{CCD}$ triple is regular of type $\mathrm{CCD}$ for every parameter choice. The dephased entries recover $s,t$, and a ratio in the second row recovers $a$, giving 3 independent parameters modulo the finite permutation ambiguities.

A point with no profile ambiguity is obtained by taking $s=t=1$ and $a=w^{6}$. Its 2 nontrivial rows are $w$ raised entrywise to
$$\begin{pmatrix}
0&6&12&18&24&0&10&20&6&16&26\\
0&12&24&6&18&0&10&20&21&1&11
\end{pmatrix}.$$
The exact cycle-cover check finds only $\mathrm{C}$ for $u$, only $\mathrm{C}$ for $v$, and only $\mathrm{D}$ for $u/v$. Each alternative cycle cover imposes monomial equations in the parameters. Since it fails at this point, its locus is a proper closed subset. There are finitely many covers, so the same uniqueness holds on an open dense parameter set.

\medskip
The common-pentagon $\mathrm{CCD}$ family also has larger regular extensions.

\par\addvspace{\medskipamount}\noindent
\begin{minipage}{\linewidth}
\begin{proposition}\label{n11_mixed:fiveconstruction}
The displayed family of $\mathrm{CCD}$ triples extends to a $6$-parameter family of regular $5\times11$ partial Hadamard matrices. Generically its $10$ triples consist of $6$ of type $\mathrm{CCD}$ and $4$ of type $\mathrm{CCC}$.
\end{proposition}
\end{minipage}
\par\addvspace{\medskipamount}

\begin{proof}
Write $r_k=(1,\omega^k,\omega^{2k})$ and set
$$B(a,b)=
\begin{pmatrix}
r_0&r_0\\
r_1&a r_1\\
r_1&-a r_1\\
r_2&b r_2\\
r_2&-b r_2
\end{pmatrix},
\qquad a,b\in\mathbb{T}.$$
This is a $5\times6$ partial Hadamard matrix. Between different Fourier modes each 3-column block sums to zero separately. Between the 2 rows of the same mode, the 2 block contributions cancel. Consequently
\[
H(a,b,s,t,p,q)= \left[\ F_5\ \middle|\ \operatorname{diag}(1,s,t,p,q)\,B(a,b)\ \right], \qquad a,b,s,t,p,q\in\mathbb{T},
\]
is a regular $5\times11$ matrix. Its first 3 rows are the displayed $\mathrm{CCD}$ triple.

Numbering rows from zero, the pairs $(1,2)$ and $(3,4)$ have type $\mathrm{D}$. All other pairs have type $\mathrm{C}$. These are the unique types for generic parameters. Hence its 10 3-row submatrices consist generically of
$$6\text{ of type }\mathrm{CCD},\qquad 4\text{ of type }\mathrm{CCC}.$$
The 4 phases $s,t,p,q$ are read from the first column of the 6-column block; ratios in its 2 3-column halves recover $a,b$. Thus the family $H(a,b,s,t,p,q)$ of Proposition~\ref{n11_mixed:fiveconstruction} has 6 genuine parameters. All 10 quotient decompositions and parameter coefficients were checked exactly.
\end{proof} 

The $\mathrm{CCC}$ triples inside this five-row construction connect its extension problem with the generic $\mathrm{CCC}$ theorem.

\par\addvspace{\medskipamount}\noindent
\begin{minipage}{\linewidth}
\begin{corollary}\label{ccc_common:mixedmax}
The $6$-parameter family $H(a,b,s,t,p,q)$ of Proposition~\ref{n11_mixed:fiveconstruction} is generically maximal at $5$ rows among regular partial Hadamard matrices.
\end{corollary}
\end{minipage}
\par\addvspace{\medskipamount}

\begin{proof}
Its rows $0,1,3$ form $S_3(s,sa,p,pb)$. The map $(a,b,s,t,p,q)\mapsto(s,sa,p,pb)$ is onto $\mathbb T^4$, with the 4 displayed coordinates independent. The inverse image of the exceptional set in Theorem~\ref{ccc_common:extension} is therefore a proper finite union of monomial loci. Outside it, a sixth row would contradict that theorem.
\end{proof}

\medskip
With the same $w=e^{\pi i/15}$, take first row all ones and let the other rows be
\[
\begin{aligned}
u(q)&=(w^{U_0},\ldots,w^{U_4}, qw^{U_5},\ldots,qw^{U_{10}}),\\
v(q)&=(w^{V_0},\ldots,w^{V_4}, qw^{V_5},\ldots,qw^{V_{10}}), \qquad q\in\mathbb{T},
\end{aligned}
\]
where
$$\begin{pmatrix}U\\ V\end{pmatrix}
=
\begin{pmatrix}
0&6&12&18&24&0&20&10&7&17&27\\
0&6&18&12&24&20&5&1&16&27&12
\end{pmatrix}.$$
The cycle supports are
$$\begin{array}{c|c|l}
\text{quotient}&\text{profile}&\text{cycles, with columns numbered from zero}\\ \hline
u&C&(0,1,2,3,4),\ (5,6,7),\ (8,9,10)\\
v&D&(0,1,2,3,4),\ (5,6),\ (7,8),\ (9,10)\\
u/v&F&(0,5,9),\ (1,6),\ (2,7),\ (3,8),\ (4,10).
\end{array}$$
Multiplication by $q$ rotates the last-block cycles of $u,v$, while the quotient $u/v$ is unchanged. These supports therefore certify regularity for all $q$. The cross ratio from the first 2 rows and columns $0,5$ recovers $q$, so this parameter is genuine.

At $q=1$, the only available profiles are exactly $\mathrm{C},\mathrm{D},\mathrm{F}$. This was verified by exhausting all exact covers by prime cycles; it can also be read directly from the exponent multisets. All entries are thirtieth roots, so only prime lengths $2,3,5$ can occur. The $u$ multiset has at most one disjoint antipodal pair and at most 2 disjoint triangles, excluding $\mathrm{D},\mathrm{E},\mathrm{F}$. The $v$ multiset has no triangle, excluding $\mathrm{C},\mathrm{E},\mathrm{F}$. The $u/v$ multiset has no pentagon and at most one disjoint triangle, excluding $\mathrm{C},\mathrm{D},\mathrm{E}$. Thus the example cannot be made homogeneous merely by choosing other cycle decompositions. Generic $q$ has the same unique profiles.

Unlike the displayed $\mathrm{CCD}$ triple, the orthogonality of $u$ and $v$ here is not separate on the 5-column and 6-column blocks. The triangle $(0,5,9)$ and the 4 pairs of $u/v$ cross their boundary. This is why the immediate concatenation argument used for the family $H(a,b,s,t,p,q)$ of Proposition~\ref{n11_mixed:fiveconstruction} does not apply. Extensions of the displayed $\mathrm{CDF}$ triple are not addressed by this calculation.

\medskip
We now give the exhaustive covers under the common-pentagon restrictions stated above.

\par\addvspace{\medskipamount}\noindent
\begin{minipage}{\linewidth}
\begin{proposition}\label{n11_mixed:CCDcover}
In the common-pentagon $\mathrm{CCD}$ case there is an exhaustive cover by $2052$ affine descriptions of dimension $2$ or $3$, all with seeds over $\mu_{30}$. These descriptions are taken in the support and label normalization below; they do not enumerate ordinary equivalence classes.
\end{proposition}
\end{minipage}
\par\addvspace{\medskipamount}

\begin{proof}
Put the common pentagon on columns $0,\ldots,4$ and use column permutations and row phases to set its entries in $u$ to $w^{6j}$, and in $v$ to $w^{6\pi(j)}$, where $\pi(0)=0$. For $u/v$ to have its pentagon on these same columns, both $\pi$ and $j-\pi(j)$ must be permutations modulo 5. Checking the $24$ permutations fixing zero gives exactly
$$\pi=(0,2,4,1,3),\quad(0,3,1,4,2),\quad(0,4,3,2,1).$$

Fix the 2 triangles of $u$ to be $(5,6,7),(8,9,10)$. There are 10 partitions of the remaining 6 columns into the 2 triangles of $v$, and $15$ pairings for $u/v$. For each of these $150$ support choices there are $16$ triangle orientation choices. 4 triangle phases remain after normalization.

Let $x\in(\mathbb R/\mathbb Z)^4$ be those phases. If column $j$ belongs to triangle $k$ of $u$ and triangle $\ell$ of $v$, write its quotient phase as
$$W_jx+c_j,\qquad W_jx=x_k-x_{2+\ell},\qquad c_j=\frac{\sigma_k h_j-\sigma_{2+\ell}h'_j}{3},$$
where $h_j,h'_j\in\{0,1,2\}$ are its positions in the triangles. For every prescribed pair $(a,b)$ of $u/v$ impose
$$(W_b-W_a)x=\frac12-c_b+c_a\pmod1.$$
These are 3 integer phase equations in 4 variables.

Unimodular changes of variables solve them exactly over the torus. At a step with coefficient gcd $g$, all branches $(\text{right-hand side}+k)/g$, $0\leq k<g$, are retained. Thus completeness does not depend on a finite-grid assumption.

The result is:
$$\begin{array}{l|r}
\text{support choices tested}&150\\
\text{feasible support choices}&96\\
\text{2-dimensional phase components}&648\\
\text{3-dimensional phase components}&36.
\end{array}$$
Each of the $684$ components permits all 3 pentagon-label assignments above, giving $2052$ affine descriptions. Every seed is in $\mu_{30}$. These are redundant descriptions, not matrix-equivalence classes.

An independent verifier tests all $150\cdot16\cdot6^4=3\,110\,400$ points of the sixth-root phase grids. It checks exact agreement with the reported components, their disjointness for fixed supports and signs, their dimensions, and primitivity of their integer parameter bases. Continuous completeness comes from the unimodular branch argument.
\end{proof}

\medskip
For $\mathrm{CDF}$, the common pentagon leaves a different condition on the final six entries. We again retain all admissible labels and their phase solutions.

\par\addvspace{\medskipamount}\noindent
\begin{minipage}{\linewidth}
\begin{proposition}\label{n11_mixed:CDFcover}
Every common-pentagon $\mathrm{CDF}$ triple is equivalent to a member of a $1$-parameter family with a seed over $\mu_{30}$. The normalization below gives $168$ distinct ordered seeds and an exhaustive, possibly redundant cover of these triples.
\end{proposition}
\end{minipage}
\par\addvspace{\medskipamount}

\begin{proof}
Normalize the common pentagon of $u,v$ in the same way, but now all $24$ permutations $\pi$ fixing zero must be allowed, because $u/v$ has type $\mathrm{F}$, not a prescribed pentagon. Its first 5 entries lie in $\mu_5$.

A pair or triangle meets $\mu_5$ in at most one entry: the ratio of 2 different entries cannot have both order dividing 5 and order dividing 2 or 3. Since $\mathrm{F}$ has exactly 5 blocks, each block must therefore contain exactly one of the first 5 columns. The remaining 6 columns supply 2 additional triangle entries and 4 pair partners. There are
$$5\binom62\,4!=1800$$
such supports.

Fix the 2 triangles of $u$ outside the pentagon. A matching into the 3 pairs of $v$ has either one or 3 pairs crossing between these triangles. Under the stabilizer of the triangles these are exactly 2 orbits, represented by
$$((5,6),(7,8),(9,10)),\qquad ((5,8),(6,9),(7,10)).$$
The 3-crossing case has a short obstruction. Let $x,y$ be the 2 triangle phases of $u$, and put $q=y/x$. Write $z=u/v$. All 6 outside entries of $z$ lie in $\mu_{30}$: 2 are the additional entries of its triangle and 4 are pair partners. Thus $z^{15}$ equals $+1$ at exactly 2 positions and $-1$ at 4. For a $v$ pair $(a,b)$ crossing from the first $u$ triangle to the second, antipodality gives
$$q=-\frac{z_b}{z_a}\,\omega^{\sigma h_a-\tau h_b}\in\mu_{30}.$$
Taking fifteenth powers shows that all 3 pairs must have the same sign ratio $z_b^{15}/z_a^{15}$. They cannot all have opposite signs, since the sign counts are 2 and 4. Hence the 2 positive entries are paired together and lie on opposite $u$ triangles. Multiplying the 3 displayed equations and raising to the tenth power now gives
$$1=q^{30}=\frac{\prod_b z_b^{10}}{\prod_a z_a^{10}}.$$
But the 4 pair-partner entries have tenth power $1$, whereas the 2 remaining entries have tenth powers $\omega,\omega^2$ and lie on opposite sides. The ratio is therefore $\omega$ or $\omega^2$, a contradiction.

Each of the 6 columns is an edge between its $u$ triangle phase and its $v$ pair phase. The resulting graph has 5 vertices, 6 edges, and is connected in both cases. Every column prescribes a difference of its 2 incident phases, with right-hand side in $\tfrac1{30}\mathbb Z/\mathbb Z$.

Integrating along a spanning tree leaves one common free phase. The 2 remaining edges impose exact compatibility conditions. Therefore every feasible presentation in this common-pentagon case is a one-parameter family with a seed over $\mu_{30}$.

The enumeration tests
$$2\cdot1800\cdot24\cdot8=691\,200$$
support, pentagon-label and triangle-orientation choices. There are $480$ feasible support presentations and $1920$ feasible label/orientation choices. All occur in the one-crossing-pair orbit. After setting the first $u$ triangle phase to zero and deduplicating identical ordered seed matrices, $168$ distinct seeds remain. 

An independent verifier uses direct pair identities rather than graph integration. Inside a $u$ triangle, the condition that the corresponding $v$ entries are opposite directly fixes the ratio of the 2 quotient entries. A pair crossing between the triangles fixes their relative phase; all such constraints must agree. This second method reproduces the same $1920$ feasible choices, $480$ supports and $168$ seeds.
\end{proof}

\section{Remaining problems and computations}\label{sec:N11classification}\label{sec:certificates}
The square classification in the $\mathrm{A}$ branch suggests the following conjecture, which remains open.

\par\addvspace{\medskipamount}\noindent
\begin{minipage}{\linewidth}
\begin{conjecture}
Every regular complex Hadamard matrix of order $11$ is equivalent to $F_{11}$.
\end{conjecture}
\end{minipage}
\par\addvspace{\medskipamount}

Of the $35$ unordered no-$\mathrm{A}$ profile triples, the $5$ homogeneous types have exhaustive affine descriptions. The $\mathrm{CCD}$ and $\mathrm{CDF}$ descriptions are complete only under the stated common-pentagon restrictions. The other $28$ mixed types, the unrestricted $\mathrm{CCD},\mathrm{CDF}$ cases, and the general square completions remain to be treated. No reduction of every square to a homogeneous triple has been proved.

For $\mathrm{DDD}$, all $54836$ triples admitting one of the zero-dimensional chosen presentations stop before $6$ rows. All labels are expanded for the $190,41,3$ positive-dimensional support classes. Theorems~\ref{ddd_one:bound} and~\ref{ddd_one:3217} exclude square completion throughout 2 of the 190 one-parameter classes. Class $3203$ has a uniform bound of 6 rows. Class $3217$ has exact maxima 3, 4 and 5 as specified in Theorem~\ref{ddd_one:3217}. Theorem~\ref{ddd_pilots:3250} excludes square completion throughout two-parameter class $3250$. For three-parameter class $3297$, Theorem~\ref{ddd_pilots:3297} gives generic maximum 5, a uniform bound of 9 on ten full slices, and a finite projection onto $s$ for possible 10-row extensions. The other classes and the remaining exceptional values still require analysis. The prescribed families $H$ and $G$ are uniformly maximal at $5$ rows. For the labelled triple $R$, the generic region and the $4$ surfaces $b/c\in\mu_4$ are settled. The remaining exceptional loci, including the unrestricted cube-root surfaces, require further analysis. The level-$120$ theorem covers only the specified entry group. Matrix-equivalence reduction of the continuous descriptions also remains open.

For the common-support $\mathrm{BBB}$ families $T_f$, the generic extension problem is settled at $3$ or $4$ rows. The $3038$ recorded exceptional equations and the other $\mathrm{BBB}$ support presentations remain to be analysed. The explicit $\mathrm{BDD}$ example proves occurrence only; it is not a classification of that mixed profile. For the common-block $\mathrm{CCC}$ families, which cover all $4$-dimensional chosen phase presentations, the generic maximum is $5$. Their $3608$ recorded exceptional equations and the other $\mathrm{CCC}$ support configurations remain open.

For $\mathrm{EEE}$, every $4$-dimensional chosen phase presentation is covered, up to equivalence and conjugation, by the family $H(s,t,a,b)$ of Subsection~\ref{sec:EEE}. Its generic maximum is $3$, and the explicit level-$204$ example has no regular fourth row. The $5108$ sufficient exceptional equations and the lower-dimensional $\mathrm{EEE}$ presentations remain to be analysed. For $\mathrm{FFF}$, both maximal support types are covered by the 6-parameter families $H_h(a,b,s,t,p,q)$, $h=0,1$, of Subsection~\ref{sec:FFF}, with generic maximum $3$ and explicit maximal examples of level $498$. Their $21403$ sufficient exceptional equations and the lower-dimensional $\mathrm{FFF}$ presentations remain open.

\medskip
The current companion is \texttt{code020.zip}, with root directory \texttt{code020}. Its \texttt{README.txt} maps the short folder names to the original runtime components, and \texttt{run.py} prepares each calculation in a separate working directory. The following mapping uses the runtime paths relative to \texttt{cases/}; names in the first column have the prefix \texttt{n11\_}, except \texttt{EEE}.

\begin{center}\small
\begin{tabular}{@{}>{\raggedright\arraybackslash}p{.35\textwidth}>{\raggedright\arraybackslash}p{.59\textwidth}@{}}
\toprule
Component & Calculation and verification\\\midrule
\texttt{A\_branch}, \texttt{A\_completion}&Support graphs, root labels and Fourier compatibility; \texttt{verify.py} checks the saved data, with a fresh $10!$ permutation enumeration for the square case.\\
\texttt{noA\_BBB}, \texttt{noA\_CCC}, \texttt{noA\_DDD}, \texttt{EEE}, \texttt{noA\_FFF}&Support partitions, integer phase systems, root-label tests and symbolic families. Each has a full \texttt{run.py} driver.\\
\texttt{BBB\_common}&Generic common-support $\mathrm{BBB}$ extensions and the $\mathrm{BDD}$ witness; \texttt{run.py --verify-only} checks all saved row certificates and the circle reduction.\\
\texttt{CCC\_common}&Generic common-block $\mathrm{CCC}$ extensions; \texttt{run.py --verify-only} checks all row certificates, the $6$ circles and exact compatibility.\\
\texttt{EEE\_generic}&Generic four-parameter $\mathrm{EEE}$ exclusion; \texttt{run.py --verify-only} checks scan agreement, pair covers, normal forms and the level-$204$ example.\\
\texttt{FFF\_generic}&Generic exclusion for both 6-parameter $\mathrm{FFF}$ types; \texttt{run.py --verify-only} checks scan agreement, pair covers, normal forms and the level-$498$ examples.\\
\texttt{mixed\_pilots}&Common-pentagon $\mathrm{CCD},\mathrm{CDF}$ and the 6-parameter 5-row family; \texttt{run.py}.\\
\texttt{DDD\_full\_family}&Complete $H(t)$ exclusion; \texttt{run.py} regenerates both arithmetic runs, \texttt{verify.py} checks delivered certificates.\\
\texttt{DDD\_fourth\_rows}&All quintic labels, finite triple equivalence, fourth rows and compatibility; \texttt{run.py}, \texttt{verify.py}.\\
\texttt{DDD\_one}&One-parameter $\mathrm{DDD}$ classes 3203 and 3217, all labels and phases, compatibility graphs, and exact maxima for 3217. \texttt{run.py} reproduces both calculations.\\
\texttt{DDD\_three\_parameter\_five}&The family $G$ and its uniform exclusion; \texttt{run.py} regenerates the calculation; \texttt{verify.py} checks the saved certificates.\\
\texttt{DDD\_pilots}&Two-parameter class 3250 and three-parameter class 3297. \texttt{run.py} reproduces the covers, phase domains and compatibility bounds.\\
\texttt{DDD\_generic\_triple}&The $8$ generic rows and $5001$ exceptional equations; \texttt{run.py --verify-only}.\\
\texttt{DDD\_exceptional\_surfaces}, \texttt{DDD\_imaginary\_surfaces}&Complete surfaces, special loci and compatibility. Both support \texttt{run.py --verify-only}; \texttt{verify\_triangles.py} checks the simultaneous-phase contradictions in the second.\\
\texttt{DDD\_bh120\_pilot}&All parameter points and actual row graphs in $\mu_{120}$; \texttt{run.py --verify-only}.\\
\bottomrule
\end{tabular}
\end{center}

The earlier \texttt{DDD\_sixth\_pilot}, \texttt{DDD\_exceptional\_circles} and \texttt{DDD\_two\_parameters} components retain the preliminary exclusions. The circle component also contains the stronger fourth-root-grid certificate and the orthogonal nonregular example. Their intermediate bounds are superseded by the full-family theorem.

At the package root, \texttt{python3 run.py check} checks file integrity; it is not a mathematical verifier. The commands \texttt{python3 run.py verify bbb\_common} and \texttt{python3 run.py reproduce bbb\_common} check the supplied certificates and regenerate the full common-support calculation respectively. Other component names are listed by \texttt{python3 run.py list}. Generated outputs stay in the prepared working directory. Python, a C++17 compiler and, for $\mathrm{EEE}$ and $\mathrm{FFF}$, Node.js are required. No network access is needed. The exact coverage and commands of each component are recorded in its notes.

\begingroup
\hypersetup{allcolors=black}
\begin{thebibliography}{99}
\addcontentsline{toc}{section}{References}
\setlength{\itemsep}{0pt}
\bibitem{inv} T. Banica, Invitation to Hadamard matrices, \href{https://arxiv.org/abs/1910.06911}{arXiv:1910.06911} (2024).

\bibitem{bbs} T. Banica, J. Bichon and J.-M. Schlenker, Representations of quantum permutation algebras, \emph{J. Funct. Anal.} \textbf{257} (2009), 2864--2910.

\bibitem{but} A. T. Butson, Generalized Hadamard matrices, \emph{Proc. Amer. Math. Soc.} \textbf{13} (1962), 894--898.

\bibitem{regular78} GPT and Teo, Regular complex Hadamard matrices of order 7,8, \href{https://hexagonmath.org/2611.00079v1}{hexagon:2611.00079} (2026).

\bibitem{ha1} U. Haagerup, Orthogonal maximal abelian $*$-subalgebras of the $n\times n$ matrices and cyclic $n$-roots, in \emph{Operator algebras and quantum field theory}, International Press (1997), 296--323.

\bibitem{hsc} G. Hiranandani and J.-M. Schlenker, Small circulant complex Hadamard matrices of Butson type, \emph{European J. Combin.} \textbf{51} (2016), 306--314.

\bibitem{hkm} M. Hirasaka, K.-T. Kim and Y. Mizoguchi, Uniqueness of Butson Hadamard matrices of small degrees, \emph{J. Discrete Algorithms} \textbf{34} (2015), 70--77.

\bibitem{KB} R. Kannan and A. Bachem, Polynomial algorithms for computing the Smith and Hermite normal forms of an integer matrix, \emph{SIAM J. Comput.} \textbf{8} (1979), 499--507.

\bibitem{lle} T. Y. Lam and K. H. Leung, On vanishing sums of roots of unity, \emph{J. Algebra} \textbf{224} (2000), 91--109.

\bibitem{pal} R. Paley, On orthogonal matrices, \emph{J. Math. Phys.} \textbf{12} (1933), 311--320.

\bibitem{tz1} W. Tadej and K. \.Zyczkowski, A concise guide to complex Hadamard matrices, \emph{Open Syst. Inf. Dyn.} \textbf{13} (2006), 133--177.

\bibitem{prime} {\raggedright Teo B and Chat GPT, On Fourier and Hadamard matrices of prime order, \href{https://hexagonmath.org/2611.00221v1}{hexagon:2611.00221} (2026).\par}

\end{thebibliography}
\endgroup
\end{document}
